% paralleles.tex — f164v, vv. 165–208 : loci paralleli. Généré par .travail/f164v_lp_redac/rediger.py
% à partir des collectes vérifiées (paralleles_165-186.json, paralleles_187-208.json) et des parallèles
% d'Austin–Olson 2004 retrouvés dans les sources ; règles : edition/CONVENTIONS.md § 4
\lp{165 αὐτός τε καλὸς ἦν]}{Plu.~\textit{Mor.}~645e (\textit{Quaest.\,conv.}~3.1.1) τὰς μὲν ἐν τοῖς μέλεσι παραχρώσεις βδελυττόμενον καὶ κατηγοροῦντα τοῦ καλοῦ Ἀγάθωνος, ὃν πρῶτον εἰς τραγῳδίαν φασὶν ἐμβαλεῖν καὶ ὑπομῖξαι τὸ χρωματικὸν, ὅτε τοὺς Μυσοὺς ἐδίδασκεν ; Ael.~\textit{VH}~13.4 Εὐριπίδης … Ἀγάθωνα τὸν τῆς τραγῳδίας ποιητὴν περιλαβὼν κατεφίλει, τετταράκοντα ἐτῶν που γεγονότα … οὐ γὰρ μόνον τὸ ἔαρ τῶν καλῶν κάλλιστον, ἀλλὰ καὶ τὸ μετόπωρον ; Suid.~α~124 Ἀγάθων· ὄνομα κύριον. τραγικὸς δὲ ἦν· διεβέβλητο δὲ ἐπὶ μαλακίᾳ … οὗτος ὁ Ἀγάθων ἀγαθὸς ἦν τὸν τρόπον, ποθεινὸς τοῖς φίλοις καὶ τὴν τράπεζαν λαμπρός … ἐκωμῳδεῖτο δὲ εἰς θηλύτητα ; Plu.~\textit{Alc.}~1.3 οὐ γάρ, ὡς Εὐριπίδης ἔλεγε, πάντων τῶν καλῶν καὶ τὸ μετόπωρον καλόν ἐστιν ; Plu.~\textit{Mor.}~177a–b τοῦ δὲ Εὐριπίδου τὸν καλὸν Ἀγάθωνα περιλαμβάνοντος ἐν τῷ συμποσίῳ καὶ καταφιλοῦντος ἤδη γενειῶντα … τῶν γὰρ καλῶν καὶ τὸ μετόπωρον καλόν ἐστιν ; Erasm.~\textit{Adag.}~1.3.72 Pulchrorum etiam autumnus pulcher est. Τῶν καλῶν καὶ ὀπώρη καλή … Metaphora proverbialis nata ex Archelai apophthegmate, quod ab eo dictum Plutarchus refert in Euripidem, qui jam pubescentem atque exoletum Agathonem in convivio suaviabatur ; Luc.~\textit{Rh.Pr.}~11 πάγκαλον ἄνδρα, διασεσαλευμένον τὸ βάδισμα, ἐπικεκλασμένον τὸν αὐχένα, γυναικεῖον τὸ βλέμμα, μελιχρὸν τὸ φώνημα … πάναβρόν τινα Σαρδανάπαλλον ἢ Κινύραν ἢ αὐτὸν Ἀγάθωνα, τὸν τῆς τραγῳδίας ἐπέραστον ἐκεῖνον ποιητήν ; Ath.~13.603e–604d οὐκ εὖ εἴρηκε Φρύνιχος πορφυρέας εἰπὼν τὰς γνάθους τοῦ καλοῦ ; Suid.~κ~251 Καλὸς κἀγαθός· λέγεται κατὰ συναλοιφήν, οὐχὶ καλὸς καὶ ἀγαθός. καὶ Καλός, ὁ ἐρώμενος … τὸ μὲν καλὸς ἐπὶ τῆς ἐν σώματι ὥρας, τὸ δὲ ἀγαθὸς ἐπὶ τῆς ἐν ψυχῇ ; Plu.~\textit{Mor.}~770c ἔφη γὰρ Ἀγάθωνα τὸν καλὸν ἤδη γενειῶντα περιβάλλων καὶ κατασπαζόμενος, ὅτι τῶν καλῶν καὶ τὸ μετόπωρον ; Ar.~\textit{Ach.}~143–144 ὑμῶν τ᾽ ἐραστὴς ἦν ἀληθὴς ὥστε καὶ / ἐν τοῖσι τοίχοις ἔγραφ᾽, Ἀθηναῖοι καλοί ; \textit{Il.}~2.673–674 Νιρεύς, ὃς κάλλιστος ἀνὴρ ὑπὸ Ἴλιον ἦλθε / τῶν ἄλλων Δαναῶν μετ᾽ ἀμύμονα Πηλεΐωνα ; Suid.~κ~265 Καλοί· Ἀριστοφάνης· ἐν τοῖσι τοίχοις ἔγραφον· Ἀθηναίοις καλός. ἴδιον ἐραστῶν ἦν τὸ τὰ τῶν ἐρωμένων ὀνόματα γράφειν ἐν τοῖς τοίχοις … οὕτως· ὁ δεῖνα καλός ; Luc.~\textit{DMort.}~9.4 ὑπεξυρημένος μὲν τὸ γένειον καὶ βαρβαρίζων, Κόδρου δὲ εὐγενέστερος καὶ Νιρέως καλλίων καὶ Ὀδυσσέως συνετώτερος λεγόμενος εἶναι}
\lp{165 αὐτός τε καλὸς ἦν καὶ καλῶς ἠμπέσχετο]}{Ar.~\textit{Thesm.}~148 ἐγὼ δὲ τὴν ἐσθῆθ᾽ ἅμα γνώμῃ φορῶ ; Ar.~\textit{Thesm.}~159–163 ἄλλως τ᾽ ἄμουσόν ἐστι ποιητὴν ἰδεῖν / ἀγρεῖον ὄντα καὶ δασύν· σκέψαι δ᾽ ὅτι / Ἴβυκος ἐκεῖνος κἀνακρέων ὁ Τήιος / κἀλκαῖος … ἐμιτροφόρουν τε καὶ διεκλῶντ᾽ Ἰωνικῶς ; Pi.~\textit{N.}~3.19 εἰ δ᾽ ἐὼν καλὸς ἔρδων τ᾽ ἐοικότα μορφᾷ ; Pl.~\textit{Smp.}~174a ταῦτα δὴ ἐκαλλωπισάμην, ἵνα καλὸς παρὰ καλὸν ἴω ; D.L.~2.28 καίτοι ἐνίοτε πρὸς τοὺς καιροὺς ἁρμοττόμενος καὶ λαμπρὰ ἠμπίσχετο· καθάπερ ἐν τῷ Πλάτωνος συμποσίῳ παρ᾽ Ἀγάθωνα βαδίζων ; Ar.~\textit{Thesm.}~191–192 σὺ δ᾽ εὐπρόσωπος λευκὸς ἐξυρημένος / γυναικόφωνος ἁπαλὸς εὐπρεπὴς ἰδεῖν ; Semon.~fr.~7.2–6~Bergk τὴν μὲν ἐξ ὑὸς τανύτριχος, / τῇ πάντ᾽ ἀν᾽ οἶκον βορβόρῳ πεφυρμένα / ἄκοσμα κεῖται … / αὐτὴ δ᾽ ἄλουτος ἀπλύτοις τ᾽ ἐν εἵμασιν / ἐν κοπρίῃσιν ἡμένη πιαίνεται ; Sapph. ap. Ath.~1.21a–d Σαπφὼ περὶ Ἀνδρομέδας σκώπτει· τίς δ᾽ ἀγροιῶτις θέλγει νόον οὐκ ἐπισταμένη τὰ βράκε᾽ ἕλκειν ἐπὶ τῶν σφυρῶν ; Hdt.~3.22.1 δολεροὺς μὲν τοὺς ἀνθρώπους ἔφη εἶναι, δολερὰ δὲ αὐτῶν τὰ εἵματα ; Pl.~\textit{Prt.}~315d–e νέον τι ἔτι μειράκιον, ὡς μὲν ἐγᾦμαι καλόν τε κἀγαθὸν τὴν φύσιν, τὴν δ᾽ οὖν ἰδέαν πάνυ καλός. ἔδοξα ἀκοῦσαι ὄνομα αὐτῷ εἶναι Ἀγάθωνα ; X.~\textit{Mem.}~2.1.22 τὴν δ᾽ ἑτέραν τεθραμμένην μὲν εἰς πολυσαρκίαν τε καὶ ἁπαλότητα, κεκαλλωπισμένην δὲ τὸ μὲν χρῶμα ὥστε λευκοτέραν τε καὶ ἐρυθροτέραν τοῦ ὄντος δοκεῖν φαίνεσθαι … ἐσθῆτα δὲ ἐξ ἧς ἂν μάλιστα ὥρα διαλάμποι ; Ath.~12.543c–544a οὗτος γὰρ παρὰ μέλος ὑπὲρ τὴν γραφικὴν τρυφήσας λόγῳ τῆς ἀρετῆς ἀντελαμβάνετο καὶ ἐπέγραφεν … ἁβροδίαιτος ἀνὴρ ἀρετήν τε σέβων τάδ᾽ ἔγραψεν ; Ael.~\textit{VH}~9.11 Παρράσιος ὁ ζωγράφος ὅτι μὲν πορφυρίδα ἐφόρει καὶ χρυσοῦν στέφανον περιέκειτο μαρτυροῦσι καὶ ἄλλοι καὶ τὰ ἐπιγράμματα}
\lp{165 ἠμπέσχετο]}{Ar.~\textit{Eq.}~893 καὶ τοῦτό ⟨γ᾽⟩ ἐπίτηδές σε περιήμπεσχ᾽, ἵνα σ᾽ ἀποπνίξῃ ; Ar.~\textit{Ran.}~1063–1064 πρῶτον μὲν τοὺς βασιλεύοντας ῥάκι᾽ ἀμπισχών, ἵν᾽ ἐλεινοὶ / τοῖς ἀνθρώποις φαίνοιντ᾽ εἶναι ; Ar.~\textit{Eccl.}~540 ἔπειθ᾽ ἵν᾽ ἀλεαίνοιμι, τοῦτ᾽ ἠμπεσχόμην ; \textit{idem} Synag. Σᵇ~α~951 (s.v. ἀλεαίνειν ; οὐκοῦν ἵν᾽ ἀλεαινέμην) ; Eur.~\textit{Med.}~1159 λαβοῦσα πέπλους ποικίλους ἠμπέσχετο ; Hsch.~η~552 ἠμπέσχετο· ἐνεδύσατο. ἐφόρησεν. περιεβάλλετο ; \textit{idem} Synag.~η~88 (om. ἐνεδύσατο. ἐφόρησεν.), Phot.~η~177 (om. ἐνεδύσατο. ἐφόρησεν.), Suid.~η~362 (om. ἐνεδύσατο. ἐφόρησεν. ; add. καὶ ἠμπίσχετο· δοτικῇ. ἠμπείχετο δέ) ; Ar.~\textit{Lys.}~1156 τὸν δῆμον ὑμῶν χλαῖναν ἠμπέσχον πάλιν ; Eur.~\textit{Ion}~1159 ἤμπισχεν ἄλλα βαρβάρων ὑφάσματα ; Ath.~12.512a–d ἁλουργῆ μὲν γὰρ ἠμπίσχοντο ἱμάτια, ποικίλους δ᾽ ὑπέδυνον χιτῶνας ; Hsch.~η~551 ἤμπισχε· περιέσχε. περιέβαλεν ; Zonar.~994 Ἠμπίσχετο· ἐνεδύετο. καὶ ἠμπέσχετο. ἠμπείχετο δὲ δίφθογγον ; Ar.~\textit{Ach.}~1024 ὦ τρισκακόδαιμον εἶτα λευκὸν ἀμπέχει ; Ar.~\textit{Av.}~1089–1090 οἳ χειμῶνος μὲν / χλαίνας οὐκ ἀμπισχνοῦνται ; Eur.~\textit{Supp.}~165 πίτνων πρὸς οὖδας γόνυ σὸν ἀμπίσχειν χερί ; Pl.~\textit{Phd.}~87b τεκμήριον δὲ παρέχοιτο θοἰμάτιον ὃ ἠμπείχετο αὐτὸς ὑφηνάμενος ; Ael.~\textit{VH}~7.9 ἠμπείχετο δὲ πρώτῃ μὲν τῇ σωφροσύνῃ, δευτέροις γε μὲν τοῖς παροῦσιν ; Hsch.~α~3788 ἀμπέχειν· περιβάλλεσθαι ; Hsch.~α~3791 ἀμπεχόμενος· περιβαλλόμενος ; Hsch.~α~3792 ἀμπέχονον· σύμμετρον περίβλημα (Ar. fr. 332, 7 K.–A.) ; Hsch.~α~3793 ἀμπεχόνην· λεπτὸν ἱμάτιον (Plat. Charm. 173b) ; Hsch.~α~3801–3802 ἀμπίσχειν· περιβαλεῖν. ἀμπισχόμενον· περιβαλλόμενον ; Hsch.~η~553 ἠμπισχόμενον· περιβαλλόμενον ; Phot.~η~182 ἠμπείχετο (Pl. Phaed. 87b)· περιεβέβλητο ; Zonar.~159 Ἀμπίσχω. τὸ ἐνδύομαι. [καὶ ἠμπισχημένος, καὶ ἀμπέχω, καὶ ἀμπέχομαι, ἠμπίσχετο. ἠμπειχόμην δὲ καὶ ἠμπείχετο, δίφθογγα.]}
\lp{165–166 καλὸς … κάλ᾽ (ὅ τι καλὸν φίλον)]}{Thgn.~17 ὅττι καλόν, φίλον ἐστί· τὸ δ᾽ οὐ καλὸν οὐ φίλον ἐστίν ; \textit{idem} Eur.~\textit{Ba.}~881 (ὅ τι καλὸν φίλον ἀεί), Erasm.~\textit{Adag.} (LB II~1117 F) (Quod pulcrum, idem amicum. Τὸ καλὸν φίλον), Pl.~\textit{Ly.}~216c (κατὰ τὴν ἀρχαίαν παροιμίαν τὸ καλὸν φίλον εἶναι)}
\lp{166 διὰ τοῦτ᾽ ἄρ᾽ … καί]}{Ar.~\textit{Av.}~486 διὰ ταῦτ᾽ ἄρ᾽ ἔχων καὶ νῦν ὥσπερ βασιλεὺς ὁ μέγας διαβάσκει ; Σ Ar.~\textit{Nub.}~319 ταῦτ᾽· διὰ ταῦτα. λείπει ἡ διά. ἀντὶ τοῦ διὰ ταῦτα, Ἀττικῶς ; Ar.~\textit{Pax}~891–892 τουτὶ δ᾽ ὁρᾶτ᾽ ὀπτάνιον ἡμῖν ὡς καλόν. / διὰ ταῦτα καὶ κεκάπνικ᾽ ἄρ᾽ ; Pl.~\textit{Prt.}~341b τί ἔλεγεν, ὦ Πρόδικε, τὸ χαλεπὸν Σιμωνίδης; κακόν, ἔφη. διὰ ταῦτ᾽ ἄρα καὶ μέμφεται, ἦν δ᾽ ἐγώ, ὦ Πρόδικε, τὸν Πιττακόν ; Σ Ar.~\textit{Nub.}~335 ταῦτ᾽ ἄρ᾽ ἐποίουν· λείπει ἡ διά ; Pl.~\textit{Grg.}~520d διὰ ταῦτ᾽ ἄρα, ὡς ἔοικε, τὰς μὲν ἄλλας συμβουλὰς συμβουλεύειν λαμβάνοντα ἀργύριον ; D.~21.187 διὰ τοῦτ᾽ ἄρα ; Hsch.~α~6922b ἄρα· λοιπόν. τάχα. δῆτα. δή}
\lp{166 κάλ᾽ ἦν τὰ δράματα]}{Ar.~\textit{Thesm.}~49–57 μέλλει γὰρ ὁ καλλιεπὴς Ἀγάθων / πρόμος ἡμέτερος … δρυόχους τιθέναι δράματος ἀρχάς. / κάμπτει δὲ νέας ἁψῖδας ἐπῶν, / τὰ δὲ τορνεύει, τὰ δὲ κολλομελεῖ ; Ar.~\textit{Thesm.}~149–150 χρὴ γὰρ ποιητὴν ἄνδρα πρὸς τὰ δράματα / ἃ δεῖ ποιεῖν πρὸς ταῦτα τοὺς τρόπους ἔχειν ; Ar.~\textit{Thesm.}~155–156 ἃ δ᾽ οὐ κεκτήμεθα, / μίμησις ἤδη ταῦτα συνθηρεύεται ; Ar.~\textit{Ran.}~83–84 Ἀγάθων δὲ ποῦ ᾽στιν; — ἀπολιπών μ᾽ ἀποίχεται, / ἀγαθὸς ποιητὴς καὶ ποθεινὸς τοῖς φίλοις ; Pl.~\textit{Smp.}~198a–b ὡς πρεπόντως τοῦ νεανίσκου εἰρηκότος καὶ αὑτῷ καὶ τῷ θεῷ … μετὰ καλὸν οὕτω καὶ παντοδαπὸν λόγον ῥηθέντα … τὸ δὲ ἐπὶ τελευτῆς τοῦ κάλλους τῶν ὀνομάτων καὶ ῥημάτων τίς οὐκ ἂν ἐξεπλάγη ἀκούων ; Ar.~\textit{Ach.}~470 ἀπολεῖς μ᾽. ἰδού σοι. φροῦδά μοι τὰ δράματα ; Arist.~\textit{Po.}~9.1451b19–23 ἐν ἐνίαις δὲ οὐθέν, οἷον ἐν τῷ Ἀγάθωνος Ἀνθεῖ· ὁμοίως γὰρ ἐν τούτῳ τά τε πράγματα καὶ τὰ ὀνόματα πεποίηται, καὶ οὐδὲν ἧττον εὐφραίνει ; Ael.~\textit{VH}~14.13 πολλοῖς καὶ πολλάκις χρῆται τοῖς ἀντιθέτοις ὁ Ἀγάθων … ἀλλὰ σύ γε, γενναῖε, λέληθας σεαυτὸν τὸν Ἀγάθωνα ἐκ τοῦ Ἀγάθωνος ἀφανίζων ; Synag.~δ~360 δρᾶμα· ποίημα. ὡς καὶ δρᾶσαι πρᾶξαι. λέγεται δὲ δρᾶμα καὶ τὰ ὑπὸ τῶν θεατρικῶν μιμηλῶς γινόμενα, ὡς ἐν ὑποκρίσει ; \textit{idem} Suid.~δ~1498 (add. πρᾶγμα) ; Hsch.~δ~2310 δράματα· ποιήματα. … πρᾶξις ἢ λόγος περιέχων ὑπόθεσίν τινα ; Hsch.~δ~2311 δραματουργεῖν· πανουργεῖν. συντιθέναι. πράττειν ; Hsch.~κ~377 καλα· ὀξυτόνως τὸ σύνηθες, καὶ ἀντὶ ἐπιρρήματος τοῦ καλῶς· παροξυτόνως δὲ τὰ ξύλα}
\lp{166 κάλ᾽ ἦν τὰ δράματα (Ἀγαθώνειος αὔλησις)]}{Hsch.~α~281 Ἀγαθώνειον αὔλησιν τὴν μαλακὴν (Ar. fr. 178 K.–A.). Ἀγαθὼν γὰρ ὁ τραγικὸς ἐπὶ μαλακίᾳ διεβάλλετο ; \textit{idem} Phot.~α~83 (Ἀγαθώνειος αὔλησις· ἡ μαλακή) ; Zen.~1.2 Ἀγαθώνιος αὔλησις· ἡ μαλακή, καὶ μήτε πικρὰ μήτε χαλαρά, ἀλλ᾽ εὔκρατος καὶ ἡδίστη. Εἴρηται δὲ ἀπὸ Ἀγάθωνος αὐλητοῦ, ὃς ἐπὶ μαλακίᾳ ἐκωμῳδεῖτο ; \textit{idem} Suid.~α~125 (ἡ μαλακὴ καὶ ἐκλελυμένη· ἢ ἡ μήτε χαλαρά, μήτε πικρά ; om. εἴρηται … ἐκωμῳδεῖτο), Diogenian.~1.7 (Ἀγαθώνειος ; μαλθακὴ καὶ μήτε χλιαρὰ μήτε πικρά ; om. εὔκρατος καὶ et εἴρηται … ἐκωμῳδεῖτο), Diogenian. Vind.~1.6 (Ἀγαθώνειος ; ἀγαθὴ καὶ μέση καὶ καλή· Ἀγάθων γάρ τις δόκιμος εἰς αὐλητικήν pro ἡ μαλακή … ἐκωμῳδεῖτο), Greg. Cypr.~1.33 (Ἀγαθώνειος ; ἡ ἡδίστη καὶ εὐφραντή pro ἡ μαλακή … ἐκωμῳδεῖτο), Apostol.~1.7 (Ἀγαθώνειος ; ἡ μήτε χαλαρὰ μήτε πικρά, ἀλλ᾽ εὔκρατος καὶ ἡδίστη, μαλακὴ καὶ ἐκλελυμένη ; οὗτος ἀγαθὸς ἦν τὸν τρόπον καὶ ποθεινὸς τοῖς φίλοις. ἐκωμῳδεῖτο δὲ εἰς θηλύτητα pro ὃς ἐπὶ μαλακίᾳ ἐκωμῳδεῖτο) ; Erasm.~\textit{Adag.}~2.4.63 Agathonia cantio. Ἀγαθώνιος αὔλησις … Recte dicetur de oratione blanda magis, quam frugifera. Agathon tibicen quispiam erat, cantilenarum suavitate majorem in modum aures deliniens. Nec hujus mores dissimiles musicae fuisse perhibentur. Laboravit enim infamia mollitiei … Hujus multis modis Aristophanes in Thesmophoriis carminum mollitiem deridet}
\lp{167 ὅμοια γὰρ ποιεῖν ἀνάγκη τῇ φύσει]}{Ar.~\textit{Ach.}~410–413 ἀναβάδην ποιεῖς, / ἐξὸν καταβάδην; οὐκ ἐτὸς χωλοὺς ποιεῖς. / ἀτὰρ τί τὰ ῥάκι᾽ ἐκ τραγῳδίας ἔχεις, / ἐσθῆτ᾽ ἐλεινήν; οὐκ ἐτὸς πτωχοὺς ποιεῖς ; Ar.~\textit{Ran.}~1058–1059 ἀλλ᾽ ὦ κακόδαιμον ἀνάγκη / μεγάλων γνωμῶν καὶ διανοιῶν ἴσα καὶ τὰ ῥήματα τίκτειν ; Eur.~\textit{Ba.}~369 τοῖς πράγμασιν δέ· μῶρα γὰρ μῶρος λέγει ; Pi.~\textit{N.}~1.25 τέχναι δ᾽ ἑτέρων ἕτεραι· χρὴ δ᾽ ἐν εὐθείαις ὁδοῖς στείχοντα μάρνασθαι φυᾷ ; Pi.~\textit{O.}~2.86 σοφὸς ὁ πολλὰ εἰδὼς φυᾷ· μαθόντες δὲ λάβροι ; Pl.~\textit{Ap.}~22b–c ὅτι οὐ σοφίᾳ ποιοῖεν ἃ ποιοῖεν, ἀλλὰ φύσει τινὶ καὶ ἐνθουσιάζοντες ; Pl.~\textit{R.}~3.395d ὅτι αἱ μιμήσεις, ἐὰν ἐκ νέων πόρρω διατελέσωσιν, εἰς ἔθη τε καὶ φύσιν καθίστανται καὶ κατὰ σῶμα καὶ φωνὰς καὶ κατὰ τὴν διάνοιαν ; Pl.~\textit{R.}~3.400d τί δ᾽ ὁ τρόπος τῆς λέξεως, ἦν δ᾽ ἐγώ, καὶ ὁ λόγος; οὐ τῷ τῆς ψυχῆς ἤθει ἕπεται ; Pl.~\textit{R.}~8.549e καὶ ἄλλα δὴ ὅσα καὶ οἷα φιλοῦσιν αἱ γυναῖκες περὶ τῶν τοιούτων ὑμνεῖν. καὶ μάλ᾽, ἔφη ὁ Ἀδείμαντος, πολλά τε καὶ ὅμοια ἑαυταῖς ; Arist.~\textit{Po.}~4.1448b24 διεσπάσθη δὲ κατὰ τὰ οἰκεῖα ἤθη ἡ ποίησις· οἱ μὲν γὰρ σεμνότεροι τὰς καλὰς ἐμιμοῦντο πράξεις … οἱ δὲ εὐτελέστεροι τὰς τῶν φαύλων ; Arist.~\textit{Po.}~17.1455a30–33 πιθανώτατοι γὰρ ἀπὸ τῆς αὐτῆς φύσεως οἱ ἐν τοῖς πάθεσίν εἰσιν … διὸ εὐφυοῦς ἡ ποιητική ἐστιν ἢ μανικοῦ ; Arist.~\textit{Po.}~25.1460b33 οἷον καὶ Σοφοκλῆς ἔφη αὐτὸς μὲν οἵους δεῖ ποιεῖν, Εὐριπίδην δὲ οἷοι εἰσίν ; Plb.~12.24.1 φησὶ γὰρ τοὺς ποιητὰς καὶ συγγραφέας διὰ τῶν ὑπεράνω πλεονασμῶν ἐν τοῖς ὑπομνήμασι διαφαίνειν τὰς ἑαυτῶν φύσεις ; Σ Pi.~\textit{N.}~1.37 χρὴ καθ᾽ ἣν γεγέννηταί τις φύσιν, ταύτῃ ἀκολουθεῖν καὶ μὴ βιάζεσθαι αὐτὴν εἰς ἄλλα τρέποντα ; Eur.~\textit{Hipp.}~701 πρὸς τὰς τύχας γὰρ τὰς φρένας κεκτήμεθα ; Eur.~\textit{Supp.}~180–183 τόν θ᾽ ὑμνοποιὸν αὐτὸς ἃν τίκτῃ μέλη / χαίροντα τίκτειν· ἢν δὲ μὴ πάσχῃ τόδε, / οὔτοι δύναιτ᾽ ἂν οἴκοθέν γ᾽ ἀτώμενος / τέρπειν ἂν ἄλλους ; \textit{Od.}~8.169–177 ἄλλος μὲν γάρ τ᾽ εἶδος ἀκιδνότερος πέλει ἀνήρ, / ἀλλὰ θεὸς μορφὴν ἔπεσι στέφει … / ἄλλος δ᾽ αὖ εἶδος μὲν ἀλίγκιος ἀθανάτοισιν, / ἀλλ᾽ οὔ οἱ χάρις ἀμφιπεριστέφεται ἐπέεσσιν ; Pi.~\textit{O.}~9.100 τὸ δὲ φυᾷ κράτιστον ἅπαν ; Catull.~16.5–6 nam castum esse decet pium poetam / ipsum, versiculos nihil necesse est ; Pl.~\textit{R.}~3.401a καὶ ἡ μὲν ἀσχημοσύνη καὶ ἀρρυθμία καὶ ἀναρμοστία κακολογίας καὶ κακοηθείας ἀδελφά ; [Arist.]~\textit{Phgn.}~805a1 ὅτι διάνοιαι ἕπονται τοῖς σώμασι, καὶ οὐκ εἰσὶν αὐταὶ καθ᾽ ἑαυτὰς ἀπαθεῖς οὖσαι τῶν τοῦ σώματος κινήσεων ; Dialex.~3.10 (Dissoi Logoi,~DK~90) ἐν γὰρ τραγωιδοποιίαι καὶ ζωγραφίαι ὅστις ⟨κα⟩ πλεῖστα ἐξαπατῆι ὅμοια τοῖς ἀληθινοῖς ποιέων, οὗτος ἄριστος ; Longin.~9.3 οὐδὲ γὰρ οἷόν τε μικρὰ καὶ δουλοπρεπῆ φρονοῦντας καὶ ἐπιτηδεύοντας παρ᾽ ὅλον τὸν βίον θαυμαστόν τι καὶ τοῦ παντὸς αἰῶνος ἐξενεγκεῖν ἄξιον ; Tosi n°~2186 Vestis virum reddit ; Ar.~\textit{Ran.}~939–944 ἀλλ᾽ ὡς παρέλαβον τὴν τέχνην παρὰ σοῦ τὸ πρῶτον εὐθὺς / οἰδοῦσαν ὑπὸ κομπασμάτων καὶ ῥημάτων ἐπαχθῶν, / ἴσχνανα μὲν πρώτιστον αὐτὴν καὶ τὸ βάρος ἀφεῖλον ; Hdt.~3.81.3 ἀρίστων δὲ ἀνδρῶν οἰκὸς ἄριστα βουλεύματα γίνεσθαι ; Hsch.~α~4231 ἀναγκαίῃ· ἀνάγκῃ ; Hsch.~ο~755 ὅμοια· ἴσα ; Hsch.~φ~1057 φύσις· γένος. οὐσία. προαίρεσις}
\lp{167 ὅμοια … τῇ φύσει (οἷος ὁ τρόπος, τοιοῦτος ὁ λόγος)]}{Isoc.~3.7 λόγος ἀληθὴς καὶ νόμιμος καὶ δίκαιος ψυχῆς ἀγαθῆς καὶ πιστῆς εἴδωλόν ἐστιν ; Cic.~\textit{Tusc.}~5.47 qualis cuiusque animi adfectus esset, talem esse hominem; qualis autem homo ipse esset, talem eius esse orationem; orationi autem facta similia, factis vitam ; Sen.~\textit{Ep.}~114.1 Hoc quod audire vulgo soles, quod apud Graecos in proverbium cessit: talis hominibus fuit oratio qualis vita ; \textit{idem} Erasm.~\textit{Adag.}~1.6.50 (Qualis vir, talis oratio (lemme) ; om. hoc quod audire vulgo soles, quod ; add. inquit) ; Sen.~\textit{Ep.}~114.4 Quomodo Maecenas vixerit notius est, quam ut narrari nunc debeat … Non oratio eius aeque soluta est quam ipse discinctus? ; Quint.~11.1.30 profert enim mores plerumque oratio et animi secreta detegit. nec sine causa Graeci prodiderunt, ut vivat, quemque etiam dicere ; Aristid.~\textit{Or.}~45 (p.~133 Dind. ;~2.392 Behr) οὐκ ἀποστατεῖ δὲ οὐδὲ ἡ παροιμία τούτων ἡ λέγουσα οἷος ὁ τρόπος, τοιοῦτον εἶναι καὶ τὸν λόγον· καὶ πάλιν τὸ ἕτερον ὡσαύτως ; \textit{idem} Apostol.~12.42c (οἷος ὁ τρόπος, τοιοῦτος ὁ λόγος), Tosi n°~1711 (οἷον ὁ τρόπος, τοιοῦτος ὁ λόγος) ; D.L.~1.58 Ἔλεγε δὲ τὸν μὲν λόγον εἴδωλον εἶναι τῶν ἔργων ; Ph.~\textit{Mos.}~1.29 ἵν᾽ οἷος ὁ λόγος τοιοῦτος ὁ βίος καὶ οἷος ὁ βίος τοιοῦτος ὁ λόγος ; Men.~\textit{Mon.}~26~Meineke Ἀνδρὸς χαρακτὴρ ἐκ λόγου γνωρίζεται}
\lp{167 τῇ φύσει (τὸ ἐμφυὲς ἦθος)]}{Ar.~\textit{Pax}~1083 οὔποτε ποιήσεις τὸν καρκίνον ὀρθὰ βαδίζειν ; \textit{idem} Macar.~6.79 (οὔτε … βαδίζειν διδάξεις), Apostol.~13.46a (add. οὐδ᾽ ἐπὶ τῷ πραχθέντι ποιήσεις ὕστερον οὐδέν· Ἀριστοφάνους ἐν Εἰρήνῃ), Tosi n°~2026 ; Eur.~fr.~75~N.² ὦ παῖ Κρέοντος, ὡς ἀληθὲς ἦν ἄρα, / ἐσθλῶν ἀπ᾽ ἀνδρῶν ἐσθλὰ γίγνεσθαι τέκνα, / κακῶν δ᾽ ὅμοια τῇ φύσει τῇ τοῦ πατρός ; \textit{idem} Apostol.~18.62f ; Pi.~\textit{O.}~11.19–20 τὸ γὰρ ἐμφυὲς οὔτ᾽ αἴθων ἀλώπηξ οὔτ᾽ ἐρίβρομοι λέοντες διαλλάξαντο ἦθος ; Pi.~\textit{O.}~13.13 ἄμαχον δὲ κρύψαι τὸ συγγενὲς ἦθος ; Men.~\textit{Mon.}~531~Meineke Φύσιν πονηρὰν μεταβαλεῖν οὐ ῥᾴδιον ; Ar.~\textit{Av.}~766–767 εἰ δ᾽ ὁ Πεισίου προδοῦναι τοῖς ἀτίμοις τὰς πύλας / βούλεται, πέρδιξ γενέσθω, τοῦ πατρὸς νεοττίον ; Isoc.~15.138 ὀρθῶς μὲν ἔφασκέ με λέγειν, οὐ μὴν οἷός τ᾽ ἦν τὴν φύσιν μεταβαλεῖν ; Σ Pi.~\textit{O.}~11.19c τὸ γὰρ ἦθος τὸ ἐκ τῆς φύσεως οὐκ ἂν μεταβάλοι οὔτε ἡ δολία ἀλώπηξ οὔτε οἱ φοβεροὶ λέοντες ; Macar.~8.43 Τοῦ πατρὸς τὸ παιδίον: ἐπὶ τοῦ ὁμοίου ; Apostol.~12.66 Ὁ λύκος τὴν τρίχα οὐ τὴν γνώμην ἀλλάττει ; Tosi n°~375 Talis pater talis filius ; Tosi n°~1659 Naturam expellas furca, tamen usque recurret}
\lp{168–170 αἰσχρὸς … κακὸς … ψυχρός]}{Arist.~\textit{Rh.}~3.3.4 (1406b5–10) καὶ ἔτι τέταρτον τὸ ψυχρὸν ἐν ταῖς μεταφοραῖς γίνεται … οἷον Γοργίας χλωρὰ καὶ ἄναιμα τὰ πράγματα, σὺ δὲ ταῦτα αἰσχρῶς μὲν ἔσπειρας κακῶς δὲ ἐθέρισας· ποιητικῶς γὰρ ἄγαν}
\lp{168 ταῦτ᾽ ἄρ᾽]}{Ar.~\textit{Nub.}~335 ταῦτ᾽ ἄρ᾽ ἐποίουν ὑγρᾶν Νεφελᾶν στρεπταιγλᾶν δάιον ὁρμάν ; Ar.~\textit{Nub.}~353–354 ταῦτ᾽ ἄρα ταῦτα Κλεώνυμον αὗται τὸν ῥίψασπιν χθὲς ἰδοῦσαι, / ὅτι δειλότατον τοῦτον ἑώρων, ἔλαφοι διὰ τοῦτ᾽ ἐγένοντο ; Ar.~\textit{Nub.}~394 ταῦτ᾽ ἄρα καὶ τὠνόματ᾽ ἀλλήλοιν βροντὴ καὶ πορδὴ ὁμοίω ; Ar.~\textit{Ach.}~90 ταῦτ᾽ ἄρ ἐφενάκιζες σὺ δύο δραχμὰς φέρων ; Ar.~\textit{Nub.}~319 ταῦτ᾽ ἄρ᾽ ἀκούσασ᾽ αὐτῶν τὸ φθέγμ᾽ ἡ ψυχή μου πεπότηται ; Ar.~\textit{Thesm.}~649 ὦ μιαρὸς οὗτος· ταῦτ᾽ ἄρ᾽ ὑπὲρ Εὐριπίδου ; Ar.~\textit{Eq.}~125 ὦ μιαρὲ Παφλαγὼν ταῦτ᾽ ἄρ᾽ ἐφυλάττου πάλαι ; Ar.~\textit{Pax}~414 ταῦτ᾽ ἄρα πάλαι τῶν ἡμερῶν παρεκλεπτέτην ; Eur.~\textit{IT}~932 ταῦτ᾽ ἆρ᾽ ἐπ᾽ ἀκταῖς κἀνθάδ᾽ ἠγγέλης μανείς}
\lp{168 Φιλοκλέης]}{\textit{Test.} Suid.~φ~379 Φιλοκλῆς, κωμῳδίας ποιητής, αἰσχροπρόσωπος. ἦν δὲ προκέφαλος, ὡς ἔποψ, ἤγουν ὀξυκέφαλος· Αἰσχύλου ἀδελφῆς υἱός. οἱ δὲ Ἁλμίωνα αὐτόν φασιν ἐπιθετικῶς, διὰ τὸ πικρὸν εἶναι· ἅλμη γὰρ ἡ πικρία ; Ar.~\textit{Vesp.}~461–462 εἴπερ ἔτυχον τῶν μελῶν τῶν Φιλοκλέους βεβρωκότες ; Ar.~\textit{Av.}~281–282 οὑτοσὶ μέν ἐστι Φιλοκλέους / ἐξ ἔποπος ; Ar.~\textit{Av.}~1295 κορυδὸς Φιλοκλέει, χηναλώπηξ Θεογένει ; Ar.~\textit{Ran.}~151 ἢ Μορσίμου τις ῥῆσιν ἐξεγράψατο ; Cratin.~fr.~292~Kock ὅνπερ Φιλοκλέης τὸν λόγον διέφθορεν ; \textit{idem} Σ Soph.~\textit{Ant.}~404 (praem. Κρατῖνος) ; Telecl.~fr.~14~Kock ἀλλ᾽ ἡ τάλαινα Φιλοκλέα * * * * οὖν, εἰ δ᾽ ἐστὶν Αἰσχύλου φρόνημ᾽ ἔχων ; Σ Ar.~\textit{Av.}~281 ἔστι δὲ ὁ Φιλοκλῆς τραγῳδίας ποιητής, καὶ Φιλοπείθους υἱὸς ἐξ Αἰσχύλου ἀδελφῆς. ὅσοι δὲ Ἁλμίωνος αὐτόν φασιν, ἐπιθετικῶς λέγουσι διὰ τὸ πικρὸν εἶναι. ἅλμη γὰρ ἡ πικρία ; Σ Ar.~\textit{Av.}~1295 κορυδός· μήποτε ὀξυκέφαλος ἦν εἰς τὸ ἄνω καὶ ὀρνιθώδης τὴν κεφαλήν. ὅτι δὲ αἰσχρὸς ἦν καὶ ἐν Θεσμοφοριαζούσαις δηλοῖ ; Σ Ar.~\textit{Thesm.}~168 ἐκ τούτου δηλονότι καὶ ἐν τοῖς Ὄρνισιν οὐκ ἄλλως κορυδῷ αὐτὸν ὡμοίου ἀλλ᾽ ἢ ὅτι μικρὸς καὶ αἰσχρός … τραγῳδίας ποιητὴς αἰσχρὸς τὸ πρόσωπον καὶ σιμὸς καὶ μικρὸς τὸ σῶμα ; Σ Ar.~\textit{Vesp.}~462 ὡς αὐστηροῦ κατὰ τὴν μελοποιίαν ὄντος. … ὡς τοῦ Φιλοκλέους ἀγρίου ὄντος ἐν τῇ μελοποιίᾳ. ὁ Σοφοκλῆς γὰρ ἡδύς· διὸ καὶ μέλιττα ἐκαλεῖτο. ἔπαιξεν οὖν, ὡσεὶ εἶπεν· εἴπερ τὴν πικρίαν αὐτοῦ καὶ τὴν σκληρότητα εἶχον, οὐκ ἂν αὐτοὺς ῥᾳδίως διέφυγες ; Σ Eur.~\textit{Andr.}~32 Φιλοκλῆς δὲ ὁ τραγῳδοποιὸς καὶ Θέογνις προεκδοθῆναί φασιν ὑπὸ Τυνδάρεω τὴν Ἑρμιόνην τῷ Ὀρέστῃ ; Suid.~η~358 Ἡ Μορσίμου τις ῥῆσιν ἐξεγράψατο: ἐπὶ τῶν ἀποβλήτων καὶ διαβεβλημένων. ὁ γὰρ Μόρσιμος τραγῳδὸς ἦν, ποιητὴς ὑπόψυχρος, Φιλοκλέους, ὀφθαλμῶν ἰατρός, μικρός ; Suid.~μ~1261 Μόρσιμος καὶ Μελάνθιος, ποιηταὶ τραγικοί. Μόρσιμος δὲ Φιλοκλέους τοῦ τραγικοῦ υἱός, ποιητὴς ψυχρός· ἦν δὲ καὶ ἰατρός ; Suid.~φ~378 Φιλοκλῆς, Πολυπείθους, Ἀθηναῖος, τραγικός, τοῖς χρόνοις μετὰ Εὐριπίδην. ἐπεκαλεῖτο δὲ Χολὴ διὰ τὸ πικρόν … Αἰσχύλου δὲ τοῦ τραγικοῦ ἦν ἀδελφιδοῦς καὶ ἔσχεν υἱὸν Μόρσιμον, τὸν τραγικόν ; Ar.~\textit{Eq.}~400–401 εἴ σε μὴ μισῶ, γενοίμην ἐν Κρατίνου κῴδιον, / καὶ διδασκοίμην προσᾴδειν Μορσίμου τραγῳδίᾳ ; Ar.~\textit{Pax}~801 ἑζομένη κελαδῇ, χορὸν δὲ μὴ ᾽χῃ Μόρσιμος ; Aristid.~\textit{Or.}~46 (§~256 Perseus) Σοφοκλῆς Φιλοκλέους ἡττᾶτο ἐν Ἀθηναίοις τὸν Οἰδίπουν, ὦ Ζεῦ καὶ θεοὶ … αἰσχύνη μὲν οὖν αὐτῷ τοσοῦτον ἀκοῦσαι, ὅτι βελτίων Φιλοκλέους ; Σ Ar.~\textit{Av.}~279 διαβάλλει δὲ τὸν Φιλοκλέα ὡς ὀξυκέφαλον ; Hyp. Soph.~\textit{OT} (ap.~Meineke, FCG II.2, p.~905 n.) χαριέντως δὲ τύραννον ἅπαντες αὐτὸν ἐπέγραφον, ὡς ἐξέχοντα πάσης τῆς Σοφοκλέους ποιήσεως, καίπερ ἡττηθέντα ὑπὸ Φιλοκλέους ; Suid.~μ~1262 Μόρσιμος, τραγῳδίας ποιητής, ὀφθαλμῶν ἰατρός, μικρός, ὑπόψυχρος ; Suid.~π~2464 Προκέφαλος· προκέφαλος ἦν ὁ Φιλοκλῆς, ὡς ἔποψ. ἦν δὲ κωμῳδίας ποιητής. εἴρηται ἐπὶ τῶν ὀξυκεφάλων ; Ar.~\textit{Nub.}~70 ὅταν σὺ μέγας ὢν ἅρμ᾽ ἐλαύνῃς πρὸς πόλιν, / ὥσπερ Μεγακλέης, ξυστίδ᾽ ἔχων ; Ath.~2.65f–66c καὶ γὰρ Φιλοκλῆς τε ἐγκέφαλόν φησιν οὐδ᾽ ἂν ἐγκέφαλον ἔσθων λίποι ; Suid.~α~4264 Ἀστυδάμας, ὁ πρεσβύτης, υἱὸς Μορσίμου τοῦ Φιλοκλέους, τραγικῶν ἀμφοτέρων, Ἀθηναῖος, τραγικός}
\lp{168 Φιλοκλέης (κορυδός)]}{Zen.~2.38 Ἀετοῦ γῆρας; κορύδου νεότης: παρόσον καὶ γηράσκων ἀετὸς ἀμείνων ἐστὶ παντὸς νεάζοντος ὄρνιθος ; \textit{idem} Diogenian.~1.56 (κορύδου pro παντός ; om. ὄρνιθος), Greg. Cypr.~1.4 (om. παρόσον … ὄρνιθος), Suid.~α~577 (ὁ ἀετός ; om. νεάζοντος), Tosi n°~1105 (om. παρόσον … ὄρνιθος) ; Diogenian.~4.56 Ἐν ἀμούσοις καὶ κόρυδος φθέγγεται ; \textit{idem} Zen.~3.77 (ἐν γὰρ ἀμούσοις καὶ κορυδὸς [κύκνου] φθέγγετ᾽ [ἀοιδότερος]), Suid.~ε~1112, Suid.~ε~1154 (καὶ γὰρ ἐν ἀμούσοις), Eust.~1072.40 (ad Il.~16.492) (καθὰ παροιμία, καὶ ὡς ἐν ἀμούσοις) ; Diogenian.~6.50 Μετὰ Μουσῶν κόρυδος: ἐπὶ τῶν ἀπαιδεύτων πεπαιδευμένοις συνόντων. Ὁμοία, Καρικὴ Μοῦσα. Καὶ, Κόρχορος ἐν λαχάνοις ; \textit{idem} Apostol.~11.28 (ἐπὶ τοῦ πεπαιδευμένου· ἢ κοινῇ ἐπὶ τοῦ συνόντος ἀπαιδεύτου τοῖς πεπαιδευμένοις. καθ᾽ ὑπόκρισιν δὲ ἀναγνωστέον pro ἐπὶ τῶν ἀπαιδεύτων … λαχάνοις) ; Erasm.~\textit{Adag.}~2.2.92 Inter indoctos etiam Corydus sonat. Ἐν ἀμούσοις καὶ Κόρυδος φθέγγεται … Corydus vilissimum aviculae genus, minimeque canorum, strepit tamen utcunque inter aves mutas, apud lusciniam canens ferri nequaquam posset ; Ar.~\textit{Av.}~471–476 οὐδ᾽ Αἴσωπον πεπάτηκας, / ὃς ἔφασκε λέγων κορυδὸν πάντων πρώτην ὄρνιθα γενέσθαι, / … τὸν πατέρ᾽ αὑτῆς ἐν τῇ κεφαλῇ κατορύξαι ; Ath.~6.241a–e ἦν δὲ καὶ ὁ Κόρυδος τῶν δι᾽ ὀνόματος παρασίτων … Εὐκράτην αὐτὸν καλεῖσθαι κυρίως φάσκων ; Σ Ar.~\textit{Av.}~475 ἐν τῇ κεφαλῇ κατορύξαι: ἐπεὶ λόφον ἔχει ἡ κορυδός. θηλυκῶς δὲ εἴρηκε τὴν κορυδόν. Πλάτων δὲ ἀρσενικῶς ; Hsch.~κ~3680 κορυδαλλός· ἢ πτηνὸν ζῷον. καὶ δῆμος τῆς Ἀττικῆς. ἢ κόρυδος παρ᾽ ἐνίων ; Greg. Cypr.~\textit{M.}~2.26 Ἄλλο γλαῦξ, ἄλλο κόρυδος φθέγγεται ; Erasm.~\textit{Adag.}~1.9.57 Aquilae senecta. Corydi juventa. Ἀετοῦ γῆρας. Κορύδου νεότης … De vivida viridique senecta, quae praestantior sit aliquorum juventa. Nam anus aquila praestat Corydo aviculae, etiam aetate integra ; Erasm.~\textit{Adag.}~3.4.96 Inter caecos regnat strabus … In regione caecorum rex est luscus. Inter indoctos, qui semidoctus est, doctissimus habetur ; Simon.~fr.~68~Bergk (Plu.~\textit{Mor.}~91e) ἐπεὶ δὲ πάσαις κορυδαλλίσι χρὴ λόφον ἐγγενέσθαι κατὰ τὸν Σιμωνίδην ; \textit{idem} Erasm.~\textit{Adag.}~3.3.67 (Oportet omnibus corydalis (lemme) ; πάσῃσι κορυδαλλοῖσι χρὴ λόφον ἐγγενέσθαι) ; Gal.~\textit{De\,simpl.\,med.\,temp.\,ac\,fac.}~11.1.37 ἔχει δὲ ἐπὶ τῆς κεφαλῆς τοῦτο τὸ ζῶον ὥσπερ τινὰ λόφον ἐκ τῶν τριχῶν αὐτοφυῆ, δι᾽ ὃν καὶ ὁ μῦθος ὃν Ἀριστοφάνης ὁ κωμικὸς ἔγραψεν ἐπλάσθη ; Apostol.~7.74 τοιαῦτα ἄττα καὶ Ἀθηναῖοι περὶ κορύδου τερατεύονται· ᾧ καὶ Ἀριστοφάνης ἀκολουθεῖ ἐν Ὄρνισι λέγων}
\lp{168 Φιλοκλέης (Κορυδέως εἰδεχθέστερος)]}{Zen.~4.59 Κορυδέως εἰδεχθέστερος: οὗτος ὁ Κορυδεὺς ἐπὶ δυσμορφίᾳ ἐκωμῳδεῖτο, καὶ οἱ παῖδες αὐτοῦ, ὧν ἕνα τὸν Ἀρχέστρατον φασίν ; \textit{idem} Com. adesp. fr.~21~Kock (Κορυδέως … εἰδεχθέστερος), Diogenian.~5.57 (ἐπὶ τῶν δυσμόρφων· οὗτος γὰρ ἐπὶ ἀμορφίᾳ διεβάλλετο), Apostol.~9.92 (ἐπὶ δυσμορφίᾳ οὗτος διεβάλλετο), Hsch.~κ~3681 (Κορυδεύς· ὄνομα κύριον. οὗτος ἐκωμῳδεῖτο ἐπὶ δυσμορφίᾳ, καὶ οἱ παῖδες αὐτοῦ) ; Erasm.~\textit{Adag.}~3.2.8 Corytheo deformior. Κορυθέως εἰδεχθέστερος … Ajunt Corytheum quempiam fuisse hominem insignite stultum ac vecordem, ad hoc forma infelicissima foedissimaque, cui liberi item fuerint omnes patrem referentes … Refertur adagium à Zenodoto}
\lp{168 αἰσχρὸς ὢν αἰσχρῶς ποιεῖ]}{Soph.~\textit{El.}~621 αἰσχροῖς γὰρ αἰσχρὰ πράγματ᾽ ἐκδιδάσκεται ; Hsch.~α~2151 αἰσχρά· μυσαρά, κακά. Εὐριπίδης Τρῳάσιν ; Eur.~\textit{Tro.}~1177 ὀστέων ῥαγέντων φόνος, ἵν᾽ αἰσχρὰ μὴ λέγω ; \textit{Il.}~2.216 αἴσχιστος δὲ ἀνὴρ ὑπὸ Ἴλιον ἦλθε ; Hsch.~α~2148 αἴσχιστος· αἰσχρότατος (Β 216) … κακός, μυσαρός ; Hsch.~α~2150 αἶσχος· αἰσχρότης, κακότης, κακία (Eur. Tr. 1114). λοιδορία. μῦσος ; App. Prov.~3.19 Θερσίτειον βλέμμα, καὶ, Θερσίτειον εἴδωλον: παροιμία ἐπὶ τῶν πάνυ δυσειδῶν, ἐπεὶ καὶ ὁ Θερσίτης τοιοῦτος ἦν ; \textit{idem} Suid.~θ~257 (Θερσίτης· ὄνομα κύριον. καὶ Θερσίτειον βλέμμα) ; Erasm.~\textit{Adag.}~4.3.80 Thersitae facies. Θερσίτειον βλέμμα. i. e. Thersitica facies. De prodigiose deformi dici solitum, quod Homerus scripserit hunc omnium, qui ad Trojam venissent, foedissimum fuisse … Si corpus adspicias, Nireus est, si animum, plusquam Θερσίτειον βλέμμα, reperies. Adagii mentionem fecit Suidas ; Max.Tyr.~1.5 ἀσθενέστερος μὲν ὢν τοῦ Ἐπειοῦ, αἰσχίων δὲ τοῦ Θερσίτου, μικρότερος δὲ τοῦ Τυδέως ; Epict.~2.23.32 οἱ λέγοντες μηδεμίαν εἶναι παραλλαγὴν κάλλους πρὸς αἶσχος. εἶτα ὁμοίως ἦν κινηθῆναι τὸν Θερσίτην ἰδόντα καὶ τὸν Ἀχιλλέα ; Luc.~\textit{DMort.}~25.1 Νιρεὺς καὶ Θερσίτης. πότερος οὖν ὁ Νιρεὺς καὶ πότερος ὁ Θερσιτης; οὐδέπω γὰρ τοῦτο δῆλον ; Max.Tyr.~26.5 Πάλιν αὖ Θερσίτης πεποίηται αὐτῷ αἰσχρὸς ἰδεῖν, φωνὴν ἐπεσβόλος, γνώμην ἄτακτος ; Hsch.~α~2154 αἰσχροῖς· τοῖς αἰσχύνην ἐνεγκεῖν δυναμένοις}
\lp{169 ὁ δ᾽ αὖ]}{Ar.~\textit{Ach.}~552 τὸ νεώριον δ᾽ αὖ κωπέων πλατουμένων ; Ar.~\textit{Av.}~516 ἡ δ᾽ αὖ θυγάτηρ γλαῦχ᾽, ὁ δ᾽ Ἀπόλλων ὥσπερ θεράπων ἱέρακα ; Ar.~\textit{Av.}~582 οἱ δ᾽ αὖ κόρακες τῶν ζευγαρίων, οἷσιν τὴν γῆν καταροῦσιν ; Ar.~\textit{Lys.}~91 Κορινθία δ᾽ αὖ ; Ar.~\textit{Lys.}~426 ποῖ δ᾽ αὖ σὺ βλέπεις ; Ar.~\textit{Eccl.}~667 οὐδ᾽ αὖ κλέπτης οὐδεὶς ἔσται}
\lp{169 Ξενοκλέης]}{Ar.~\textit{Vesp.}~1510–1511 ὁ πινοτήρης οὗτός ἐστι τοῦ γένους, / ὁ σμικρότατος, ὃς τὴν τραγῳδίαν ποιεῖ ; Ar.~\textit{Thesm.}~441 ὥστ᾽ ἂν εἰ λέγοι παρ᾽ αὐτὴν / Ξενοκλέης ὁ Καρκίνου, δοκεῖν / ἂν αὐτόν, ὡς ἐγᾦμαι, / πᾶσιν ὑμῖν / ἄντικρυς μηδὲν λέγειν ; Ar.~\textit{Ran.}~86 ὁ δὲ Ξενοκλέης; — ἐξόλοιτο νὴ Δία ; Pl. Com. fr.~134~Kock Ξενοκλῆς ὁ δωδεκαμήχανος ὁ Καρκίνου παῖς τοῦ θαλαττίου ; Σ Ar.~\textit{Pax}~792 μηχανοδίφας : ἀπὸ μέρους τοῦτο· Ξενοκλῆς γὰρ ὁ Καρκίνου δοκεῖ μηχανὰς καὶ τερατείας εἰσάγειν ἐν τοῖς δράμασι ; Σ Ar.~\textit{Ran.}~86 Ξενοκλῆς : υἱὸς Καρκίνου, ὃς κωμῳδεῖται ὡς ἄξεστος ἐν τῇ ποιήσει (καὶ ἀλληγορικός) ; Σ Ar.~\textit{Thesm.}~441 διὰ τὸ εὐμήχανον ἐν τοῖς δράμασι λέγεσθαι τὸν Ξενοκλέα πρὸς τὰς τῆς γυναικὸς μηχανὰς παρελήφθη … ἤτοι δὲ διαβάλλει αὐτὸν ὡς ἀφυᾶ ; Suid.~α~2802 Ἄξεστος· Ξενοκλῆς ὁ Καρκίνου ἐκωμῳδεῖτο ὡς ἄξεστος ποιητὴς καὶ ἀλληγορικός ; Ar.~\textit{Nub.}~1264–1265 ὦ σκληρὲ δαῖμον, ὦ τύχαι θραυσάντυγες / ἵππων ἐμῶν, ὦ Παλλὰς ὥς μ᾽ ἀπώλεσας ; Ael.~\textit{VH}~2.8 ἀντηγωνίσαντο ἀλλήλοις Ξενοκλῆς καὶ Εὐριπίδης. καὶ πρῶτός γε ἦν Ξενοκλῆς, ὅστις ποτὲ οὗτός ἐστιν ; Σ Ar.~\textit{Pax}~795 δρᾶμα γὰρ ποιήσας ὁ Καρκίνος ἢ ὁ Ξενοκλῆς παρῆλθε μέγα καυχώμενος ἐπ᾽ αὐτῷ, καὶ παρελθὼν ἡττήθη ; Σ Ar.~\textit{Vesp.}~1509 εἰς σμικρὸν τὸν Ξενοκλέα. καὶ τὸ φαλάγγιον μικρὸν καὶ συνεστραμμένον ; Hsch.~α~5617 ἀξέστους· τραχείας. Σοφοκλῆς Σίνωνι ; Hsch.~μ~1301 μηχαναί· τέχναι. ἐπίνοιαι. βουλεύματα ; Hsch.~μ~1303 μηχανορράφος· κατασκευαστής, ἐπινοητὴς κακῶν}
\lp{169 Ξενοκλέης (ὁ Καρκίνου)]}{Ar.~\textit{Pax}~788–790 ὄρτυγας οἰκογενεῖς γυλιαύχενας ὀρχηστὰς / ναννοφυεῖς σφυράδων ἀποκνίσματα μηχανοδίφας ; \textit{idem} Suid.~κ~396 (praem. Καρκίνος, ποιητὴς Ἀττικός. σκώπτει δέ τις Καρκίνου παῖδας οὕτως) ; Σ Ar.~\textit{Nub.}~1261 οἱ δὲ Καρκίνου παῖδες Ξενοκλῆς καὶ Ξενότιμος, καὶ ὁ μὲν χορευτής, Ξενοκλῆς δὲ τραγῳδίας ποιητής ; Hsch.~π~2318 πινοτήρης· πινοφύλαξ λεγόμενος. πῖνα δὲ εἶδος ὀστρέου ; Ar.~\textit{Nub.}~1261 τῶν Καρκίνου τις δαιμόνων ἐφθέγξατο ; Ar.~\textit{Vesp.}~1505 ἕτερος τραγῳδὸς Καρκινίτης ἔρχεται ; Ar.~\textit{Pax}~781–786 ἢν δέ σε Καρκίνος ἐλθὼν / ἀντιβολῇ μετὰ τῶν παίδων χορεῦσαι, / μήθ᾽ ὑπάκουε μήτ᾽ ἔλθῃς / συνέριθος αὐτοῖς ; Men.~fr.~525~Kock (=~415~K.-A.) Καρκίνου ποιήματα ; Σ Ar.~\textit{Pax}~864 κακοδαιμονέστερος, ἐν εἰρωνείᾳ· τῶν Καρκίνου στροβίλων: στροβίλους φησὶ τούτου ⟨τοὺ⟩ς ⟨παῖδας⟩ διὰ τὸ βραχὺ τοῦ σώματος … τῶν Καρκίνου στροβίλων: διὰ τὸ ἐν τῇ ὀρχήσει στροβεῖσθαι ; Σ Ar.~\textit{Vesp.}~1510 πιννοτήρης δὲ οὐκ ἄλλως, ἀλλ᾽ ἐπεὶ καρκίνιόν τί ἐστι σύννομον καὶ ἀεὶ ταύτῃ προσδεχόμενον. Σοφοκλῆς Ἀμφιαράῳ· ὁ πιννοτήρης τοῦδε μάντεως χοροῦ ; Harp. s.v.~Καρκίνος ἔστι δὲ καὶ κύριον ὄνομα … περὶ δὲ τοῦ τῆς τραγῳδίας ποιητοῦ τοῦ Ξενοκλέους υἱοῦ Λυσίας ἐν τῷ πρὸς Μνησίμαχόν φησι ; Zen.~3.77 Ἐν γὰρ ἀμηχανίῃ καὶ Καρκίνος ἔμμορε τιμῆς ; Suid.~σ~1208 Στρόβιλοι: Καρκίνου τοῦ ποιητοῦ υἱοί, ὡς Ἀριστοφάνης· οὓς καὶ γυλιοτραχήλους εἴρηκε. καὶ παροιμία· εὐδαιμονέστερος τῶν Καρκίνου στροβίλων. ἀντὶ τοῦ κακοδαιμονέστερος· ἐν εἰρωνείᾳ ; Erasm.~\textit{Adag.}~1.2.68 Fortunatior Strobilis Carcini … Ironia proverbialis in infelicem \& improbum. Carcinus Poeta quispiam fuit, cujus filii dicti sunt Strobili ab Aristophane … Adagium recensetur apud Suidam \& Zenodotum, exstat apud Aristophanem in Pace ; Phot.~κ~193 Καρκίνου ποιήματα· Μένανδρος Ψευδηρακλεῖ (fr. 415 K.-A.)· αἰνιγματώδη· ὁ γὰρ Καρκίνος Ὀρέστην ἀπὸ †Ἡλίου† ἀναγκαζόμενον ὁμολογῆσαι, ὅτι ἐμητροκτόνησεν, ἐποίησε δι᾽ αἰνιγμάτων ἀποκρινόμενον ; \textit{idem} Suid.~κ~397 (ἀντὶ τοῦ αἰνιγματώδη ; Ἰλίου ; om. ὁμολογῆσαι … ἀποκρινόμενον) ; Ar.~\textit{Pax}~864 εὐδαιμονέστερος φανεῖ τῶν Καρκίνου στροβίλων ; Th.~2.23.2 ἐστρατήγει δὲ Καρκίνος τε ὁ Ξενοτίμου ; Hsch.~κ~832 καρκίνος· τὸ νῦν καρκίνωμα λεγόμενον. καὶ ὑποδήματα κοῖλα. καὶ δεσμός τις. καὶ ἡ πυράγρα. ἢ εἶδος ζῴου. καὶ ἄστρον ; Suid.~κ~394 Καρκίνος, Ἀκραγαντῖνος, τραγικός. καὶ Καρκίνος, Θεοδέκτου ἢ Ξενοκλέους, Ἀθηναῖος, τραγικός. δράματα ἐδίδαξεν ρξ΄, ἐνίκησε δὲ α΄}
\lp{169 κακὸς κακῶς]}{Ar.~\textit{Nub.}~553–554 Εὔπολις μὲν τὸν Μαρικᾶν πρώτιστον παρείλκυσεν / ἐκστρέψας τοὺς ἡμετέρους Ἱππέας κακὸς κακῶς ; Eur.~\textit{Med.}~1386 σὺ δ᾽, ὥσπερ εἰκός, κατθανῇ κακὸς κακῶς ; Zen.~4.82 Κακοῦ κόρακος κακὸν ὠόν ; \textit{idem} Diogenian.~5.39 (add. ἔφυ· ἐπὶ τῶν ἐκ κακῶν γονέων ὄντων), Greg. Cypr.~2.75, Tosi n°~1664 (ᾠόν) ; Ar.~\textit{Ran.}~1008–1010 τίνος οὕνεκα χρὴ θαυμάζειν ἄνδρα ποιητήν; / δεξιότητος καὶ νουθεσίας, ὅτι βελτίους τε ποιοῦμεν / τοὺς ἀνθρώπους ἐν ταῖς πόλεσιν. — τοῦτ᾽ οὖν εἰ μὴ πεποίηκας ; Eub.~fr.~116–117~Kock κακὸς κακῶς ἀπόλοιθ᾽ ὅστις γυναῖκα δεύτερος ἔγημε ; Eur.~\textit{Tro.}~446 ἦ κακὸς κακῶς ταφήσῃ νυκτός, οὐκ ἐν ἡμέρᾳ ; Soph.~\textit{Aj.}~1177 κακὸς κακῶς ἄθαπτος ἐκπέσοι χθονός ; Soph.~\textit{OT}~248 κακὸν κακῶς νιν ἄμορον ἐκτρῖψαι βίον ; Plu.~\textit{Mor.}~218f φαῦλα δὲ ποιήματα γράφοντος, τί δήποτε, ὦ Περίανδρε εἶπεν ἀντὶ χαρίεντος ἰατροῦ κακὸς ποιητὴς καλεῖσθαι ἐπιθυμεῖς ; Hsch.~κ~333 κακοί· ἄνανδροι, δειλοί, δυστυχεῖς, κακοδαίμονες ; Erasm.~\textit{Adag.}~1.9.25 Mali corvi malum ovum. Κακοῦ κόρακος κακὸν ᾠόν. i. e. Mali corvi malum ovum. Apte usurpabitur, quoties à malo praeceptore discipulus malus proficiscitur: ex improbo patre filius improbus … denique facinus sceleratum ab homine scelesto ; Ar.~\textit{Plu.}~65 εἰ μὴ φράσεις γάρ— — ἀπό σ᾽ ὀλῶ κακὸν κακῶς ; Ar.~\textit{Plu.}~418 ἐγὼ γὰρ ὑμᾶς ἐξολῶ κακοὺς κακῶς ; Eur.~\textit{Cyc.}~268 ἢ κακῶς οὗτοι κακοὶ / οἱ παῖδες ἀπόλοινθ᾽ ; Hsch.~κ~352 κακός· δειλός … δεινός. πονηρός ; Hsch.~κ~375 κακῶς· δεινῶς. ταπεινῶς ; Ammon.~\textit{Diff.}~263 Nickau κακὸς πονηροῦ διαφέρει ὥσπερ ὁ ἄκακος τοῦ ἀγαθοῦ. κακὸς μὲν γὰρ ὁ πανοῦργος, πονηρὸς δὲ ὁ δραστικὸς κακοῦ}
\lp{170 Θέογνις]}{Ar.~\textit{Ach.}~11 ὁ δ᾽ ἀνεῖπεν, εἴσαγ᾽ ὦ Θέογνι τὸν χορόν ; \textit{idem} Suid.~α~2384 (s.v. ἀνεῖπεν ; ἀνεῖπ᾽) ; Ar.~\textit{Ach.}~138–140 εἰ μὴ κατένειψε χιόνι τὴν Θρᾴκην ὅλην / καὶ τοὺς ποταμοὺς ἔπηξ᾽, ὑπ᾽ αὐτὸν τὸν χρόνον, / ὅτ᾽ ἐνθαδὶ Θέογνις ἠγωνίζετο ; Σ Ar.~\textit{Ach.}~11 Θέογνις δὲ οὗτος τραγῳδίας ποιητὴς πάνυ ψυχρός, ἐκ τῶν τριάκοντα, ὃς καὶ Χιὼν ἐλέγετο ; \textit{idem} Suid.~θ~137 (τραγῳδοποιητής ; λ΄) ; Σ Ar.~\textit{Ach.}~140 οὗτος ὁ Θέογνις τραγῳδίας ποιητὴς ψυχρός. κωμῳδῶν οὖν αὐτὸν παρίστησιν ἡμῖν τὴν πολλὴν χιόνα, διὰ τὴν τούτου περὶ τὰ ποιήματα ψυχρότητα. τὸν ψυχρὸν οὖν τῷ ψυχρῷ παρέβαλεν. πάνυ δὲ χαριέντως ἔσκωψε ; \textit{idem} Suid.~ει~148 (praem. εἰ μὴ κατένιψε τὴν Θρᾴκην χιόνι πολλῇ … ἠγωνίζετο ; ἦν ; τῆς … ψυχρότητος ; om. τὸν ψυχρὸν … ἔσκωψε), Suid.~ψ~178 (praem. ἡνίκα Θέογνις ἠγωνίζετο, κατένιψε … χρόνον ; παρίστησι ; τῆς … ψυχρότητος ; om. ἡμῖν et τὸν ψυχρὸν … ἔσκωψε) ; Hor.~\textit{Sat.}~2.5.40–41 seu pingui tentus omaso / Furius hibernas cana nive conspuet Alpes ; X.~\textit{Hell.}~2.3.2 καὶ ᾑρέθησαν οἵδε· … Ἐρατοσθένης, Χαρικλῆς, Ὀνομακλῆς, Θέογνις ; Arist.~\textit{Rh.}~3.11.11 (1413a1) οἷον ἡ ἀσπίς, φαμέν, ἐστι φιάλη Ἄρεως, καὶ ⟨τὸ⟩ τόξον φόρμιγξ ἄχορδος ; Demetr.~\textit{Eloc.}~85 ὡς ὁ Θέογνις παρατίθεται τῷ τόξῳ φόρμιγγα ἄχορδον ἐπὶ τοῦ τῷ τόξῳ βάλλοντος ; Harp. s.v.~Θέογνις Θέογνις· Λυσίας ἐν τῷ κατ᾽ Ἐρατοσθένους. οὗτος εἷς ἦν τῶν παρ᾽ Ἀθηναίοις λ τυράννων, καθά φασιν ἄλλοι τε καὶ Ξενοφῶν ἐν β Ἑλληνικῶν ; Lys.~12.6 Θέογνις γὰρ καὶ Πείσων ἔλεγον ἐν τοῖς τριάκοντα περὶ τῶν μετοίκων}
\lp{170 ψυχρὸς ὢν ψυχρῶς ποιεῖ]}{Ar.~\textit{Av.}~935 ἔχε τὴν σπολάδα· πάντως δέ μοι ῥιγῶν δοκεῖς ; Ar.~\textit{Av.}~950–953 κλῇσον ὦ χρυσόθρονε τὰν τρομερὰν κρυεράν· / νιφόβολα πεδία πολύπορά τ᾽ ἤλυθον ἀλαλάν ; Ar.~\textit{Av.}~954–955 νὴ τὸν Δί᾽ ἀλλ᾽ ἤδη πέφευγας ταυταγὶ / τὰ κρυερὰ τονδὶ τὸν χιτωνίσκον λαβών ; Ar.~\textit{Av.}~1385 ἀεροδονήτους καὶ νιφοβόλους ἀναβολάς ; Ar.~\textit{Thesm.}~848 οὐ τὸν Παλαμήδην ψυχρὸν ὄντ᾽ αἰσχύνεται ; Eup.~fr.~244~Kock τοῦτ᾽ ἔστι σοι τὸ σκῶμμα ἀσελγὲς καὶ Μεγαρικὸν καὶ σφόδρα ψυχρόν ; Alex.~fr.~179~Kock πρᾶγμα δ᾽ ἔστι μοι μέγα φρέατος ἔνδον ψυχρότερον Ἀραρότος ; Macho ap. Ath.~13.580a παρὰ Γναθαίνῃ Δίφιλος πίνων ποτὲ / ψυχρὸν γ᾽, ἔφη, τἀγγεῖον, ὦ Γνάθαιν᾽, ἔχεις. / τῶν σῶν γάρ, εἶπεν, ἐπιμελῶς, ὦ Δίφιλε, / εἰς αὐτὸ γ᾽ αἰεὶ δραμάτων ἐμβάλλομεν ; Theophil.~fr.~4~Kock ψυχρός ἐστιν, ἄπαγε, φησί· ῥητόρων οὐ γεύομαι ; Pl.~\textit{Euthd.}~284e καὶ τοὺς μεγάλους, ἔφη ὁ Εὐθύδημος, μεγάλως λέγουσι καὶ τοὺς θερμοὺς θερμῶς;—μάλιστα δήπου, ἔφη ὁ Κτήσιππος· τοὺς γοῦν ψυχροὺς ψυχρῶς λέγουσί τε καὶ φασὶν διαλέγεσθαι ; X.~\textit{Cyr.}~8.4.23 πολύ γ᾽, ἔφη, μάλιστα τούτου σε, ὦ Κῦρε, ζηλῶ ἐν τῇ βασιλείᾳ. τίνος; ἔφη ὁ Κῦρος. ὅτι δύνασαι καὶ ψυχρὸς ὢν γέλωτα παρέχειν ; Arist.~\textit{Rh.}~3.3.1 τὰ δὲ ψυχρὰ ἐν τέτταρσι γίγνεται κατὰ τὴν λέξιν ; Plu.~\textit{Mor.}~853b (Comp. Ar. et Men.) τούτοις γὰρ ὁ μὲν μετὰ τοῦ προσήκοντος λόγου καὶ ὀλιγάκις χρῆται ἐπιμελείας αὐτὰ ἀξιῶν, ὁ δὲ καὶ πολλάκις καὶ οὐκ εὐκαίρως καὶ ψυχρῶς ; Demetr.~\textit{Eloc.}~114 ὁρίζεται δὲ τὸ ψυχρὸν Θεόφραστος οὕτως, ψυχρόν ἐστι τὸ ὑπερβάλλον τὴν οἰκείαν ἀπαγγελίαν ; Hsch.~η~687 ἠπίαλος· ῥῖγος πρὸ πυρετοῦ. ἐκαλοῦντο δὲ οὕτω καὶ οἱ ψυχροί ; Hsch.~ψ~300 ψυχροκομψώματα· ψυχροφαντάσματα ; Hsch.~ψ~303 ψυχρολόγος· μηδὲν χρήσιμον λέγων ; Suid.~λ~808 Λύκις: λέγεται καὶ Λύκος. κωμῳδεῖται δὲ ὡς ψυχρὸς ποιητής … οὗτοι οὖν οἱ τρεῖς κωμικοὶ ὑπόψυχροι ; Suid.~μ~495 Μέλητος: τραγῳδίας ποιητὴς ὑπόψυχρος … κωμῳδεῖται δὲ ὡς ψυχρὸς ἐν τῇ ποιήσει καὶ ὡς πονηρὸς τὸν τρόπον ; Suid.~ψ~176 Ψυχρός: ἀσθενής. ἦν δὲ ὁ Μόρσιμος Φιλοκλέους τοῦ ποιητοῦ υἱός, τραγῳδίας ποιητὴς ψυχρός. ἦν δὲ καὶ ἰατρός ; Ar.~\textit{Av.}~946 ξυνῆχ᾽ ὅτι βούλει τὸν χιτωνίσκον λαβεῖν ; Ar.~\textit{Thesm.}~47 θηρῶν τ᾽ ἀγρίων πόδες ὑλοδρόμων / μὴ λυέσθων ; Ar.~\textit{Thesm.}~67–69 καὶ γὰρ μελοποιεῖν ἄρχεται· χειμῶνος οὖν / ὄντος κατακάμπτειν τὰς στροφὰς οὐ ῥᾴδιον, / ἢν μὴ προίῃ θύρασι πρὸς τὸν ἥλιον ; Ar.~\textit{Ran.}~12–14 τί δῆτ᾽ ἔδει με ταῦτα τὰ σκεύη φέρειν, / εἴπερ ποιήσω μηδὲν ὧνπερ Φρύνιχος / εἴωθε ποιεῖν καὶ Λύκις κἀμειψίας ; Pl.~\textit{Lg.}~7.802d τραφεὶς δ᾽ ἐν τῇ κοινῇ καὶ γλυκείᾳ, ψυχρὰν καὶ ἀηδῆ ; X.~\textit{Symp.}~6.7 εἰ δὲ ψυχρὰ λέγω, σὺ αἴτιος, ἔφη, πράγματά μοι παρέχων ; D.~18.256 καί μου πρὸς Διὸς μηδεμίαν ψυχρότητα καταγνῷ μηδείς ; D.~21.112 καὶ τἀδικήμαθ᾽ ἕωλα τὰ τούτων ὡς ὑμᾶς καὶ ψύχρ᾽ ἀφικνεῖται ; Arist.~\textit{Rh.}~3.3.3 (1406a18–32) διὸ τὰ Ἀλκιδάμαντος ψυχρὰ φαίνεται· οὐ γὰρ ὡς ἡδύσματι χρῆται ἀλλ᾽ ὡς ἐδέσματι τοῖς ἐπιθέτοις … καὶ οὐ τὸ σῶμα παρήμπισχεν ἀλλὰ τὴν τοῦ σώματος αἰσχύνην … διὸ ποιητικῶς λέγοντες τῇ ἀπρεπείᾳ τὸ γελοῖον καὶ τὸ ψυχρὸν ἐμποιοῦσι ; Thphr.~\textit{Char.}~2.4 καὶ σκώψαντι ψυχρῶς ἐπιγελάσαι ; D.H.~\textit{Isoc.}~3 σχηματίζει τε φορτικῶς καὶ τὰ πολλὰ γίνεται ψυχρὸς ; Longin.~4.1 Θατέρου δὲ ὧν εἴπομεν, λέγω δὲ τοῦ ψυχροῦ, πλήρης ὁ Τίμαιος ; Σ Ar.~\textit{Ach.}~13 ὁ δὲ Δεξίθεος ἄριστος κιθαρῳδὸς καὶ Πυθιονίκης. οἱ δὲ ψυχρὸν αὐτὸν εἶναι φασιν ; Σ Ar.~\textit{Av.}~950–951 950. διαβάλλει τοὺς ποιητὰς ὡς τῷ κόμπῳ τῶν λέξεων ἀπατῶντας, μηδὲν δὲ λέγοντας … 951. τρομερὰν, κρυεράν … κρυερὰν δέ, τὴν κρύους μεταλαμβάνουσαν, διὰ τὸ ἐπ᾽ αὐτὴν καταράσσειν καὶ τὴν δρόσον καὶ τὴν ψυχρότητα ; Hsch.~β~440 βεβρός· ψυχρός, τετυφωμένος ; Hsch.~ε~7736 ἕωλον· χθιζόν, ἐχθεσινόν. διὰ δασέος. [δὲ] ψυχρόν ; Hsch.~ψ~302 ψυχρολογία· ψευδολογία ; Hsch.~ψ~305 ψυχρός· ἀσθενής ; Suid.~ψ~175 Ψυχρολογία: ψευδολογία, ἀκαιρολογία. καὶ Ψυχρολόγος ἄνθρωπος, ὁ μηδὲν ἀγαθὸν λέγων ; Ar.~\textit{Plu.}~263 ψυχροῦ βίου καὶ δυσκόλου ζήσειν ἀπαλλαγέντας ; Soph.~\textit{Ant.}~88 θερμὴν ἐπὶ ψυχροῖσι καρδίαν ἔχεις ; Pi.~fr.~123 (Ath.~13.601d) ἐξ ἀδάμαντος ἢ σιδάρου κεχάλκευται μέλαιναν καρδίαν ψυχρᾷ φλογί ; Cic.~\textit{De\,or.}~2.260 Haec aut frigida sunt aut tum salsa, cum aliud est exspectatum ; Hsch.~γ~301 γελανδρόν· ψυχρόν ; Hsch.~ψ~299 ψυχρά· τὴν ἣν ἡμεῖς λεπτὴν λέγομεν ; Hsch.~ψ~301 ψυχρὸν κακόν· ψύξεως αἴτιον}
\lp{170 ψυχρός (Ἰαλέμου ψυχρότερος)]}{Zen.~4.39 Ἰαλέμου ψυχρότερος: παροιμία, ὡς εἴ τις εἴποι, Γυμνότερος τοῦ κακοδαίμονος· ἢ ὅτι οἱ θρηνοῦντες τῷ ἰαλέμῳ χρῶνται ; \textit{idem} Apostol.~8.97 (Ἰάλεμον τὸν θρῆνόν φασι. λέγουσι δὲ αὐτὸν υἱὸν Καλλιόπης γενέσθαι κακοδαίμονα καὶ ὑστερούμενον καὶ ὀρφανόν pro παροιμία … χρῶνται) ; Diogenian.~2.74 Ἆκις ποταμός: ἐπὶ τῶν ἄγαν ψυχρῶν. Ψυχρὸς γάρ ἐστιν οὗτος ; Hsch.~ι~29 ἰαλεμώδη· ψυχρά, οὐδενὸς ἄξια. καὶ παροιμία· ψυχρότερος καὶ οἰκτρότερος Ἰαλέμου ; Hsch.~λ~217 Λακρατίδης· Ἀριστοφάνης φησὶ παλαιὸν Λακρατίδην, τὰ ψυχρὰ βουλόμενος δηλοῦν· ψυχροὶ γὰρ οἱ γέροντες ; Synag.~ι~2 ἰαλέμῳ· θρήνῳ. ἰάλεμος γὰρ θρήνου εἶδος. — λέγεται δὲ καὶ Ἰάλεμος υἱὸς Καλλιόπης καὶ ὁ καλὸς καὶ ὀρφανός. ἔνιοι δὲ τὸν κακὸν δαίμονα. καὶ ἰαλεμώδη τὰ ψυχρὰ καὶ οὐδενὸς λόγου ἄξια ; App. Prov.~3.24 Ἰάλεμος: εἶδος θρήνου, ἡ ἐπὶ τοῖς ἀπολωλόσιν ἀνία, καὶ ὁ υἱὸς Καλλιόπης, καὶ ὁ ὀρφανός ; Erasm.~\textit{Adag.}~2.10.86 Ialemo frigidior. Ἰαλέμου ψυχρότερος. i. e. Ialemo frigidior. De cantione, oratione, argumento, joco aut ejusmodi re quapiam supra modum frigida, recte dicetur ; Demetr.~\textit{Eloc.}~119 Καὶ καθόλου ὁποῖόν τί ἐστιν ἡ ἀλαζονεία, τοιοῦτον καὶ ἡ ψυχρότης … καὶ ὁποῖόν τι τὸ ἐν τῇ παροιμίᾳ κοσμούμενον ὕπερον, τοιοῦτόν τί ἐστι καὶ τὸ ἐν τῇ ἑρμηνείᾳ ἐξηρμένον ἐν μικροῖς πράγμασιν ; Diogenian.~3.88 Γέλως Μεγαρικός: ἐπὶ τῶν ἀώρως θρυπτομένων. Ἤκμασε γὰρ ἡ Μεγαρικὴ κωμῳδία ἐπὶ χρόνον ; App. Prov.~5.4 Ὕπερον κοσμεῖς: ὅμοιον τῷ, Χύτραν ποικίλλεις}
\lp{171 ἅπασ᾽ ἀνάγκη]}{Ar.~\textit{Pax}~372–373 νῦν ἆρά με — ἅπασ᾽ ἀνάγκη ᾽στ᾽ ἀποθανεῖν ; Ar.~\textit{Eccl.}~1029 καὶ ταῦτ᾽ ἀνάγκη μοὐστί; — Διομήδειά γε ; Polyzel.~fr.~3~Kock τριῶν κακῶν γοῦν ἦν ἑλέσθ᾽ αὐτῷ τι πᾶσ᾽ ἀνάγκη ; Soph.~\textit{El.}~1497 ἦ πᾶσ᾽ ἀνάγκη τήνδε τὴν στέγην ἰδεῖν ; Pl.~\textit{R.}~2.381c ἅπασα, ἔφη, ἀνάγκη ἔμοιγε δοκεῖ ; Pl.~\textit{Tht.}~203d προγιγνώσκειν τὰ στοιχεῖα ἅπασα ἀνάγκη τῷ μέλλοντί ποτε γνώσεσθαι συλλαβήν ; Ar.~\textit{Thesm.}~167 ὅμοια γὰρ ποιεῖν ἀνάγκη τῇ φύσει ; Eur.~\textit{Hec.}~396 πολλή γ᾽ ἀνάγκη θυγατρὶ συνθανεῖν ἐμέ ; Eur.~\textit{Med.}~1013 πολλή μ᾽ ἀνάγκη, πρέσβυ ; Aesch.~\textit{Supp.}~439–440 ἢ τοῖσιν ἢ τοῖς πόλεμον αἴρεσθαι μέγαν / πᾶσ᾽ ἔστ᾽ ἀνάγκη ; Soph.~\textit{OT}~986 ζῇ, πᾶσ᾽ ἀνάγκη, κεἰ καλῶς λέγεις, ὀκνεῖν ; Pl.~\textit{R.}~4.441d καὶ τοῦτο πᾶσα ἀνάγκη ; D.~20.26 οὔκουν ὁ πολλὰ κεκτημένος, οὗτος, ὅστις ἂν ᾖ, πόλλ᾽ εἰς ταῦτα συντελεῖ; πᾶσ᾽ ἀνάγκη ; Eur.~\textit{Alc.}~378 πολλή μ᾽ ἀνάγκη, σοῦ γ᾽ ἀπεστερημένοις ; Eur.~\textit{Phoen.}~1674 πολλή σ᾽ ἀνάγκη· ποῖ γὰρ ἐκφεύξῃ λέχος ; Soph.~\textit{Trach.}~295 πολλή ᾽στ᾽ ἀνάγκη τῇδε τοῦτο συντρέχειν ; Damocr. ap. Gal.~\textit{De\,antid.}~1.15 Ἅπασ᾽ ἀνάγκη ; Hdt.~5.52.2 ἐπ᾽ ᾧ πύλαι τε ἔπεισι, τὰς διεξελάσαι πᾶσα ἀνάγκη ; Pl.~\textit{Epin.}~977d ἀριθμὸν μὲν ἀνάγκη πᾶσα ὑποτίθεσθαι ; Pl.~\textit{Epin.}~982a ἡ ψυχῆς δὲ ἀνάγκη νοῦν κεκτημένης ἁπασῶν ; Pl.~\textit{Grg.}~507a ἡ δέ γε κοσμία σώφρων;—πολλὴ ἀνάγκη ; Hsch.~α~4234 ἀνάγκη· ἡ δικαστικὴ κλεψύδρα}
\lp{171 ἀνάγκη (ἀνάγκῃ οὐδὲ θεοὶ μάχονται)]}{Eur.~\textit{Hel.}~513–514 λόγος γάρ ἐστιν οὐκ ἐμός, σοφὸν δ᾽ ἔπος, / δεινῆς ἀνάγκης οὐδὲν ἰσχύειν πλέον ; \textit{idem} Zen.~3.9 (ἰσχυρότερον ; add. ἀνάγκη γὰρ καὶ θεοὺς βιάζεται) ; Pl.~\textit{Prt.}~345d (Simon.) πάντας δ᾽ ἐπαίνημι καὶ φιλέω / ἑκὼν ὅστις ἕρδῃ / μηδὲν αἰσχρόν· ἀνάγκῃ δ᾽ οὐδὲ θεοὶ μάχονται ; \textit{idem} Zen.~1.85 (om. δ᾽), Greg. Cypr.~1.52 (om. δ᾽), Suid.~α~1826 (om. δ᾽), Tosi n°~1554 ; Zen.~3.8 Διομήδειος ἀνάγκη; ἐπὶ τῶν κατὰ ἀνάγκην τινὰ πραττόντων. Μέμνηται αὐτῆς Ἀριστοφάνης ; \textit{idem} Suid.~δ~1164 (λέγεται καὶ ἵππος. παροιμία pro ἐπὶ τῶν … Ἀριστοφάνης)}
\lp{171–172 γνοὺς … ἐμαυτόν]}{Hsch.~γ~743 γνῶθι σαυτόν· ἀπόφθεγμα Χείλωνός φασι ; \textit{idem} Tosi n°~1992 (γνῶθι σεαυτόν ; om. ἀπόφθεγμα Χείλωνός φασι)}
\lp{172 ἐμαυτὸν ἐθεράπευσα]}{Eur.~\textit{Ba.}~932 ἀλλ᾽ αὐτὸν ἡμεῖς, οἷς σε θεραπεύειν μέλει, / πάλιν καταστελοῦμεν … ἰδού, σὺ κόσμει ; Pl.~\textit{Grg.}~465b τῇ δὲ γυμναστικῇ κατὰ τὸν αὐτὸν τρόπον τοῦτον ἡ κομμωτική, κακοῦργός τε καὶ ἀπατηλὴ καὶ ἀγεννὴς καὶ ἀνελεύθερος, σχήμασιν καὶ χρώμασιν καὶ λειότητι καὶ ἐσθῆσιν ἀπατῶσα ; Hsch.~η~293 ἠκεσάμην· ἐθεράπευσα ; Hsch.~η~893 ἤσκει· ἐθεράπευεν. ἐπονεῖτο. ἐκόσμει ; Ar.~\textit{Eq.}~1261 καὶ μὴν ἐγώ σ᾽ ὦ Δῆμε θεραπεύσω καλῶς ; Eur.~\textit{Hipp.}~186 κρεῖσσον δὲ νοσεῖν ἢ θεραπεύειν ; Eur.~\textit{Or.}~222 ἀδέλφ᾽ ἀδελφῇ χειρὶ θεραπεύειν μέλη ; Pl.~\textit{Grg.}~517e ἣ δὴ τῷ ὄντι γε ἐστὶν σώματος θεραπεία ; Hsch.~η~294 ἠκέσατο· ἰάσατο, ἐθεράπευσεν ; Hsch.~θ~330 θεραπευτής· ἰατρός ; Hsch.~θ~332 θεραπηΐας· τὰς βωμολοχίας ; Hsch.~κ~3440 κομέειν· κομεῖν. ἐπιμελεῖσθαι, θεραπεύειν (λ 250). ὅθεν καὶ κόμη καὶ ἱπποκόμος ; Ar.~\textit{Eq.}~59 ἄλλον θεραπεύειν, ἀλλὰ βυρσίνην ἔχων ; Ar.~\textit{Eq.}~799 πάντως δ᾽ αὐτὸν θρέψω γ᾽ὼ καὶ θεραπεύσω ; Men.~fr.~591~Kock λύπην γὰρ εὔνους οἶδε θεραπεύειν λόγος ; Hsch.~θ~329 θεραπευσία· θεραπεία. ἱκετεία ; Hsch.~π~3056 πορσαίνει· κοσμεῖ. θεραπεύει. κατασκευάζει ; Suid.~θ~227 Θεραπεύω· αἰτιατικῇ}
\lp{172 ἐμαυτὸν ἐθεράπευσα (ἰατρέ, θεράπευσον σεαυτόν)]}{Eur.~fr.~1086~N.² ἄλλων ἰατρὸς αὐτὸς ἕλκεσιν βρύων ; \textit{idem} Apostol.~2.26 (ἕλκεσι ; add. ἐκ τοῦ τραγῳδοῦ, καὶ σαφής) ; Aesch.~\textit{PV}~473–475 κακὸς δ᾽ ἰατρὸς ὥς τις ἐς νόσον / πεσὼν ἀθυμεῖς καὶ σεαυτὸν οὐκ ἔχεις / εὑρεῖν ὁποίοις φαρμάκοις ἰάσιμος ; Ev.Luc.~4.23 Ἰατρέ, θεράπευσον σεαυτόν ; Tosi n°~1729 Medice, cura te ipsum ! ; Erasm.~\textit{Adag.}~2.5.38 Aliorum Medicus, ipse ulceribus scates … Ἄλλων ἰατρὸς, αὐτὸς ἕλκεσι βρύων. i. e. Aliis medens, at ipsus ulceribus scatens. In eos torquendus est, qui errata castigant incessuntque aliena, intercutibus vitiis ipsi madentes. Illud vulgo tritissimum: Medice, tibi ipsi medicus esto}
\lp{172 ἐμαυτὸν ἐθεράπευσα (ἐμαυτῷ βαλανεύσω)]}{Ar.~\textit{Pax}~1103 κἀγὼ ᾽μαυτῷ βαλανεύσω ; \textit{idem} Zen.~3.58 (ἐμαυτῷ ; add. παροιμία), Greg. Cypr.~2.32 (ἐμαυτῷ ; add. ἐμαυτῷ διακονήσω), App. Prov.~2.53 (ἐμαυτῷ ; add. ἀντὶ τοῦ διακονήσω), Hsch.~ε~2272 (ἐμαυτῷ ; add. ἐμαυτῷ διακονήσω), Suid.~ε~932 (ἐμαυτῷ ; add. οἱονεὶ ἐμαυτῷ διακονήσω)}
\lp{172 πῶς πρὸς τῶν θεῶν;]}{Soph.~\textit{OT}~1009 πῶς, ὦ γεραιέ; πρὸς θεῶν δίδασκέ με ; Ar.~\textit{Nub.}~200 πρὸς τῶν θεῶν τί γὰρ τάδ᾽ ἐστίν; εἰπέ μοι ; Ar.~\textit{Nub.}~481 τί δέ; τειχομαχεῖν μοι διανοεῖ πρὸς τῶν θεῶν ; Ar.~\textit{Av.}~69 ἀτὰρ σὺ τί θηρίον ποτ᾽ εἶ πρὸς τῶν θεῶν ; Ar.~\textit{Thesm.}~228 μηδαμῶς πρὸς τῶν θεῶν ; Ar.~\textit{Eq.}~1390 ὦ Ζεῦ πολυτίμηθ᾽ ὡς καλαί· πρὸς τῶν θεῶν ; Ar.~\textit{Pax}~376 μὴ πρὸς τῶν θεῶν ; Ar.~\textit{Av.}~996 πρὸς τῶν θεῶν ; Ar.~\textit{Eccl.}~562 μηδαμῶς πρὸς τῶν θεῶν ; Ar.~\textit{Plu.}~1176 ἡ δ᾽ αἰτία τίς ἐστιν ὦ πρὸς τῶν θεῶν ; Eur.~\textit{Hipp.}~219 πρὸς θεῶν, ἔραμαι κυσὶ θωΰξαι}
\lp{173 παῦσαι βαΰζων]}{\textit{Test.} Suid.~β~192 Βαΰζων: ὑλακτῶν. Ἀριστοφάνης· βαΰζων· καὶ γὰρ ἐγὼ τοιοῦτος ἦν. ἀντὶ τοῦ ἤμην ; \textit{idem} Zonar.~379 (βαΰζων γὰρ καὶ ἐγώ ; om. ἀντὶ τοῦ ἤμην) ; Ar.~\textit{Vesp.}~903 αὗ αὗ ; Ar.~\textit{Av.}~1381 παῦσαι μελῳδῶν, ἀλλ᾽ ὅ τι λέγεις εἰπέ μοι ; Ar.~\textit{Thesm.}~895 βάυζε τοὐμὸν σῶμα βάλλουσα ψόγῳ ; Cratin.~fr.~6~Kock εἶδες τὴν Θασίαν ἅλμην οἷ᾽ ἄττα βαύζει ; Eur.~\textit{Alc.}~707 παῦσαι δέ, πρέσβυ, παῖδα σὸν κακορροθῶν ; Aesch.~\textit{Ag.}~449 τάδε σῖγά τις βαΰ- / ζει, φθονερὸν δ᾽ ὑπ᾽ ἄλγος ἕρ- / πει προδίκοις Ἀτρείδαις ; Aesch.~\textit{Pers.}~13 πᾶσα γὰρ ἰσχὺς Ἀσιατογενὴς / ᾤχωκε, νέον δ᾽ ἄνδρα βαΰζει ; Lyc.~1453 βαύζω, κενὸν ψάλλουσα μάστακος κρότον ; Heraclit.~fr.~97~DK κύνες γὰρ καταβαΰζουσιν ὧν ἂν μὴ γινώσκωσι ; \textit{idem} Plu.~\textit{Mor.}~787c–788a (καὶ βαΰζουσιν ὅν ; add. καθ᾽ Ἡράκλειτον) ; Hsch.~β~357 βαΰζειν· ὑλακτεῖν ἀσαφῶς λέγειν ; Hsch.~υ~145 ὑλακτηθέντων· φλυαρηθέντων ; Poll.~5.86 εἰρήκασι δέ τινες τῶν ποιητῶν καὶ βαΰζειν τοὺς κύνας ; Ar.~\textit{Ach.}~1107 ὦνθρωπε παῦσαι καταγελῶν μου τῶν ὅπλων ; Ar.~\textit{Vesp.}~904 ἀγαθός γ᾽ ὑλακτεῖν καὶ διαλείχειν τὰς χύτρας ; Ar.~\textit{Vesp.}~1402 θρασεῖα καὶ μεθύση τις ὑλάκτει κύων ; Ar.~\textit{Eccl.}~129 Ἀρίφραδες παῦσαι λαλῶν ; Ar.~\textit{Plu.}~360 παῦσαι φλυαρῶν ὦγάθ᾽ ; Ar.~\textit{Plu.}~369 σὺ μὲν οἶδ᾽ ὃ κρώζεις ; Pl. Com. fr.~224~Kock παῦσαι δυσωνῶν ; Philetaer.~fr.~6~Kock παῦσαι γέρων ὢν τοὺς τρόπους ; Eur.~\textit{Ba.}~809 σὺ δὲ παῦσαι λέγων ; Eur.~\textit{Hipp.}~706 παῦσαι λέγουσα ; Aesch.~\textit{Ag.}~1307–1308 φεῦ φεῦ. / τί τοῦτ᾽ ἔφευξας ; Aesch.~\textit{Pers.}~574 τεῖνε δὲ δυσβάυκτον / βοᾶτιν τάλαιναν αὐδάν ; Soph.~\textit{Ichn.}~399 [τὸ]ν παῖδα παῦσαι τοῦ Διὸς [κακῶς λέγων ; Theoc.~6.10 ἅ τοι τᾶν ὀίων ἕπεται σκοπός· ἁ δὲ βαΰσδει ; Luc.~\textit{Merc.Cond.}~34 κυνίδιον ἐκ τοῦ ἱματίου προκῦπτον … καὶ βαΰζον λεπτῇ τῇ φωνῇ ; Ath.~3.99e–100b ἀλλὰ μὴ βάυζε, εἶπεν, ὦ ἑταῖρε, μηδὲ ἀγριαίνου τὴν κυνικὴν προβαλλόμενος λύσσαν ; Σ Ar.~\textit{Thesm.}~231 ὅτι τὸ μύζειν παρῆκται ἀπὸ τοῦ μῦ πολλοῖς ἄλλοις ὁμοίως. μύζειν δέ ἐστι τὸ τοῖς μυκτῆρσι ποιὸν ἦχον ἀποτελεῖν ; Σ Soph.~\textit{Aj.}~334 ἀπηνεστέρῳ χρήσασθαι τῷ φθέγματι καὶ κυνικώτερον βαύζειν ; Hsch.~ρ~33 ῥάζειν· τρώγειν. κυρίως ἐπὶ τῶν κυνῶν ; Hsch.~ρ~475 ῥύζειν· ὑλακτεῖν ; Hsch.~ρ~476 ῥύζουσι· διαμωκῶνται. μισοῦσι. γογγύζουσιν ; Hsch.~υ~146 ὑλακτῶν· μετὰ μανίας κράζων ; Ioann. Alex.~\textit{Praec.\,Ton.}~32.23 Dind. Καὶ τὸ βαὺ κατὰ μίμησιν κυνὸς ὀξύνεται, βαὺ βαὺ καὶ κυνὸς φωνὴν ἵεις· ἐξ οὗ καὶ τὸ βαύζω ῥῆμα ; \textit{idem} Hdn.~\textit{Cath.\,Pros.} (19) (ἱείς) ; Ar.~\textit{Pax}~1229 παῦσαί μ᾽ ὑβρίζων τοῖς ἐμοῖσι χρήμασιν ; Ar.~\textit{Av.}~859 παῦσαι σὺ φυσῶν ; Ar.~\textit{Thesm.}~231 τί μύζεις; πάντα πεποίηται καλῶς ; Philem.~fr.~213~Kock ὦ Κλέων, παῦσαι φλυαρῶν ; Eur.~\textit{Tro.}~145 τύφεται Ἴλιον, αἰάζωμεν ; Σ Il.~22.69 ἵνα βαΰζοντες δηλοῖεν τοῖς ἔνδον τοὺς ξένους ; Hsch.~β~1 βαβάζειν· τὸ ⟨μὴ⟩ διηρθρωμένα λέγειν. ἔνιοι δὲ βοᾶν ; Hsch.~υ~142 ὑλακή· ὑλαγμός, βοή ; Hdn.~\textit{Περὶ\,μον.\,λέξ.}~2.61 ϲαλπίζω, βαύζω, χαιτίζω}
\lp{173 καὶ γὰρ ἐγὼ τοιοῦτος ἦν]}{Eur.~\textit{Hec.}~284 κἀγὼ γὰρ ἦ ποτ᾽, ἀλλὰ νῦν οὐκ εἴμ᾽ ἔτι ; Soph.~\textit{Phil.}~96–97 ἐσθλοῦ πατρὸς παῖ, καὐτὸς ὢν νέος ποτὲ / γλῶσσαν μὲν ἀργόν, χεῖρα δ᾽ εἶχον ἐργάτιν ; Arist.~\textit{Po.}~18.1456a25–30 καὶ τὸν χορὸν δὲ ἕνα δεῖ ὑπολαμβάνειν τῶν ὑποκριτῶν … καὶ συναγωνίζεσθαι μὴ ὥσπερ Εὐριπίδῃ ἀλλ᾽ ὥσπερ Σοφοκλεῖ … διὸ ἐμβόλιμα ᾄδουσιν πρώτου ἄρξαντος Ἀγάθωνος τοῦ τοιούτου ; Σ Ar.~\textit{Thesm.}~173 τοιοῦτος ἦ· κακὸς ἤμην ; Σ Eur.~\textit{Hec.}~13 νεώτατος δ᾽ ἦν: ἀντὶ τοῦ ἤμην φησίν. Ἀττικῶς δὲ ἦν. καὶ χωρὶς δὲ τοῦ ν ἦ, ἀντὶ τοῦ ἔα. οὕτω Δίδυμος ; Moer. p.~197,~32~Bekker ἦν Ἀττικοί, ἤμην Ἕλληνες ; Phot. s.v.~ἦν ἦν· ἤμην. … ἦν· ἀντὶ τοῦ ἤμην· τὸ δὲ παλαιὸν ἔλεγον ἦ κατὰ συναλιφήν, οἷον ἦα καὶ ἔα ; Suid.~η~5 Ἦ: ἀντὶ τοῦ ἤμην. καὶ χωρὶς τοῦ ν. Ἀριστοφάνης Πλούτῳ· ἃ κρύπτειν ἦ παρεσκευασμένος ; Ar.~\textit{Thesm.}~63 ἦ που νέος γ᾽ ὢν ἦσθ᾽ ὑβριστὴς ὦ γέρον ; Ar.~\textit{Eccl.}~998 οἶδ᾽ οἶδ᾽ ὅ τι βούλει. — καὶ γὰρ ἐγώ σε νὴ Δία ; Ar.~\textit{Eq.}~1092 νὴ Δία καὶ γὰρ ἐγώ ; Ar.~\textit{Thesm.}~478 ὅτε νύμφη μὲν ἦν τρεῖς ἡμέρας ; Ar.~\textit{Plu.}~29 κακῶς ἔπραττον καὶ πένης ἦν ; Ar.~\textit{Plu.}~695 ἔπειτ᾽ ἐπειδὴ μεστὸς ἦν, ἀνεπαλλόμην ; Ar.~\textit{Plu.}~821–822 οὐχ οἷός τε γὰρ / ἔνδον μένειν ἦν ; Eur.~\textit{HF}~1416 ὡς ἐς τὸ λῆμα παντὸς ἦν ἥσσων ἀνήρ ; Eur.~\textit{Hipp.}~1012 μάταιος ἆρ᾽ ἦν}
\lp{173 καὶ γὰρ ἐγὼ τοιοῦτος ἦν (πάλαι ποτ᾽ ἦσαν ἄλκιμοι Μιλήσιοι)]}{Ar.~\textit{Plu.}~1002 πάλαι ποτ᾽ ἦσαν ἄλκιμοι Μιλήσιοι ; \textit{idem} Zen.~5.80, Diogenian.~5.3 (ἦσάν ποτ᾽ ἦσαν ; add. ἐπὶ τῶν ποτὲ εὖ, εἶθ᾽ ἑτέρως γεγονότων. ὁμοία, ἄμμες ποτ᾽ ἦμες), Greg. Cypr.~2.62 (ἦσαν ποτ᾽ ἦσαν), Suid.~π~61 (add. ἐπὶ τῶν ποτε ἀκμασάντων, νῦν δὲ ἀσθενῶν ὄντων) ; Ar.~\textit{Vesp.}~1060 ὦ πάλαι ποτ᾽ ὄντες ἡμεῖς ἄλκιμοι μὲν ἐν χοροῖς ; Plu.~\textit{Lyc.}~21.2 ἄμμες πόκ᾽ ἦμες ἄλκιμοι νεανίαι ; Erasm.~\textit{Adag.}~1.9.49 Fuere quondam strenui Milesii. Πάλαι ποτ᾽ ἦσαν ἄλκιμοι Μιλήσιοι … Hoc senario notabant eos, qui à pristina fortuna defluxissent, aut à majorum moribus degenerassent. Breviter, in omnes quadrabit, qui desierunt esse, quod erant, facti deteriores, in senem … Allusit eodem in Vespis, jocans in senes ; Erasm.~\textit{Adag.}~1.9.50 Fuimus Troes. Ἦμεν Τρῶες. i. e. Fuimus Troes. Hic loquendi color in Graecanicis tragoediis adeo passim obvius est, ut tanquam in adagium abierit, ubi perisse quempiam significant. Euripides in Hecuba: Κἀγὼ γὰρ ἦν ποτέ. i. e. Et quondam ego fui}
\lp{174 ὢν τηλικοῦτος]}{Ar.~\textit{Eq.}~881 τονδὶ δ᾽ ὁρῶν ἄνευ χιτῶνος ὄντα τηλικοῦτον ; Ar.~\textit{Eccl.}~1039 παρὰ σοὶ καθεύδειν τηλικοῦτος ὤν ; Antiph.~fr.~261~Kock ὅστις δ᾽ ἐρυθριᾷ τηλικοῦτος ὢν ἔτι / πρὸς τοὺς ἑαυτοῦ γονέας, οὐκ ἔστιν κακός ; Soph.~\textit{El.}~614 καὶ ταῦτα τηλικοῦτος ; Hsch.~τ~775 τηλίκος· τηλικοῦτος ἢ νήπιος ; Phot.~τ~247 τηλικοῦτος· ἐπὶ ἡλικίας τίθεται. οὕτως Μένανδρος (Asp. 258, Dysc. 696) ; Philetaer.~fr.~6~Kock παῦσαι γέρων ὢν τοὺς τρόπους ; Aesch.~\textit{Ag.}~1619–1620 γνώσῃ γέρων ὢν ὡς διδάσκεσθαι βαρὺ / τῷ τηλικούτῳ ; Soph.~\textit{Ant.}~767 νοῦς δ᾽ ἐστὶ τηλικοῦτος ἀλγήσας βαρύς ; Pl.~\textit{Ap.}~25d τοσοῦτον σὺ ἐμοῦ σοφώτερος εἶ τηλικούτου ὄντος τηλικόσδε ὤν ; Pl.~\textit{Prt.}~318b κἂν σύ, καίπερ τηλικοῦτος ὢν καὶ οὕτως σοφός ; Luc.~\textit{DMort.}~9.2 ἐραστὰς σὺ τηλικοῦτος ὤν, ὀδόντας τέτταρας ἔχων ; Hsch.~τ~1644 τυννός· μικρός. τηλικοῦτος (Callim. fr. 471, 1 Pf. Theocr. 24, 139) ; Soph.~\textit{OC}~751 πτωχῷ διαίτῃ, τηλικοῦτος, οὐ γάμων ; Is.~3.55 ὄντων αὐτῷ παίδων ἤδη τηλικούτων ἐκ τῆς γυναικός ; Lys.~21.20 οἳ περὶ ἀσεβείας μὲν ἀγωνιζόμενοι τηλικοῦτοι γεγόνασιν}
\lp{174 ἡνίκ᾽ ἠρχόμην ποιεῖν]}{Eur.~\textit{Heracl.}~889 ὃ πρῶτον εἶπας, ἡνίκ᾽ ἠρχόμην λόγου ; Vit. Eur.~1 ἤρξατο δὲ διδάσκειν ἐπὶ Καλλίου ἄρχοντος κατὰ ὀλυμπιάδα πά΄ ἔτει α΄, πρῶτον δὲ ἐδίδαξε τὰς Πελιάδας, ὅτε καὶ τρίτος ἐγένετο ; Phot. s.v.~ἡνίκα ἡνίκα· ἀνταποδίδοται τῷ τηνικαῦτα καὶ τηνικάδε· σημαίνει δὲ τὸ μὲν ἡνίκα καθ᾽ ἣν ὥραν ; LXX \textit{Ge.}~48.7 ἐγὼ δὲ ἡνίκα ἠρχόμην ἐκ Μεσοποταμίας ; Hsch.~η~599 ἡνίκα· ὅτε, ἢ ὅταν}
\lp{174 ὢν τηλικοῦτος, ἡνίκ᾽ ἠρχόμην ποιεῖν (ἐξ εὐρυπρώκτων)]}{Ar.~\textit{Eq.}~423–426 ὥστ᾽ εἶπ᾽ ἀνὴρ τῶν ῥητόρων ἰδών με τοῦτο δρῶντα· / οὐκ ἔσθ᾽ ὅπως ὁ παῖς ὅδ᾽ οὐ τὸν δῆμον ἐπιτροπεύσει ; Ar.~\textit{Eq.}~878–880 οὔκουν σε δῆτα ταῦτα δεινόν ἐστι πρωκτοτηρεῖν / παῦσαί τε τοὺς βινουμένους; κοὐκ ἔσθ᾽ ὅπως ἐκείνους / οὐχὶ φθονῶν ἔπαυσας, ἵνα μὴ ῥήτορες γένωνται ; Ar.~\textit{Nub.}~1091–1092 τί δαί; τραγῳδοῦσ᾽ ἐκ τίνων; / ἐξ εὐρυπρώκτων. — εὖ λέγεις ; Ar.~\textit{Eccl.}~112–113 λέγουσι γὰρ καὶ τῶν νεανίσκων ὅσοι / πλεῖστα σποδοῦνται, δεινοτάτους εἶναι λέγειν ; Pl.~\textit{Smp.}~192a μέγα δὲ τεκμήριον· καὶ γὰρ τελεωθέντες μόνοι ἀποβαίνουσιν εἰς τὰ πολιτικὰ ἄνδρες οἱ τοιοῦτοι. ἐπειδὰν δὲ ἀνδρωθῶσι, παιδεραστοῦσι ; Ar.~\textit{Eq.}~1242 ἠλλαντοπώλουν καί τι καὶ βινεσκόμην ; Ar.~\textit{Nub.}~1089–1090 συνηγοροῦσιν ἐκ τίνων; / ἐξ εὐρυπρώκτων ; Ar.~\textit{Nub.}~1093–1094 δημηγοροῦσι δ᾽ ἐκ τίνων; / ἐξ εὐρυπρώκτων ; Ar.~\textit{Vesp.}~687–688 ὅταν εἰσελθὸν μειράκιόν σοι κατάπυγον, Χαιρέου υἱός, / ὡδὶ διαβὰς διακινηθεὶς τῷ σώματι καὶ τρυφερανθείς}
\lp{175 μὰ τὸν Δί᾽ οὐ ζηλῶ σε τῆς παιδεύσεως]}{\textit{Test.} Suid.~ζ~63 Ζηλῶ: μακαρίζω. Ἀριστοφάνης Θεσμοφοριαζούσαις· οὐ ζηλῶ σε τῆς παιδεύσεως. ζηλῶ σε τῆς εὐβουλίας, μᾶλλον δὲ τῆς εὐωχίας ; Ar.~\textit{Ach.}~1008–1009 ζηλῶ σε τῆς εὐβουλίας, / μᾶλλον δὲ τῆς εὐωχίας ; Ar.~\textit{Eq.}~836–837 ὦ πᾶσιν ἀνθρώποις φανεὶς μέγιστον ὠφέλημα, / ζηλῶ σε τῆς εὐγλωττίας ; Ar.~\textit{Vesp.}~1450–1451 ζηλῶ γε τῆς εὐτυχίας / τὸν πρέσβυν οἷ μετέστη ; Ar.~\textit{Thesm.}~1118 ἐμὲ δὲ καὐτὸν τῆς κόρης / ταύτης ἔρως εἴληφεν. — οὐ ζηλῶσί σε ; Eur.~\textit{IA}~677 ζηλῶ σὲ μᾶλλον ἢ ᾽μὲ τοῦ μηδὲν φρονεῖν ; Eur.~\textit{Med.}~60 ζηλῶ σ᾽· ἐν ἀρχῇ πῆμα κοὐδέπω μεσοῖ ; Eur.~\textit{Or.}~521–522 οὐδὲ σὲ ζηλῶ, κακῆς / γυναικὸς ἐλθόνθ᾽ οὕνεκ᾽ ἐς Τροίας πέδον ; Aesch.~\textit{PV}~332 ζηλῶ σ᾽ ὁθούνεκ᾽ ἐκτὸς αἰτίας κυρεῖς ; Soph.~\textit{El.}~1027 ζηλῶ σε τοῦ νοῦ, τῆς δὲ δειλίας στυγῶ ; Hsch.~ζ~141 ζηλῶ· μακαρίζω. Σοφοκλῆς Φιλοκτήτῃ ἐν Τροίᾳ ; \textit{idem} Phot. s.v.~ζηλῶ (om. Σοφοκλῆς Φιλοκτήτῃ ἐν Τροίᾳ) ; Hsch.~ζ~142 ζηλῶ σε· μακαρίζω σε. μιμοῦμαί σε ; Moer. p.~197,~23–25~Bekker ζηλῶ Ἀττικοί, ζηλοτυπῶ Ἕλληνες. ζήλωσις Ἀττικοί, μίμησις Ἕλληνες. ζηλῶ Ἀττικοί, μακαρίζω Ἕλληνες ; Eur.~\textit{IA}~16 ζηλῶ σέ, γέρον ; Eur.~\textit{IA}~1406 ζηλῶ δὲ σοῦ μὲν Ἑλλάδ᾽, Ἑλλάδος δὲ σέ ; Eur.~\textit{Or.}~1673 ζηλῶ δέ σε ; Aesch.~\textit{Pers.}~712 νῦν τέ σε ζηλῶ θανόντα ; Soph.~\textit{Aj.}~552 καίτοι σε καὶ νῦν τοῦτό γε ζηλοῦν ἔχω ; Archil. ap. Plu.~\textit{Mor.}~470c οὐδ᾽ εἷλὲ πώ με ζῆλος ; X.~\textit{Symp.}~4.45 νὴ τὴν Ἥραν, ἔφη, τά τε ἄλλα ζηλῶ σε τοῦ πλούτου ; Ammon.~\textit{Diff.}~211 ζῆλος φθόνου διαφέρει. ζῆλος μὲν γάρ ἐστιν ἡ δι᾽ ἐπιθυμίαν γινομένη μίμησις ; Σ Eur.~\textit{Or.}~521 οὐδὲ σὲ ζηλῶ: ζηλῶ τὸ μακαρίζω· ἐνταῦθα δὲ ἀντὶ τοῦ ἐπαινῶ ; Hsch.~ζ~136 ζηλοῖ· φθονεῖ. ἐρίζει. μισεῖ. μιμεῖται ; Hsch.~ζ~144 ζηλωτός· μακαριστός ; \textit{idem} Synag.~ζ~7 (add. ἐπαινετός), Phot. s.v.~ζηλωτός (add. ἐπαινετός) ; Suid.~ζ~59 Ζηλοτυποῦν· … ζηλοῦν γὰρ τὸ μακαρίζειν, ζηλοτυπεῖν δὲ τὸ εἰς ἀπέχθειαν ἐληλυθέναι ; Ar.~\textit{Ran.}~1183 μὰ τὸν Δί᾽ οὐ δῆτ᾽, ἀλλὰ κακοδαίμων φύσει}
\lp{175 παίδευσις]}{Ar.~\textit{Nub.}~986 ἐξ ὧν ἄνδρας Μαραθωνομάχας ἡμὴ παίδευσις ἔθρεψεν ; Pl.~\textit{Smp.}~184d–e ὁ μὲν δυνάμενος εἰς φρόνησιν καὶ τὴν ἄλλην ἀρετὴν συμβάλλεσθαι, ὁ δὲ δεόμενος εἰς παίδευσιν καὶ τὴν ἄλλην σοφίαν κτᾶσθαι ; Ar.~\textit{Nub.}~937 παίδευσιν, ὅπως ἂν ἀκούσας σφῷν ; Ar.~\textit{Nub.}~1043 σκέψαι δὲ τὴν παίδευσιν ᾗ πέποιθεν ὡς ἐλέγξω ; Pl.~\textit{Prt.}~349a ἀπέφηνας παιδεύσεως καὶ ἀρετῆς διδάσκαλον ; Antipho Soph. fr.~60 πρῶτον, οἶμαι, τῶν ἐν ἀνθρώποις ἐστὶ παίδευσις ; Hsch.~π~59 παιδεύειν· τρέφειν, παιδοτροφεῖν ; \textit{idem} Suid.~π~851 (om. τρέφειν) ; Hsch.~π~1422 πεπαιδεῦσθαι· πεπαιδοποιῆσθαι}
\lp{176 ἀλλ᾽ ὧνπερ οὕνεκ᾽ ἦλθον]}{\textit{Test.} Σ Ar.~\textit{Thesm.}~500 ὅτι καὶ συστολαῖς βιαίοις χρῶνται καὶ ἐκτάσεσι. καὶ ἄνω « ἀλλ᾽ ὧνπερ οὕνεχ᾽ ἦλθον ἀλλ᾽ εἰπεῖν βούλομαι » ; Ar.~\textit{Nub.}~238 ἵνα με διδάξῃς ὧνπερ οὕνεκ᾽ ἐλήλυθα ; Ar.~\textit{Ran.}~108–109 ἀλλ᾽ ὧνπερ ἕνεκα τήνδε τὴν σκευὴν ἔχων / ἦλθον κατὰ σὴν μίμησιν ; Ar.~\textit{Plu.}~1200 ὧν δ᾽ οὕνεκ᾽ ἦλθον; — πάντα σοι πεπράξεται ; Eur.~\textit{Alc.}~1019–1020 ὧν δ᾽ οὕνεχ᾽ ἥκω δεῦρ᾽ ὑποστρέψας πάλιν / λέξω. γυναῖκα τήνδε μοι σῷσον λαβών ; Eur.~\textit{Andr.}~1238 ὧν δ᾽ οὕνεκ᾽ ἦλθον σημανῶ, σὺ δ᾽ ἐνδέχου ; Eur.~\textit{Hec.}~787 ἀλλ᾽ ὧνπερ οὕνεκ᾽ ἀμφὶ σὸν πίπτω γόνυ ; Eur.~\textit{Hel.}~144–145 ὧν δ᾽ οὕνεκ᾽ ἦλθον τούσδε βασιλείους δόμους, / τὴν θεσπιῳδὸν Θεονόην χρῄζων ἰδεῖν ; Eur.~\textit{Med.}~932–933 ἀλλ᾽ ὧνπερ οὕνεκ᾽ εἰς ἐμοὺς ἥκεις λόγους, / τὰ μὲν λέλεκται, τῶν δ᾽ ἐγὼ μνησθήσομαι ; Ar.~\textit{Ran.}~1414 οὐδὲν ἄρα πράξεις ὧνπερ ἦλθες οὕνεκα ; Plaut.~\textit{Rud.}~31 nunc, huc qua causa veni, argumentum eloquar ; Soph.~\textit{OT}~447 εἰπὼν ἄπειμ᾽ ὧν οὕνεκ᾽ ἦλθον ; Phot.~ε~245 εἵνεκα καὶ ἕνεκα καὶ εἵνεκεν· οἱ δὲ παλαιοὶ καὶ οὕνεκα ; Eur.~\textit{IT}~1388 ἔχομεν γὰρ ὧνπερ οὕνεκ᾽ ἄξενον πόρον ; Eur.~\textit{Or.}~611 ὧν εἵνεκ᾽ ἦλθον θυγατρὶ κοσμήσων τάφον ; Hsch.~ο~1787 οὕνεκα· οὗ χάριν ; Ammon.~\textit{Diff.}~362 Nickau οὕνεκα καὶ εἵνεκα διαφέρει. οὕνεκα μὲν σημαίνει τὸ ὅτι, εἵνεκα δὲ χάριν}
\lp{176 ἔα μ᾽ εἰπεῖν]}{Ar.~\textit{Nub.}~932 δεῦρ᾽ ἴθι, τοῦτον δ᾽ ἔα μαίνεσθαι ; Ar.~\textit{Thesm.}~64 ὦ δαιμόνιε τοῦτον μὲν ἔα χαίρειν, σὺ δὲ ; Ar.~\textit{Eccl.}~564 ὦ δαιμόνι᾽ ἀνδρῶν τὴν γυναῖκ᾽ ἔα λέγειν ; Ar.~\textit{Eccl.}~784 ὦ δαιμόνι᾽ ἀνδρῶν ἔα με τῶν προὔργου τι δρᾶν ; Eur.~\textit{Hipp.}~323 ἔα μ᾽ ἁμαρτεῖν ; Soph.~\textit{Ant.}~95 ἀλλ᾽ ἔα με καὶ τὴν ἐξ ἐμοῦ δυσβουλίαν ; Hsch.~ε~4 ἔα· ἄφες ; Ar.~\textit{Thesm.}~1077 ὦγάθ᾽ ἔασόν με μονῳδῆσαι ; Com. adesp. fr.~211~Kock ἔα μ᾽ ἀπολέσθαι ; Aesch.~\textit{PV}~386 ἔα με τῇδε τῇ νόσῳ νοσεῖν ; Soph.~\textit{OC}~625 ἔα μ᾽ ἐν οἷσιν ἠρξάμην ; Soph.~\textit{OT}~1451 ἀλλ᾽ ἔα με ναίειν ὄρεσιν ; Soph.~\textit{Phil.}~1397 ἔα με πάσχειν ταῦθ᾽ ἅπερ παθεῖν με δεῖ}
\lp{176 λέγε]}{Ar.~\textit{Vesp.}~944 ἀνάβαιν᾽, ἀπολογοῦ. τί σεσιώπηκας; λέγε ; Ar.~\textit{Ran.}~1442 ἐγὼ μὲν οἶδα καὶ θέλω φράζειν. — λέγε ; Ar.~\textit{Lys.}~214 ὅστις πρὸς ἐμὲ πρόσεισιν ἐστυκώς. λέγε}
\lp{177–178 σοφοῦ πρὸς ἀνδρός … συντέμνειν λόγους]}{\textit{Test.} Suid.~οι~154 καὶ αὖθις· σοφοῦ πρὸς ἀνδρός, ὅστις ἐν βραχεῖ πολλοὺς καλῶς οἷός τε συντέμνειν λόγους ; \textit{idem} Suid.~σ~792 (Ἀριστοφάνης ; ὅς τις ; συντεμεῖν ; add. ἐγὼ δὲ κοινῇ ξυμφορᾷ πεπληγμένος ἱκέτης ἀφῖγμαι πρὸς σέ) ; Men.~fr.~3~Mein. (Stob.~3.35.5) Τοῖς ἀναιδέσιν βοηθεῖ γὰρ λόγοις τοῦθ᾽ ἓν μόνον, / ἂν βραχεῖς αὐτοὺς ποῇ τις τόν τε καιρὸν εὖ λάβῃ ; Nicostr. Com. fr.~27~Kock εἰ τὸ συνεχῶς καὶ πολλὰ καὶ ταχέως λαλεῖν / ἦν τοῦ φρονεῖν παράσημον, αἱ χελιδόνες / ἐλέγοντ᾽ ἂν ἡμῶν σωφρονέστεραι πολύ ; Nicostr. Com. fr.~28~Kock οὐκ ἔστιν ὅστις πάντ᾽ ἀνὴρ εὐδαιμονεῖ. / νὴ τὴν Ἀθηνᾶν συντόμως γε, φίλτατε / Εὐριπίδη, τὸν βίον ἔθηκας εἰς στίχον ; Eur.~fr.~28~N.² (Stob.~3.35.3) Παῖδες, σοφοῦ πρὸς ἀνδρός, ὅστις ἐν βραχεῖ / πολλοὺς λόγους οἷός τε συντέμνειν καλῶς ; Soph.~fr.~98~N.² (Stob.~3.35.4) Βραχεῖ λόγῳ δὲ πολλὰ πρόσκειται σοφά ; Stob.~3.35 (titre) Περὶ βραχυλογίας ; Pythag. ap. Stob.~3.35.8 Μὴ ἐν πολλοῖς ὀλίγα λέγε, ἀλλ᾽ ἐν ὀλίγοις πολλά ; Apostol.~8.41c Ἡ βραχυλογία ἐγγύς ἐστι τοῦ σιγᾶν:. Λυκούργου ἀπόφθεγμα ; Eup.~\textit{Dem.} fr.~18~Mein. (Stob.~3.35.2) Καὶ τοῦ μὲν ⟨ἐν⟩ κύκλῳ γε παύσομαι λόγου, / φράσω δέ σοι τὸ πρᾶγμα διὰ τῶν χωρίων ; Philonid. ap. Stob.~3.35.7 Οὐκ ἄν τις εἴπῃ πολλά, θαυμασθήσεται, / ὁ μικρὰ δ᾽ εἰπὼν μᾶλλον, ἂν ᾖ χρήσιμα ; P.-L.~\textit{Courier,\,Pamphlet\,des\,pamphlets} (1824) et la chose est bien dite dès que chacun l'entend ; d'autant mieux dite qu'elle l'est plus brièvement ; mérite non commun, savez-vous ? ni facile, de clore en peu de mots beaucoup de sens ; Lycurg. ap. Stob.~3.35.9 Λυκοῦργος πρὸς τὸν εἰπόντα «διὰ τί Λακεδαιμόνιοι τὴν βραχυλογίαν ἀσκοῦσιν;» εἶπεν «ὅτι ἐγγύς ἐστι τοῦ σιγᾶν»}
\lp{177–178 σοφοῦ … λόγους (Εὐριπίδου Αἰόλος)]}{Ar.~\textit{Ran.}~1475 τί δ᾽ αἰσχρόν, ἢν μὴ τοῖς θεωμένοις δοκῇ ; Ar.~\textit{Nub.}~1371–1372 ὁ δ᾽ εὐθὺς ᾖσ᾽ Εὐριπίδου ῥῆσίν τιν᾽, ὡς ἐκίνει / ἀδελφὸς ὦλεξίκακε τὴν ὁμομητρίαν ἀδελφήν ; Ar.~\textit{Pax}~114–119 ὦ πάτερ ὦ πάτερ ἆρ᾽ ἔτυμός γε / δώμασιν ἡμετέροις φάτις ἥκει … / δοξάσαι ἔστι κόραι, τὸ δ᾽ ἐτήτυμον ἄχθομαι ὑμῖν ; Ar.~\textit{Ran.}~849–850 ὦ Κρητικὰς μὲν συλλέγων μονῳδίας, / γάμους δ᾽ ἀνοσίους ἐσφέρων ἐς τὴν τέχνην ; Ar.~\textit{Ran.}~1081 καὶ τικτούσας ἐν τοῖς ἱεροῖς, / καὶ μιγνυμένας τοῖσιν ἀδελφοῖς}
\lp{177–178 σοφοῦ … λόγους (βραχυλογία)]}{Ar.~\textit{Ach.}~416 δεῖ γάρ με λέξαι τῷ χορῷ ῥῆσιν μακράν ; \textit{Il.}~3.213–215 ἤτοι μὲν Μενέλαος ἐπιτροχάδην ἀγόρευε, / παῦρα μὲν ἀλλὰ μάλα λιγέως, ἐπεὶ οὐ πολύμυθος / οὐδ᾽ ἀφαμαρτοεπής ; \textit{idem} Suid.~π~816 (s.v. παῦρα (v. 214 tantum) ; om. οὐδ᾽ ἀφαμαρτοεπής) ; Pi.~\textit{I.}~6.58–59 τὸν Ἀργείων τρόπον / εἰρήσεταί που κἀν βραχίστοις ; Pi.~\textit{P.}~9.76–78 ἀρεταὶ δ᾽ αἰεὶ μεγάλαι πολύμυθοι· / βαιὰ δ᾽ ἐν μακροῖσι ποικίλλειν, / ἀκοὰ σοφοῖς ; Thal. ap. D.L.~1.35 οὔ τι τὰ πολλὰ ἔπη φρονίμην ἀπεφήνατο δόξαν· / ἕν τι μάτευε σοφόν ; Pl.~\textit{Grg.}~449c καὶ γὰρ αὖ καὶ τοῦτο ἕν ἐστιν ὧν φημι, μηδένα ἂν ἐν βραχυτέροις ἐμοῦ τὰ αὐτὰ εἰπεῖν. … ἀλλὰ ποιήσω, καὶ οὐδενὸς φήσεις βραχυλογωτέρου ἀκοῦσαι ; Pl.~\textit{Prt.}~329b Πρωταγόρας δὲ ὅδε ἱκανὸς μὲν μακροὺς λόγους καὶ καλοὺς εἰπεῖν, ὡς αὐτὰ δηλοῖ, ἱκανὸς δὲ καὶ ἐρωτηθεὶς ἀποκρίνασθαι κατὰ βραχὺ ; Pl.~\textit{Prt.}~343b ὅτι οὗτος ὁ τρόπος ἦν τῶν παλαιῶν τῆς φιλοσοφίας, βραχυλογία τις Λακωνική ; Plu.~\textit{De\,Pyth.\,or.}~29 (\textit{Mor.}~408d–409c) τὰ προγεγραμμένα τῶν σοφῶν τό γνῶθι σαυτόν καὶ τὸ μηδὲν ἄγαν ἀποδέχεσθε καὶ θαυμάζειν οὐχ ἣκιστα διὰ τὴν βραχυλογίαν ὡς πυκνὸν καὶ σφυρήλατον νοῦν ἐν ὀλίγῳ περιέχουσαν ὄγκῳ ; Plu.~\textit{De\,garr.}~17 (\textit{Mor.}~510e–511d) σοφώτεροι δοκοῦσι τῶν ἐξηνίων τούτων καὶ φερομένων οἱ στρογγύλοι καὶ βραχυλόγοι, καὶ ὧν πολὺς νοῦς ἐν ὀλίγῃ λέξει συνέσταλται ; Plu.~\textit{Lyc.}~19.1 τὸ δὲ τοῦ λόγου νόμισμα τοὐναντίον ἀπ᾽ εὐτελοῦς καὶ ὀλίγης λέξεως εἰς πολλὴν καὶ περιττὴν κατεσκεύασε διάνοιαν ; Demetr.~\textit{Eloc.}~7 δεινότερον γὰρ τὸ ἐν ὀλίγῳ πολὺ ἐμφαινόμενον καὶ σφοδρότερον, διὸ καὶ οἱ Λάκωνες βραχυλόγοι ὑπὸ δεινότητος ; Σ Ar.~\textit{Ach.}~416 Τὸ μακρὸν οὐκ ἀργῶς αὐτῷ λέλεκται, ἀλλ᾽ ἐπειδὴ καὶ τοὺς ἀγγέλους καὶ τοὺς προλόγους μακρολογοῦντας εἰσάγει Εὐριπίδης, πρὸς αὐτὸν ὁ λόγος ἀποτείνεται ; Σ Pi.~\textit{I.}~6.87a Μακρολόγοι μὲν οὖν οἱ Ἴωνες, σύντομοι δὲ οὐ μόνον Λάκωνες, ἀλλὰ καὶ Ἀργεῖοι ; Σ Pi.~\textit{P.}~9.134a ἐάν τις, φησὶ, μακρὰ ἐκλάβῃ πράγματα καὶ βαιὰ ποιήσῃ αὐτὰ τῷ λόγῳ, ἀκοὰ σοφοῖς ἐστι, τουτέστι σοφοῦ ἀκροατοῦ δεῖται ; Poll.~6.149 εἰ δ᾽ ὀλίγα μὲν λέγοι, κρίσις δέ τις αὐτῷ προσείη, βραχὺς βραχυλόγος, σύντομος, αὐτάρκης, ἀρτιεπής, σαφής … βραχέσι ῥήμασι πολὺν νοῦν ἀπαγγέλλων … πολυνοίᾳ μᾶλλον ἢ πολυλογίᾳ χαίρων ; Suid.~λ~824 βραχυλογίαν τε ἐτίμησε καὶ πενίαν, τὴν μὲν σοφίας εἶναι σημεῖον, τὴν δὲ ἀρετῆς διδάσκαλον οἰηθείς ; Suid.~χ~311 εἷς τῶν ἑπτὰ σοφῶν. ὃς ἦν βραχυλόγος … καὶ Χιλώνειος τρόπος, ἡ βραχυλογία ; Tosi n°~2097 Laconica brevitas ; Tosi n°~2098 Multa paucis ; Erasm.~\textit{Adag.}~2.10.49 Laconismus. Λακωνισμὸν veluti proverbio breviloquentiam vocant … Horum autem praecipua gratia, ut paucissimis verbis plurimum sententiae complectantur ; Pi.~\textit{N.}~4.33–34 τὰ μακρὰ δ᾽ ἐξενέπειν ἐρύκει με τεθμὸς ; Pl.~\textit{Lg.}~887b ὡς οὐδὲν ἐν τῷ παρόντι δεῖ προτιμᾶν βραχυλογίαν μᾶλλον ἢ μῆκος — οὐδεὶς γὰρ ἡμᾶς, τὸ λεγόμενον, ἐπείγων διώκει ; \textit{idem} Apostol.~18.67e (add. Πλάτων ἐν Νόμοις) ; Pl.~\textit{Prt.}~335a καὶ αὖ βραχέα οὕτως ὥστε μηδένα σοῦ ἐν βραχυτέροις εἰπεῖν· εἰ οὖν μέλλεις ἐμοὶ διαλέξεσθαι, τῷ ἑτέρῳ χρῶ τρόπῳ πρός με, τῇ βραχυλογίᾳ ; Plu.~\textit{Apophth.\,Lac.}~232b–c τοῖς ὀλίγα λέγουσιν ὀλίγων καὶ νόμων ἐστὶ χρεία ; Plu.~\textit{Lyc.}~19.3 καὶ γὰρ ὁ Λυκοῦργος αὐτὸς βραχυλόγος τις ἔοικε γενέσθαι καὶ ἀποφθεγματικός ; D.L.~1.72 Βραχυλόγος τε ἦν· ὅθεν καὶ Ἀρισταγόρας ὁ Μιλήσιος τοῦτον τὸν τρόπον Χιλώνειον καλεῖ ; Σ Il.~3.213 βραχυλόγοι γὰρ οἱ Λάκωνες ἀεί ; Σ Eur.~\textit{Hec.}~192 ἀποφθέγγομαι δὲ τὸ βραχέα καὶ συνεσταλμένα λέγω· ὅθεν καὶ ἀποφθέγματα οἱ βραχεῖς λόγοι ; Poll.~6.150 καὶ τὰ πράγματα δὲ τάχος, συντομία, αὐτάρκεια, βραχυλογία, σαφήνεια … παῦρα μὲν λέγων ἀλλὰ μάλα λιγέως ; Ammon.~\textit{Diff.}~145 φάσκει δὲ ὁ βραχυλόγος ; LXX \textit{Si.}~32.8 κεφαλαίωσον λόγον, ἐν ὀλίγοις πολλά ; Tosi n°~1963 Intelligenti pauca ; Ath.~3.1 (72a) ὅτι Καλλίμαχος ὁ γραμματικὸς τὸ μέγα βιβλίον ἴσον ἔλεγεν εἶναι τῷ μεγάλῳ κακῷ ; \textit{idem} Tosi n°~1476 (μέγα βιβλίον tantum) ; Pl.~\textit{Phdr.}~272a βραχυλογίας τε αὖ καὶ ἐλεινολογίας καὶ δεινώσεως ; Plu.~\textit{Sept.\,sap.\,conv.}~10 (\textit{Mor.}~153e–154c) ἐδεῖτο τῆς Πιττακοῦ βραχυλογίας ; Suid.~β~525 Βραχυμυθίας: βραχυλογίας}
\lp{177 σοφοῦ πρὸς ἀνδρός]}{Ar.~\textit{Eq.}~191–192 ἡ δημαγωγία γὰρ οὐ πρὸς μουσικοῦ / ἔτ᾽ ἐστὶν ἀνδρὸς οὐδὲ χρηστοῦ τοὺς τρόπους ; Ar.~\textit{Ran.}~534–535 ταῦτα μὲν πρὸς ἀνδρός ἐστι / νοῦν ἔχοντος καὶ φρένας καὶ / πολλὰ περιπεπλευκότος ; Ar.~\textit{Ran.}~540 δεξιοῦ πρὸς ἀνδρός ἐστι / καὶ φύσει Θηραμένους ; Men.~fr.~773~Kock χρηστοῦ πρὸς ἀνδρὸς μηδὲν ἐννοεῖν κακόν ; Eur.~\textit{Ba.}~641 πρὸς σοφοῦ γὰρ ἀνδρὸς ἀσκεῖν σώφρον᾽ εὐοργησίαν ; Eur.~\textit{Hel.}~950–951 καίτοι λέγουσιν ὡς πρὸς ἀνδρὸς εὐγενοῦς / ἐν ξυμφοραῖσι δάκρυ᾽ ἀπ᾽ ὀφθαλμῶν βαλεῖν ; Aesch.~\textit{Ag.}~592 ἦ κάρτα πρὸς γυναικὸς αἴρεσθαι κέαρ ; Aesch.~\textit{Ag.}~1636 τὸ γὰρ δολῶσαι πρὸς γυναικὸς ἦν σαφῶς ; Soph.~\textit{Aj.}~581–582 οὐ πρὸς ἰατροῦ σοφοῦ / θρηνεῖν ἐπῳδὰς πρὸς τομῶντι πήματι ; Men.~\textit{Mon.}~183 J. (=~121~Mein.) δὶς ἐξαμαρτεῖν ταὐτὸν οὐκ ἀνδρὸς σοφοῦ ; Ar.~\textit{Vesp.}~369 ταῦτα μὲν πρὸς ἀνδρός ἐστ᾽ ἄνοντος ἐς σωτηρίαν ; Ar.~\textit{Plu.}~354–355 οὕτως ὑπερπλουτεῖν τό τ᾽ αὖ δεδοικέναι / πρὸς ἀνδρὸς οὐδὲν ὑγιές ἐστ᾽ εἰργασμένου ; [Eur.]~\textit{Rh.}~206 σοφοῦ παρ᾽ ἀνδρὸς χρὴ σοφόν τι μανθάνειν ; Soph.~\textit{Aj.}~319–320 πρὸς γὰρ κακοῦ τε καὶ βαρυψύχου γόους / τοιούσδ᾽ ἀεί ποτ᾽ ἀνδρὸς ἐξηγεῖτ᾽ ἔχειν ; Σ Soph.~\textit{Aj.}~581 οὐκ ἔστιν ἰατροῦ σοφοῦ ἐπῳδαῖς χρῆσθαι}
\lp{177 σοφοῦ]}{Ar.~\textit{Nub.}~1377–1378 ὅστις οὐκ Εὐριπίδην ἐπαινεῖς / σοφώτατον ; Ar.~\textit{Lys.}~368 οὐκ ἔστ᾽ ἀνὴρ Εὐριπίδου σοφώτερος ποιητής ; Eur.~\textit{Med.}~580–581 ἐμοὶ γὰρ ὅστις ἄδικος ὢν σοφὸς λέγειν / πέφυκε, πλείστην ζημίαν ὀφλισκάνει ; Pl.~\textit{R.}~568a ἥ τε τραγῳδία ὅλως σοφὸν δοκεῖ εἶναι καὶ ὁ Εὐριπίδης διαφέρων ἐν αὐτῇ ; Pl.~\textit{Smp.}~174c φαῦλος ὢν ἐπὶ σοφοῦ ἀνδρὸς ἰέναι θοίνην ἄκλητος ; Pl.~\textit{Smp.}~175e ἡ δὲ σὴ λαμπρά τε καὶ πολλὴν ἐπίδοσιν ἔχουσα … ὑβριστὴς εἶ, ἔφη, ὦ Σώκρατες, ὁ Ἀγάθων … διαδικασόμεθα ἐγώ τε καὶ σὺ περὶ τῆς σοφίας ; Σ Ar.~\textit{Nub.}~144 σοφὸς Σοφοκλῆς, σοφώτερος δ᾽ Εὐριπίδης· / ἀνδρῶν δὲ πάντων Σωκράτης σοφώτατος ; \textit{idem} Suid.~σ~820 (σοφώτερος pro σοφώτατος) ; Σ Ar.~\textit{Thesm.}~177 σοφοὺς μᾶλλον ἔλεγον καὶ σοφιστὰς τοὺς περὶ μουσικὴν καὶ ποιητικήν ; Σ Eur.~\textit{Med.}~665 εὐεπίφορός ἐστιν ὁ Εὐριπίδης εἰς τὸ λέγειν σοφός καὶ σοφή πρὸς μηδὲν χρήσιμον παραλαμβάνων τὸ ὄνομα ; Hsch.~σ~1371 σοφιστήν· πᾶσαν τέχνην σοφίαν ἔλεγον, καὶ σοφιστὰς τοὺς περὶ μουσικὴν διατρίβοντας καὶ τοὺς μετὰ κιθάρας ᾄδοντας ; Ar.~\textit{Nub.}~520 οὕτω νικήσαιμί τ᾽ ἐγὼ καὶ νομιζοίμην σοφός ; Ar.~\textit{Pax}~700 τί δαί; Κρατῖνος ὁ σοφὸς ἔστιν ; Ar.~\textit{Thesm.}~1130 σκαιοῖσι γάρ τοι καινὰ προσφέρων σοφὰ ; Ar.~\textit{Ran.}~776 ὑπερεμάνησαν κἀνόμισαν σοφώτατον ; Ar.~\textit{Ran.}~1154 δὶς ταὐτὸν ἡμῖν εἶπεν ὁ σοφὸς Αἰσχύλος ; Eur.~\textit{Med.}~665 παῖ σοφοῦ Πανδίονος ; Plu.~\textit{Cons.Apoll.}~2 (\textit{Mor.}~102a–c) κατὰ γὰρ τὸν σοφὸν Εὐριπίδην ; Luc.~\textit{Par.}~4 αὐτὰ μέντοι ταῦτα καὶ ὁ σοφὸς Εὐριπίδης ; Ath.~6.99 (270a–e) παρ᾽ οὗ ὁ σοφὸς Εὐριπίδης λαβὼν ἔφη ; Orig.~\textit{Cels.}~7.6 εἰ γὰρ καὶ ἀνδρῶν ἁπάντων Σωκράτην εἶπε σοφώτατον εἶναι … ἐν τῷ· σοφὸς Σοφοκλῆς, σοφώτερος δ᾽ Εὐριπίδης ; Σ Ar.~\textit{Nub.}~331 σοφιστάς· ἐν ἄλλοις τοὺς πεπαιδευμένους … καὶ θαυμαστὸν οὐδέν, ὅπου μηδὲ τοὺς αὐλητὰς ἀπὸ τοῦ ὀνόματος τούτου προσαγορεύειν ὤκνησαν ; Hsch.~σ~1370 σοφίζεται· σοφόν τι λέγει. καὶ παρακρούεται λόγῳ ; Hsch.~σ~1373 σοφός· φρόνιμος. φιλότεχνος, καὶ ἐξευρετικός. καὶ ὁ τῶν θείων ἔμπειρος ; Phot.~σ~425–426 σοφιστής· ἀπατεών· παρὰ τὸ σοφίζεσθαι, ὅ ἐστι λόγοις ἀπατᾶν· λέγεται δὲ σοφιστὴς καὶ ὁ διδάσκαλος, ὁ σοφίζων· τὸ δὲ παλαιὸν σοφιστὴς ὁ σοφὸς ἐκαλεῖτο. — σοφιστής· πᾶς τεχνίτης· καὶ ὁ τὸ λεκτικὸν ἠσκηκώς}
\lp{177 ὅστις]}{Eur.~\textit{El.}~815–816 Ἐκ τῶν καλῶν κομποῦσι τοῖσι Θεσσαλοῖς / εἶναι τόδ᾽, ὅστις ταῦρον ἀρταμεῖ καλῶς ; Th.~4.18.4 σωφρόνων δὲ ἀνδρῶν οἵτινες τἀγαθὰ ἐς ἀμφίβολον ἀσφαλῶς ἔθεντο ; Eur.~\textit{Phoen.}~509–510 ἀνανδρία γάρ, τὸ πλέον ὅστις ἀπολέσας / τοὔλασσον ἔλαβε ; Th.~3.45.7 ἁπλῶς τε ἀδύνατον καὶ πολλῆς εὐηθείας, ὅστις οἴεται τῆς ἀνθρωπείας φύσεως ὁρμωμένης προθύμως τι πρᾶξαι ἀποτροπήν τινα ἔχειν}
\lp{177 ἐν βραχεῖ]}{Eur.~\textit{IA}~829 καλῶς ἔλεξας ἐν βραχεῖ τὰ καίρια ; Eur.~\textit{Or.}~640–641 τὰ μακρὰ τῶν σμικρῶν λόγων / ἐπίπροσθέν ἐστι καὶ σαφῆ μᾶλλον κλύειν ; Eur.~\textit{Or.}~734 οἰχόμεσθ᾽, ὡς ἐν βραχεῖ σοι τἀμὰ δηλώσω κακά ; Eur.~\textit{Phoen.}~917 ὦ πολλὰ λέξας ἐν βραχεῖ χρόνῳ κακά ; Eur.~fr.~68~N.² μητέρα κατέκταν τὴν ἐμήν, βραχὺς λόγος ; Aesch.~\textit{Supp.}~273–274 μακράν γε μὲν δὴ ῥῆσιν οὐ στέργει πόλις. / βραχὺς τορός θ᾽ ὁ μῦθος ; Soph.~\textit{El.}~673 τέθνηκ᾽ Ὀρέστης· ἐν βραχεῖ ξυνθεὶς λέγω ; Σ Eur.~\textit{Or.}~640 ἐναντιοῦται τῇ Μενελάου βραχυλογίᾳ τὴν μακρηγορίαν τῶν λόγων προκρίνων. οἱ γὰρ Λάκωνες βραχυλόγοι. καὶ Ὅμηρος ; Hsch.~β~1077 βραχύ· ἀντὶ τοῦ οὐδέν (Eur. Tr. 1248), ὀλίγον, μικρόν (Eur. Phoen. 738) ; Eur.~\textit{Hipp.}~704 μακρηγοροῦμεν· οὐκ ἐσωφρόνουν ἐγώ ; Eur.~\textit{Or.}~446 πάντων πρὸς ἀστῶν, ὡς θάνω· βραχὺς λόγος ; Eur.~\textit{Phoen.}~761 τί δεῖ μακρηγορεῖν ; Eur.~\textit{Supp.}~566 βούλῃ συνάψω μῦθον ἐν βραχεῖ σέθεν ; Eur.~\textit{Tro.}~87–88 ἡ χάρις γὰρ οὐ μακρῶν λόγων / δεῖται ; Aesch.~\textit{Pers.}~713 πάντα γάρ, Δαρεῖ᾽, ἀκούσῃ μῦθον ἐν βραχεῖ χρόνῳ ; Soph.~\textit{Phil.}~435–436 λόγῳ δέ σ᾽ ἐν βραχεῖ / τοῦτ᾽ ἐκδιδάξω ; Hsch.~μ~129 μακρηγορία· μακρὰ διάλεξις, ἢ ὑπαγορία ; Eur.~\textit{Hel.}~1522 ναύτας βέβηκεν, ὡς ἂν ἐν βραχεῖ μάθῃς ; Soph.~\textit{El.}~1335–1336 καὶ νῦν ἀπαλλαχθέντε τῶν μακρῶν λόγων ; Hsch.~β~1073 βραχεῖς· ὀλίγοι ; Hsch.~β~1081 βραχέα ἄττα· ὀλίγα τινά ; Suid.~β~526 Βραχυτελῆ: σύντομον, μικρόν}
\lp{178 πολλοὺς … συντέμνειν λόγους]}{Ar.~\textit{Ran.}~1262 εἰς ἓν γὰρ αὐτοῦ πάντα τὰ μέλη ξυντεμῶ ; Eur.~\textit{Hec.}~1180 ἅπαντα ταῦτα συντεμὼν ἐγὼ φράσω ; Eur.~\textit{Heracl.}~784–785 δέσποινα, μύθους σοί τε συντομωτάτους / κλύειν ἐμοί τε τῷδε καλλίστους φέρω ; Eur.~\textit{IA}~1249 ἓν συντεμοῦσα πάντα νικήσω λόγον ; Eur.~\textit{Tro.}~441 ὡς δὲ συντέμω ; Aesch.~\textit{Eum.}~585 πολλαὶ μέν ἐσμεν, λέξομεν δὲ συντόμως ; Aesch.~\textit{Pers.}~698 μή τι μακιστῆρα μῦθον, ἀλλὰ σύντομον λέγων ; Soph.~\textit{Ant.}~446 σὺ δ᾽ εἰπέ μοι μὴ μῆκος, ἀλλὰ συντόμως ; Soph.~\textit{OC}~1579–1582 ξυντομωτάτως μὲν ἂν / τύχοιμι λέξας Οἰδίπουν ὀλωλότα· / ἃ δ᾽ ἦν τὰ πραχθέντ᾽, οὔθ᾽ ὁ μῦθος ἐν βραχεῖ / φράσαι πάρεστιν ; Pl.~\textit{Prt.}~334d σύντεμνέ μοι τὰς ἀποκρίσεις καὶ βραχυτέρας ποίει ; Aeschin.~2.31 νυνὶ δὲ ἴσως ἀνάγκη συντέμνειν τοὺς λόγους ; Σ Ar.~\textit{Thesm.}~178 συντέμνειν λόγους· λέγειν ; Hsch.~σ~2255 συλλήβδην· ὁμοῦ … συλληπτικῶς … ὡς δι᾽ ὀλίγων πολλὰ εἰπεῖν … συντόμως ; Timocl.~fr.~13~Kock τράπεζα β. περιέργως γε, νὴ τὸν οὐρανὸν / ἐξὸν φράσαι τράπεζα συντόμως ; Eur.~\textit{Hec.}~1177 ὡς δὲ μὴ μακροὺς τείνω λόγους ; Aesch.~\textit{Ag.}~629 μακρὸν δὲ πῆμα συντόμως ἐφημίσω ; Σ Eur.~\textit{Phoen.}~917 ὦ πολλὰ λέξας ἐν βραχεῖ: τῷ βραχεῖ ἐναντίον τὸ μακρόν, τῷ δὲ ὀλίγῳ τὸ πολύ ; Hsch.~α~6778 ἀπόφθεγμα· σύντομος λόγος ; Hsch.~κ~2395 κεφαλαίωσον· συντόμισον, συνάγαγε ; Hsch.~σ~2508 συνελόντι φάναι· συντόμως λέγειν, ἢ εἰπεῖν ; \textit{idem} Phot.~σ~766 (φράσαι ; om. λέγειν, ἢ) ; Hsch.~σ~2710 συντεμών· συντελέσας ; Hsch.~σ~2720 συντεταμένου· συγκεκομ⟨μ⟩ένου. συντόμως εἰρημένου ; Hsch.~σ~2725 συντόμως· διὰ βραχέων ; [Eur.]~\textit{Rh.}~450 πρὸς οἶκον εἶμι, συντεμὼν τοὺς σοὺς πόνους ; Aesch.~\textit{Eum.}~227 τιμὰς σὺ μὴ σύντεμνε τὰς ἐμὰς λόγῳ ; Aesch.~\textit{Eum.}~415 πεύσῃ τὰ πάντα συντόμως, Διὸς κόρη ; Soph.~\textit{OT}~710 φανῶ δέ σοι σημεῖα τῶνδε σύντομα ; Σ Eur.~\textit{Hec.}~1180 συντεμών: συντόμως συλλαβών ; Hsch.~ε~3155 ἐν κεφαλαίῳ· ἐν συντόμῳ ; Suid.~σ~1633 Συντέμνω· αἰτιατικῇ}
\lp{178 οἷός τε]}{Ar.~\textit{Eq.}~343 ὁτιὴ λέγειν οἷός τε κἀγὼ καὶ καρυκοποιεῖν ; Hsch.~ο~1935 οὐχ οἷός τε εἰπεῖν· μὴ δυνάμενος εἰπεῖν ; Hsch.~ο~372 οἷόν τε· δυνατόν}
\lp{179 ἐγὼ δὲ καινῇ ξυμφορᾷ πεπληγμένος]}{\textit{Test.} Suid.~κ~2556 Κοινῇ δὲ ἀντὶ τοῦ ἀνθρωπίνῃ, μεγάλῃ καὶ θαυμαστῇ. ἐγὼ δὲ κοινῇ ξυμφορᾷ πεπληγμένος ; Σ Ar.~\textit{Thesm.}~179 μεγάλῃ καὶ θαυμαστῇ ; Hsch.~κ~3253b κοινῇ· δημοσίᾳ}
\lp{179 καινῇ ξυμφορᾷ]}{Eur.~\textit{Med.}~556 καινῆς δὲ νύμφης ἱμέρῳ πεπληγμένος ; Eur.~\textit{Tro.}~1126–1127 καινάς τινας / Πηλέως ἀκούσας συμφοράς ; Eur.~\textit{HF}~1177 ἀλλ᾽ ἄλλο πού τι καινὸν εὑρίσκω κακόν ; Eur.~\textit{Hec.}~689–690 ἄπιστ᾽ ἄπιστα, καινὰ καινὰ δέρκομαι. / ἕτερα δ᾽ ἀφ᾽ ἑτέρων κακὰ κακῶν κυρεῖ ; Eur.~\textit{Med.}~705 πρὸς τοῦ; τόδ᾽ ἄλλο καινὸν αὖ λέγεις κακόν ; Eur.~\textit{Med.}~1120 δείκνυσιν ὥς τι καινὸν ἀγγελεῖ κακόν ; Eur.~\textit{Tro.}~1259 καινόν τι κακὸν προσέσεσθαι ; Suid.~κ~1182 Καινῶς: ἀντὶ τοῦ παραδόξως ; Hsch.~κ~244 καινός· νέος, νεαρός ; \textit{idem} Suid.~κ~1173 (ὁ νέος)}
\lp{179 καινῇ (ἀεὶ Λιβύη φέρει τι καινόν)]}{Arist.~\textit{HA}~606b20 λέγεται δέ τις παροιμία, ὅτι “ἀεὶ Λιβύη φέρει τι καινόν” ; \textit{idem} Zen.~2.51 (ἀεὶ φέρει τι Λιβύη καινὸν κακόν ; add. ἐπὶ τῶν κακουργοτάτων, καὶ ἀεὶ προσεξευρισκόντων νεώτερόν τι κακόν) ; Anaxil.~fr.~27~Kock ἡ μουσικὴ δ᾽ ὥσπερ Λιβύη πρὸς τῶν θεῶν / αἰεί τι καινὸν κατ᾽ ἐνιαυτὸν θηρίον τίκτει ; Tosi n°~1138 Quid novi ex Africa ; Erasm.~\textit{Adag.}~3.7.10 Semper Africa novi aliquid apportat. Huic simillimum est illud Plinianum … Libyam semper aliquid novi adferre}
\lp{179 ξυμφορᾷ πεπληγμένος]}{Eur.~\textit{Alc.}~404–405 ὥστ᾽ ἐγὼ / καὶ σφὼ βαρείᾳ συμφορᾷ πεπλήγμεθα ; Eur.~\textit{Alc.}~856 καίπερ βαρείᾳ συμφορᾷ πεπληγμένος ; Aesch.~\textit{Cho.}~31 ξυμφοραῖς πεπληγμένων ; Hsch.~ξ~114 ξυμφορά· συντυχία. καὶ ἐπὶ ἀγαθοῦ τάσσεται ⟨ὡς παρ᾽ Αἰσχύλῳ ἐν Καβείροις⟩, καὶ ἐπὶ κακοῦ ⟨παρὰ Σοφοκλεῖ⟩ ; \textit{idem} Phot.~ξ~61 (om. καὶ ἐπὶ ἀγαθοῦ … Σοφοκλεῖ) ; Ar.~\textit{Ach.}~1218 εἰλιγγιῶ κάρα λίθῳ πεπληγμένος ; Eur.~\textit{Hipp.}~342 ἔκ τοι πέπληγμαι· ποῖ προβήσεται λόγος ; Eur.~\textit{Or.}~1–2 Οὐκ ἔστιν οὐδὲν δεινὸν ὧδ᾽ εἰπεῖν ἔπος / οὐδὲ πάθος οὐδὲ ξυμφορὰ θεήλατος ; Aesch.~\textit{Ag.}~1343 ὤμοι, πέπληγμαι καιρίαν πληγὴν ἔσω ; Aesch.~\textit{Pers.}~1008 πεπλήγμεθ᾽ οἵᾳ δι᾽ αἰῶνος τύχᾳ ; Hsch.~θ~174 θεήλατον· ὑπὸ θεοῦ … πεπληγμένον, ἢ … πεμφθέν ; Hsch.~σ~2356 συμφορά· συντυχία. σύμπτωμα. ἀτυχία ; Hsch.~τ~1233 τραγικὸν πάθος· μεστὸν συμφορᾶς ; Suid.~σ~1407 Συμφορά: ἀτυχία, κακῶν συντυχία ; Suid.~σ~1408 Συμφορά: τῶν μέσων ἐστί … δηλοῖ δὲ ἡ συμφορὰ καὶ τὸ καλὸν καὶ τὸ κακόν ; Eur.~\textit{Hel.}~1190–1191 πότερον ἐννύχοις πεπεισμένη / στένεις ὀνείροις ; Eur.~\textit{IT}~549 τέθνηκε; ποίᾳ συμφορᾷ ; Eur.~\textit{Med.}~1387 Ἀργοῦς κάρα σὸν λειψάνῳ πεπληγμένος ; Aesch.~\textit{Ag.}~544 τῶν ἀντερώντων ἱμέρῳ πεπληγμένοι ; Aesch.~\textit{Ag.}~1204 μῶν καὶ θεός περ ἱμέρῳ πεπληγμένος ; Aesch.~\textit{Pers.}~846 μάλιστα δ᾽ ἥδε συμφορὰ δάκνει ; Soph.~\textit{OT}~99 τίς ὁ τρόπος τῆς ξυμφορᾶς ; Soph.~\textit{Trach.}~931 πλευρὰν ὑφ᾽ ἧπαρ καὶ φρένας πεπληγμένην ; Hsch.~π~1399 πένθος· συμφορά ; Hsch.~π~2417 πλαγκτέ· παράφρων, καὶ πεπληγμένε τὴν διάνοιαν ; Suid.~ξ~110 Ξυμφοράν φησι Θουκυδίδης τὴν ἀπὸ περιστάσεως δυστυχίαν, καὶ οὐκ ἐξ οἰκείας ἀβουλίας κακοπραγίαν}
\lp{180 ἱκέτης ἀφῖγμαι πρὸς σέ]}{Ar.~\textit{Av.}~120 ταῦτ οὖν ἱκέται νὼ πρὸς σὲ δεῦρ᾽ ἀφίγμεθα ; Eur.~\textit{Heracl.}~94 ἱκέται σέθεν τε καὶ πόλεως ἀφιγμένοι ; Eur.~\textit{Supp.}~114 Θησεῦ, σὸς ἱκέτης καὶ πόλεως ἥκω σέθεν ; Soph.~\textit{OC}~634 ἔπειτα δ᾽ ἱκέτης δαιμόνων ἀφιγμένος ; Soph.~\textit{OT}~920 ἱκέτις ἀφῖγμαι τοῖσδε σὺν κατεύγμασιν ; Eur.~\textit{Med.}~709–713 ἀλλ᾽ ἄντομαί σε τῆσδε πρὸς γενειάδος / γονάτων τε τῶν σῶν ἱκεσία τε γίγνομαι, / οἴκτιρον οἴκτιρόν με τὴν δυσδαίμονα ; Eur.~\textit{Or.}~382–383 τῶν σῶν δὲ γονάτων πρωτόλεια θιγγάνω / ἱκέτης ; Hsch.~α~8689 Ἀφίκτορα· τὸν ἱκέσιον Δία ; Hsch.~ι~461 ἱκέτην· πρόσφυγα, ἀξιοῦντα, δεόμενον ; Ar.~\textit{Lys.}~1138–1139 οὐκ ἴσθ᾽ ὅτ᾽ ἐλθὼν δεῦρο Περικλείδας ποτὲ / ὁ Λάκων Ἀθηναίων ἱκέτης καθέζετο ; Hsch.~α~8666 ἀφῖγμαι· παραγέγονα ; \textit{idem} Suid.~α~4610 ; Hsch.~α~8685 ἀφικόμην· ἀφῖγμαι ; Hsch.~ι~462 †ἱκέτης· παῖς, δοῦλος. παράσιτος ; Suid.~π~2830 Προστρόπαιος: προστετραμμένος, ἱκέτης παρακαλῶν}
\lp{180 τοῦ χρείαν ἔχων;]}{Ar.~\textit{Ach.}~1026 εἶτα νυνὶ τοῦ δέει ; Ar.~\textit{Plu.}~54 καὶ τοῦ δεόμενος ἦλθε μετὰ νῷν ἐνθαδί ; Ar.~\textit{Plu.}~827–828 ἔπειτα τοῦ δέει; — πρὸς τὸν θεὸν / ἥκω ; Eur.~\textit{Andr.}~368 εὖ δ᾽ ἴσθ᾽, ὅτου τις τυγχάνει χρείαν ἔχων ; Eur.~\textit{Med.}~1319 εἰ δ᾽ ἐμοῦ χρείαν ἔχεις ; Eur.~\textit{Supp.}~115 τί χρῆμα θηρῶν καὶ τίνος χρείαν ἔχων ; Eur.~fr.~231~N.² ἡμῶν τί δῆτα τυγχάνεις χρείαν ἔχων; / πατέρων γὰρ ἐσθλῶν ἐλπίδας δίδως γεγώς ; Aesch.~\textit{PV}~168–170 ἦ μὴν ἔτ᾽ ἐμοῦ, καίπερ κρατεραῖς / ἐν γυιοπέδαις αἰκιζομένου, / χρείαν ἕξει μακάρων πρύτανις ; Ar.~\textit{Thesm.}~1164 χρείᾳ δὲ ποίᾳ τόνδ᾽ ἐπεσφέρεις λόγον ; Eur.~\textit{Heracl.}~95 ἱκέται σέθεν τε καὶ πόλεως ἀφιγμένοι. / τί χρέος ; Aesch.~\textit{Cho.}~481 κἀγώ, πάτερ, τοιάνδε σου χρείαν ἔχω ; Soph.~\textit{Ichn.}~257 ὅτου μάλιστα πράγματος χρείαν ἔχεις ; X.~\textit{Hell.}~3.4.5 πρῶτον μὲν Τισσαφέρνης πέμψας ἤρετο αὐτὸν τίνος δεόμενος ἥκοι ; Σ Ar.~\textit{Plu.}~827 τοῦ δέει : Τίνος χρείαν ἔχεις ; Suid.~δ~299 Δεῖ: πρέπει … ἀντὶ τοῦ χρή. καὶ ἀντὶ τοῦ χρεία ἐστί ; Hsch.~χ~682 χρεῖαι· ἔνδειαι}
\lp{180 τοῦ χρείαν ἔχων; (χρεία διδάσκει)]}{Eur.~fr.~715~N.² οὔ τἄρ᾽ Ὀδυσσεύς ἐστιν αἱμύλος μόνος· / χρεία διδάσκει, κἂν βραδύς τις ᾖ, σοφόν ; \textit{idem} Apostol.~18.41e (ἢ σοφόν ; om. ᾖ ; add. Εὐριπίδου), Men.~fr.~263~Kock (κἂν ἄμουσος ᾖ, σοφὸν / Καρχηδόνιον) ; Eur.~\textit{El.}~375–376 ἀλλ᾽ ἔχει νόσον / πενία, διδάσκει δ᾽ ἄνδρα τῇ χρείᾳ κακόν}
\lp{181 μέλλουσί μ᾽ αἱ γυναῖκες ἀπολεῖν τήμερον]}{Ar.~\textit{Ach.}~440 δεῖ γάρ με δόξαι πτωχὸν εἶναι τήμερον ; Ar.~\textit{Thesm.}~76–77 τῇδε θἠμέρᾳ κριθήσεται / εἴτ᾽ ἔστ᾽ ἔτι ζῶν εἴτ᾽ ἀπόλωλ᾽ Εὐριπίδης ; Ar.~\textit{Thesm.}~83–84 κἀν Θεσμοφόροιν μέλλουσι περί μου τήμερον / ἐκκλησιάζειν ἐπ᾽ ὀλέθρῳ ; Eur.~\textit{Or.}~857–858 ψήφῳ Πελασγῶν σὸν κασίγνητον θανεῖν / καὶ σέ, ὦ τάλαιν᾽, ἔδοξε τῇδ᾽ ἐν ἡμέρᾳ ; Ar.~\textit{Eq.}~68 εἰ μή μ᾽ ἀναπείσετ᾽, ἀποθανεῖσθε τήμερον ; Ar.~\textit{Pax}~243 καὶ πολλοδεκάκις, ὡς ἀπολεῖσθε τήμερον ; Soph.~\textit{El.}~379–380 μέλλουσι γάρ σ᾽, εἰ τῶνδε μὴ λήξεις γόων, / ἐνταῦθα πέμψειν ; Hsch.~μ~743 μέλλει· φαίνεται, ἔοικε … δοκεῖ}
\lp{181 αἱ γυναῖκες ἀπολεῖν (Λήμνιον κακόν)]}{Eur.~\textit{Hec.}~883–887 καὶ πῶς γυναιξὶν ἀρσένων ἔσται κράτος; / δεινὸν τὸ πλῆθος σὺν δόλῳ τε δύσμαχον. / … τί δ᾽; οὐ γυναῖκες εἷλον Αἰγύπτου τέκνα / καὶ Λῆμνον ἄρδην ἀρσένων ἐξῴκισαν ; Diogenian.~4.5 Γυναικῶν ὄλεθροι: ἐπὶ τῶν οἰκτρῶς ἀναιρουμένων ; \textit{idem} Apostol.~5.78, Suid.~γ~499 (γυναικὸς ὄλεθρος ; add. παρόσον πολλοὶ δεινῶς ὑπὸ γυναικῶν ἀπώλοντο) ; Ar.~\textit{Lys.}~299–300 κἄστιν γε Λήμνιον τὸ πῦρ / τοῦτο πάσῃ μηχανῇ ; Aesch.~\textit{Cho.}~631–634 κακῶν δὲ πρεσβεύεται τὸ Λήμνιον / λόγῳ … ᾔκασεν δέ τις / τὸ δεινὸν αὖ Λημνίοισι πήμασιν ; Hdt.~6.138.4 αἱ γυναῖκες τοὺς ἅμα Θόαντι ἄνδρας σφετέρους ἀποκτείνασαι, νενόμισται ἀνὰ τὴν Ἑλλάδα τὰ σχέτλια ἔργα πάντα Λήμνια καλέεσθαι ; Σ Eur.~\textit{Hec.}~887 καὶ μετὰ ταῦτα καὶ αἱ Λήμνιαι γυναῖκες τοὺς σὺν Θόαντι πάντας ἀπέκτειναν. δι᾽ ἀμφότερα οὖν ταῦτα ἡ παροιμία ἐδόθη· τὰ Λήμνια κακά ; Suid.~ε~3695 κύνας ἐπαφεῖναι αὐτῷ. οἱ δὲ ἱστόρησαν οὐχ ὑπὸ κυνῶν, ἀλλ᾽ ὑπὸ γυναικῶν νύκτωρ διασπασθῆναι, πορευόμενον ἀωρὶ πρὸς Κρατερὸν τὸν ἐρώμενον Ἀρχελάου ; Zen.~4.91 Λήμνιον κακόν: παροιμία, ἣν διαδοθῆναι φᾶσὶν ἀπὸ τῶν παρανομηθέντων εἰς τοὺς ἄνδρας ἐν Λήμνῳ ὑπὸ τῶν γυναικῶν ; Diogenian.~7.52 Προμέρου κύνες: τοῦτόν φασιν οἰκέτην ὄντα βασιλικὸν μισῆσαι καθ᾽ ὑπερβολὴν τῶν τραγῳδιῶν τὸν ποιητὴν Εὐριπίδην, ἐπαφεῖναί τε αὐτῷ κύνας ἀγρίους, οἵ τοῦτον κατεθοινήσαντο ; Erasm.~\textit{Adag.}~1.9.27 Lemnium malum. Λήμνιον κακόν. i. e. Lemnium malum, de malis immensis ac miserandis}
\lp{181 τήμερον]}{Hsch.~τ~789 τήμερον· Ἀττικοὶ τήμερος καὶ τήμερα λέγουσιν ; \textit{idem} Hdn.~\textit{Orth.}~3.19.62.1 (om. τήμερος) ; Moer. p.~210,~14–16~Bekker τήμερον εἶναι Ἀττικοί, ὡς Νεφέλαις … τὸ σήμερον ἔχον Ἕλληνες. τήμερον Ἀττικοί, σήμερον Ἕλληνες ; Suid.~τ~509 Τήμερος: ὁ σημερινός … τὸ δὲ τήμερον ἐπὶ χρόνου λέγεται ; Eust.~813.42–44 (ad Il.~10.409) οὐ πάνυ εὕρηται χαίρων ὁ ποιητὴς τῇ προφορᾷ τοῦ τ̄ ἀντὶ τοῦ σ̄ κατὰ τοὺς ὕστερον Ἀττικούς, οἳ Ὁμήρου σήμερον λέγοντος τήμερον αὐτοί φασιν ; Hsch.~σ~503 σήμερον· αὐτίκα, ταχύ ; Suid.~τ~496 ὥσπερ τὸ σήμερον, τήμερον}
\lp{182 τοῖς Θεσμοφορίοις]}{Ar.~\textit{Av.}~1519 ἀλλ᾽ ὡσπερεὶ Θεσμοφορίοις νηστεύομεν ; \textit{idem} Suid.~θ~270 (s.v. Θεσμοφόρια· ὄνομα ἑορτῆς ; ὥσπερ ἐν θεσμοφορίοισι νηστεύομεν ἂν θυηλῶν. Ἀριστοφάνης) ; Ar.~\textit{Thesm.}~80 ἐπεὶ τρίτη ᾽στὶ Θεσμοφορίων ἡ μέση ; Lys.~1.20 καὶ ὡς Θεσμοφορίοις ἐμοῦ ἐν ἀγρῷ ὄντος ᾤχετο εἰς τὸ ἱερὸν μετὰ τῆς μητρὸς τῆς ἐκείνου ; Hsch.~θ~383 θεσμοφόρια· τὰ μυστήρια τῆς Δήμητρος. ἢ ὄνομα ἑορτῆς ; Plu.~\textit{Mor.}~298b (\textit{Quaest.\,Gr.}~31) διὰ τί τοῖς θεσμοφορίοις αἱ τῶν Ἐρετριέων γυναῖκες οὐ πρὸς πῦρ ἀλλὰ πρὸς ἥλιον ὀπτῶσι τὰ κρέα ; Apollod.~1.5.1 γραῖά τις Ἰάμβη σκώψασα τὴν θεὸν ἐποίησε μειδιᾶσαι. διὰ τοῦτο ἐν τοῖς θεσμοφορίοις τὰς γυναῖκας σκώπτειν λέγουσιν ; Σ Luc.~\textit{DMeretr.}~2.1 Θεσμοφορία ἑορτὴ Ἑλλήνων μυστήρια περιέχουσα, τὰ δὲ αὐτὰ καὶ Σκιρροφορία καλεῖται ; Hsch.~α~5234 ἄνοδος· ἀνάβασις. ἡ ἑνδέκατη τοῦ Πυανεψιῶνος, ὅτε αἱ γυναῖκες ἀνέρχονται εἰς Θεσμοφόριον, οὕτω καλεῖται ; Suid.~σ~325 οὕτως οὖν καὶ τῶν γυναικῶν μελλουσῶν ἐν τοῖς Θεσμοφορίοις ἐκκλησιάζειν, σημεῖον ἐτέθη. Ἀριστοφάνης· ὡς τὸ τῆς ἐκκλησίας σημεῖον ἐν τῷ Θεσμοφόρῳ φαίνεται ; Ar.~\textit{Thesm.}~558 ὥς τ᾽ αὖ τὰ κρέ᾽ ἐξ Ἀπατουρίων ταῖς μαστροποῖς διδοῦσαι ; Ar.~\textit{Thesm.}~834 προεδρίαν τ᾽ αὐτῇ δίδοσθαι Στηνίοισι καὶ Σκίροις ; Luc.~\textit{DMeretr.}~2.1 εἶδον γὰρ αὐτὴν ἔναγχος ἐν τοῖς Θεσμοφορίοις μετὰ τῆς μητρός ; Hsch.~ζ~145 ζημία· θυσία τις ἀποδιδομένη ὑπὲρ τῶν γινομένων ἐν Θεσμοφορίοις ; Hsch.~θ~384 θεσμοφόρος· ὁ θεσμοὺς φέρων, ἤτοι νόμους}
\lp{182 ὅτι κακῶς αὐτὰς λέγω]}{Ar.~\textit{Lys.}~368–369 οὐκ ἔστ᾽ ἀνὴρ Εὐριπίδου σοφώτερος ποιητής· / οὐδὲν γὰρ οὕτω θρέμμ᾽ ἀναιδές ἐστιν ὡς γυναῖκες ; Ar.~\textit{Thesm.}~85 ὁτιὴ τραγῳδῶ καὶ κακῶς αὐτὰς λέγω ; Ar.~\textit{Thesm.}~389–394 τί γὰρ οὗτος ἡμᾶς οὐκ ἐπισμῇ τῶν κακῶν; / ποῦ δ᾽ οὐχὶ διαβέβληχ᾽ … / τὰς μοιχοτρόπους, τὰς ἀνδρεραστίας καλῶν, / τὰς οἰνοπότιδας, τὰς προδότιδας, τὰς λάλους ; Ar.~\textit{Ran.}~1043–1044 ἀλλ᾽ οὐ μὰ Δί᾽ οὐ Φαίδρας ἐποίουν πόρνας οὐδὲ Σθενεβοίας, / οὐδ᾽ οἶδ᾽ οὐδεὶς ἥντιν᾽ ἐρῶσαν πώποτ᾽ ἐποίησα γυναῖκα ; Eur.~\textit{Hec.}~1178–1182 εἴ τις γυναῖκας τῶν πρὶν εἴρηκεν κακῶς / ἢ νῦν λέγων ἔστιν τις ἢ μέλλει λέγειν, / ἅπαντα ταῦτα συντεμὼν ἐγὼ φράσω· / γένος γὰρ οὔτε πόντος οὔτε γῆ τρέφει / τοιόνδ᾽ ; Eur.~\textit{Hipp.}~664–666 μισῶν δ᾽ οὔποτ᾽ ἐμπλησθήσομαι / γυναῖκας, οὐδ᾽ εἴ φησί τίς μ᾽ ἀεὶ λέγειν· / ἀεὶ γὰρ οὖν πώς εἰσι κἀκεῖναι κακαί ; Eur.~fr.~499~N.² μάτην ἄρ᾽ εἰς γυναῖκας ἐξ ἀνδρῶν ψόγος / ψάλλει, κενὸν τόξευμα, καὶ κακῶς λέγει· / αἳ δ᾽ εἴσ᾽ ἀμείνους ἀρσένων, ἐγὼ λέγω ; Ath.~13.557b–e εἰπόντος Σοφοκλεῖ τινος ὅτι μισογύνης ἐστὶν Εὐριπίδης, ἔν γε ταῖς τραγῳδίαις, ἔφη ὁ Σοφοκλῆς· ἐπεὶ ἔν γε τῇ κλίνῃ φιλογύνης ; Vit. Eur.~1 λέγουσι δὲ καὶ ὅτι αἱ γυναῖκες διὰ τοὺς ψόγους, οὓς ἐποίει εἰς αὐτὰς διὰ τῶν ποιημάτων, τοῖς Θεσμοφορίοις ἐπέστησαν αὐτῷ βουλόμεναι ἀνελεῖν, ἐφείσαντο δ᾽ αὐτοῦ πρῶτον μὲν διὰ τὰς Μούσας, ἔπειτα δὲ βεβαιωσαμένου μηκέτι αὐτὰς κακῶς ἐρεῖν ; Ar.~\textit{Ach.}~503 ξένων παρόντων τὴν πόλιν κακῶς λέγω ; Ar.~\textit{Thesm.}~539 γυνὴ γυναῖκας οὖσα μὴ κακῶς λέγειν τὸ λοιπόν ; Hes.~\textit{Op.}~721 εἰ δὲ κακὸν εἴποις, τάχα κ᾽ αὐτὸς μεῖζον ἀκούσαις ; Gell.~15.20.6 Mulieres fere omnes in maiorem modum exosus fuisse dicitur, sive quod natura abhorruit a mulierum coetu sive quod duas simul uxores habuerat ; Hsch.~κ~301 κακήγορος· κακῶς λέγων … κακολόγος ; Suid.~ε~3695 σκυθρωπὸς δὲ ἦν τὸ ἦθος καὶ ἀμειδὴς καὶ φεύγων τὰς συνουσίας· ὅθεν καὶ μισογύνης ἐδοξάσθη ; App. Prov.~2.17 Εἰπὼν ἃ θέλεις, ἀντάκουε ἃ μὴ θέλεις: ἐπὶ τῶν ἀκουόντων κακά, εἰπόντων τοιαῦτα ; Soph.~fr.~843~N.² φιλεῖ δὲ πολλὴν γλῶσσαν ἐκχέας μάτην / ἄκων ἀκούειν οὓς ἑκὼν εἶπεν λόγους ; Chilon ap. D.L.~1.69 μὴ κακολογεῖν τοὺς πλησίον· εἰ δὲ μή, ἀκούσεσθαι ἐφ᾽ οἷς λυπήσεσθαι}
\lp{183 τίς οὖν παρ᾽ ἡμῶν ἐστιν ὠφέλειά σοι;]}{Ar.~\textit{Thesm.}~196 ἡμᾶς ὑφέξειν. καὶ γὰρ ἂν μαινοίμεθ᾽ ἄν ; Eur.~fr.~231~N.² ἡμῶν τί δῆτα τυγχάνεις χρείαν ἔχων ; Soph.~\textit{El.}~944 ἀλλ᾽ εἴ τις ὠφέλειά γ᾽, οὐκ ἀπώσομαι ; Ar.~\textit{Nub.}~648 τί δέ μ᾽ ὠφελήσουσ᾽ οἱ ῥυθμοὶ πρὸς τἄλφιτα ; Eur.~\textit{Andr.}~539 τοῖς γὰρ ἐμοῖσιν γέγον᾽ ὠφελία ; Eur.~fr.~78~N.² γυναῖκα καὶ ὠφελίαν / καὶ νόσον ἀνδρὶ φέρειν / μεγίσταν ἐδίδαξα τὠμῷ λόγῳ ; Soph.~\textit{OC}~401 ἡ δ᾽ ὠφέλησις τίς θύρασι κειμένου ; Hdn.~\textit{Orth.} s.v.~ὠφέλεια ὠφέλεια: ει δίφθογγος ἡ παράδοσις. ποιητικώτερον δὲ διὰ τοῦ ι καὶ παροξύνεται. παρὰ Σοφοκλεῖ δὲ καὶ ει δίφθογγος καὶ προπαροξύνεται ὡς προμήθεια ; Hdt.~7.158.2 ἀπ᾽ ὧν ὑμῖν μεγάλαι ὠφελίαι τε καὶ ἐπαυρέσιες γεγόνασι ; Th.~1.28.3 φίλους ποιεῖσθαι οὓς οὐ βούλονται ἑτέρους τῶν νῦν ὄντων μᾶλλον ὠφελίας ἕνεκα ; Pl.~\textit{Phdr.}~238e τίς ὠφελία ἢ βλάβη ἀπό τε ἐρῶντος καὶ μὴ τῷ χαριζομένῳ ἐξ εἰκότος συμβήσεται ; X.~\textit{Mem.}~1.4.5 ὁ ἐξ ἀρχῆς ποιῶν ἀνθρώπους ἐπ᾽ ὠφελείᾳ προσθεῖναι αὐτοῖς δι᾽ ὧν αἰσθάνονται ἕκαστα ; Lys.~7.13 ἀποφαίνοντας ἥτις ὠφέλεια τοῖς ἀδικήσασιν ἐγίγνετο ; Hsch.~ο~880 ὄνησις· ἰσχύς. ὠφέλεια}
\lp{184–186 ἐγκαθεζόμενος λάθρᾳ … ὑπεραποκρίνῃ μου]}{Ar.~\textit{Thesm.}~90–92 ἐκκλησιάσοντ᾽ ἐν ταῖς γυναιξὶ κἂν δέῃ / λέξονθ᾽ ὑπὲρ ἐμοῦ. — πότερα φανερῶς ἢ λάθρᾳ; / λάθρᾳ, στολὴν γυναικὸς ἠμφιεσμένον}
\lp{184–185 λάθρᾳ / ἐν ταῖς γυναιξίν]}{Ar.~\textit{Thesm.}~588–591 ἐκεῖνος εἴη τῶν λόγων κατάσκοπος. / καὶ πῶς λέληθεν ἐν γυναιξὶν ὢν ἀνήρ; / ἀφηῦσεν αὐτὸν κἀπέτιλ᾽ Εὐριπίδης / καὶ τἄλλ᾽ ἅπανθ᾽ ὥσπερ γυναῖκ᾽ ἐσκεύασεν ; Plu.~\textit{Caes.}~10.1 ὁ Κλώδιος οὔπω γενειῶν καὶ διὰ τοῦτο λήσειν οἰόμενος ἐσθῆτα καὶ σκευὴν ψαλτρίας ἀναλαβὼν ἐχώρει, νέᾳ γυναικὶ τὴν ὄψιν ἐοικώς ; Paus.~4.17.1 τὰς γυναῖκας ἀγούσας ἑορτὴν ἀμύνεσθαι τῶν γυναικῶν οὐκ ἄνευ τῆς θεοῦ προαχθεισῶν λαμβάνουσιν οἱ πολλοὶ τῶν Μεσσηνίων τραύματα μαχαίραις τε … τὸν δὲ Ἀριστομένην τύπτουσαι ταῖς δᾳσὶ ζῶντα αἱροῦσιν ; Suid.~θ~272 Θεσμοφόρος. ὅτι Βάττος ὁ Κυρήνην κτίσας τῆς Θεσμοφόρου τὰ μυστήρια ἐγλίχετο μαθεῖν καὶ προσῆγε βίαν λίχνοις ὀφθαλμοῖς χαριζόμενος ; Bion, \textit{Epith.\,Achill.\,et\,Deid.}~15–17 λάνθανε δ᾽ ἐν κώραις Λυκομηδίσι μοῦνος Ἀχιλλεύς, / εἴρια δ᾽ ἀνθ᾽ ὅπλων ἐδιδάσκετο … / … ἐφαίνετο δ᾽ ἠΰτε κώρα ; Hdt.~6.134.2 ἕρκος θεσμοφόρου Δήμητρος ὑπερθορεῖν, οὐ δυνάμενον τὰς θύρας ἀνοῖξαι, ὑπερθορόντα δὲ ἰέναι ἐπὶ τὸ μέγαρον ὅ τι δὴ ποιήσοντα ἐντός ; Hsch.~κ~1084 καταδύμεναι· λάθρα ὑπεισελθεῖν (Γ 241)}
\lp{184 ἡ πᾶσ᾽]}{Ar.~\textit{Nub.}~204 πότερα τὴν κληρουχικήν; / οὔκ, ἀλλὰ τὴν σύμπασαν ; Ar.~\textit{Pax}~373 νῦν ἆρά με — ἅπασ᾽ ἀνάγκη ᾽στ᾽ ἀποθανεῖν ; Soph.~\textit{El.}~301–302 ὁ πάντ᾽ ἄναλκις οὗτος, ἡ πᾶσα βλάβη, / ὁ σὺν γυναιξὶ τὰς μάχας ποιούμενος ; Soph.~\textit{Phil.}~622 ἦ κεῖνος, ἡ πᾶσα βλάβη}
\lp{184 ἐγκαθεζόμενος λάθρᾳ]}{Ar.~\textit{Thesm.}~600 λέληθεν ἡμᾶς κρυπτὸς ἐγκαθήμενος ; Ar.~\textit{Ran.}~1522–1523 μηδέποτ᾽ ἐς τὸν θᾶκον τὸν ἐμὸν / μηδ᾽ ἄκων ἐγκαθεδεῖται ; Ar.~\textit{Eccl.}~21–23 καταλαβεῖν δ᾽ ἡμᾶς ἕδρας, / … δεῖ τὰς ἑτέρας πως κἀγκαθεζομένας λαθεῖν ; Ar.~\textit{Eccl.}~98 ἢν δ᾽ ἐγκαθεζώμεσθα πρότεραι, λήσομεν ; Eur.~\textit{Ba.}~816–817 σάφ᾽ ἴσθι, σιγῇ γ᾽ ὑπ᾽ ἐλάταις καθήμενος. / ἀλλ᾽ ἐξιχνεύσουσίν σε, κἂν ἔλθῃς λάθρᾳ ; Hsch.~ε~179 ἐγκάθετος· δόλιος ; Ar.~\textit{Ach.}~343 ἀλλ᾽ ὅπως μὴ ν᾽ τοῖς τρίβωσιν ἐγκάθηνταί που λίθοι ; Ar.~\textit{Vesp.}~1114 ἀλλὰ γὰρ κηφῆνες ἡμῖν εἰσιν ἐγκαθήμενοι ; Ar.~\textit{Thesm.}~687 οὐχ ὁρῶμεν γοῦν ἔτ᾽ ἄλλον οὐδέν᾽ ἐγκαθήμενον ; Eur.~\textit{Ba.}~954–956 ἐλάταισιν δ᾽ ἐμὸν κρύψω δέμας. / κρύψῃ σὺ κρύψιν ἥν σε κρυφθῆναι χρεών, / ἐλθόντα δόλιον μαινάδων κατάσκοπον ; Th.~3.1.2 καὶ ἐγκαθεζόμενοι ἐδῄουν τὴν γῆν ; \textit{idem} Th.~4.2.1 ; Hsch.~ε~3146 ἐνκάθετος· ἐνεδρεύων … ἐπιτηρῶν, ἐντόπιος (Iob 31, 9) ; Hsch.~λ~103 λάθρᾳ· κρυφίως, ἠρέμα, ἡσύχως ; Suid.~ε~65 Ἐγκάθετος: δόλιος, κατάσκοπος. καὶ Ἐγκάθηνται, ἀντὶ τοῦ ἐγκεκρυμμένοι εἰσίν· οἱ γὰρ κρυπτόμενοι ὑποκάθηνται ; Hsch.~ε~3273 ἐν παραβύστῳ· ἐν μυστηρίῳ … ἐν κρυπτῷ, λεληθότως … λάθρα ; Hsch.~κ~3076 κλωπωμένη· λάθρα καὶ ἀψόφως περιοῦσα}
\lp{185 ὡς δοκῶν εἶναι γυνή]}{Ar.~\textit{Thesm.}~97–98 ἐγὼ γὰρ οὐχ ὁρῶ / ἄνδρ᾽ οὐδέν᾽ ἐνθάδ᾽ ὄντα, Κυρήνην δ᾽ ὁρῶ ; Ar.~\textit{Thesm.}~266–268 ἁνὴρ μὲν ἡμῖν οὑτοσὶ καὶ δὴ γυνὴ / τό γ᾽ εἶδος· ἢν λαλῇς δ᾽, ὅπως τῷ φθέγματι / γυναικιεῖς εὖ καὶ πιθανῶς ; Eur.~\textit{Ba.}~821–823 στεῖλαί νυν ἀμφὶ χρωτὶ βυσσίνους πέπλους. / τί δὴ τόδ᾽; ἐς γυναῖκας ἐξ ἀνδρὸς τελῶ; / μή σε κτάνωσιν, ἢν ἀνὴρ ὀφθῇς ἐκεῖ ; Eur.~\textit{HF}~985 ἔπτηξε κρηπῖδ᾽ ὡς λεληθέναι δοκῶν ; Eur.~fr.~698~N.² δεῖ γάρ με δόξαι πτωχὸν [εἶναι τήμερον], / εἶναι μὲν ὅσπερ εἰμί, φαίνεσθαι δὲ μή ; \textit{idem} Ar.~\textit{Ach.}~440–441 ; Ar.~\textit{Thesm.}~141–143 πότερον ὡς ἀνὴρ τρέφει; / καὶ ποῦ πέος; ποῦ χλαῖνα; ποῦ Λακωνικαί; / ἀλλ᾽ ὡς γυνὴ δῆτ᾽· εἶτα ποῦ τὰ τιτθία ; Eur.~\textit{Ba.}~925–927 τί φαίνομαι δῆτ᾽; οὐχὶ τὴν Ἰνοῦς στάσιν / ἢ τὴν Ἀγαύης ἑστάναι, μητρός γ᾽ ἐμῆς; / αὐτὰς ἐκείνας εἰσορᾶν δοκῶ σ᾽ ὁρῶν ; Aesch.~\textit{Sept.}~592 οὐ γὰρ δοκεῖν ἄριστος, ἀλλ᾽ εἶναι θέλει ; \textit{idem} Pl.~\textit{R.}~2.361b (κατ᾽ Αἰσχύλον οὐ δοκεῖν ἀλλ᾽ εἶναι ἀγαθὸν ἐθέλοντα), Plu.~\textit{Mor.}~186b (praem. Αἰσχύλου δὲ ποιήσαντος εἰς Ἀμφιάραον ; add. … πάντες εἰς Ἀριστείδην ἀπέβλεψαν) ; Thgn.~125–128 οὐδὲ γὰρ εἰδείης ἀνδρὸς νόον οὐδὲ γυναικός, / πρὶν πειρηθείης ὥσπερ ὑποζυγίου … πολλάκι γὰρ γνώμην ἐξαπατῶσ᾽ ἰδέαι ; \textit{idem} Tosi n°~2179 (v. 128 tantum) ; Pl.~\textit{R.}~2.365c (Simon.) οὐκοῦν, ἐπειδὴ τὸ δοκεῖν, ὡς δηλοῦσί μοι οἱ σοφοί, καὶ τὰν ἀλάθειαν βιᾶται καὶ κύριον εὐδαιμονίας ; Apollod.~3.13.8 κρύψασα ἐσθῆτι γυναικείᾳ ὡς παρθένον Λυκομήδει παρέθετο ; Hyp. Ar.~\textit{Eccl.} (Rutherford II, p.~510) αἱ γυναῖκες συνέθεντο πάντα μηχανήσασθαι εἰς τὸ δόξαι ἄνδρες εἶναι … αἱ δὲ μηχαναὶ τοῦ δόξαι αὐτὰς ἄνδρας εἶναι τοιαῦται· πώγωνας περιθέτους ; Sall.~\textit{Cat.}~54.6 esse quam videri bonus malebat: ita, quo minus petebat gloriam, eo magis illum adsequebatur ; Tosi n°~168 Meos tam suspicione quam crimine iudico carere}
\lp{186 ὑπεραποκρίνῃ μου]}{Ar.~\textit{Ach.}~632 ἀποκρίνασθαι δεῖται νυνὶ πρὸς Ἀθηναίους μεταβούλους ; Ar.~\textit{Vesp.}~950–951 χαλεπὸν μὲν ὦνδρες ἐστὶ διαβεβλημένου / ὑπεραποκρίνεσθαι κυνός, λέξω δ᾽ ὅμως ; Eup.~fr.~219~Kock ὡς ὑμῖν ἐγὼ / πάντως ἀποκρινοῦμαι πρὸς τὰ κατηγορούμενα ; Pl.~\textit{Grg.}~515c ὡμολογήκαμεν· ἐγὼ ὑπὲρ σοῦ ἀποκρινοῦμαι ; Hsch.~υ~667 ὑποκρίνοιτο· ἀποκρίνοιτο. ἔνθεν καὶ ὑποκριτής, ὁ ἀποκρινόμενος πρὸς τὸν χορόν ; Ar.~\textit{Thesm.}~362 ζητοῦσ᾽ ἀντιμεθιστάναι ; Ar.~\textit{Thesm.}~924 τουτὶ πονηρόν· ἀλλ᾽ ὑπαποκινητέον ; Harp. s.v.~ἀπόκρισις ἈΠΟΚΡΙΣΙΣ ἡ ἀπολογία· Λυσίας. καὶ Δημοσθένης κατ᾽ Αἰσχίνου ἀπόκριναι γὰρ δεῦρο παρελθών μοι ; Ar.~\textit{Thesm.}~951 ἐς τὰς ὥρας ξυνεπευχόμενος ; Hsch.~α~6411 ἀποκρίνου· ἀποφαίνου}
\lp{186 σαφῶς σώσεις ἐμέ]}{Ar.~\textit{Ach.}~400–401 ὦ τρισμακάρι᾽ Εὐριπίδη, / ὅθ᾽ ὁ δοῦλος οὑτωσὶ σοφῶς ὑποκρίνεται ; Ar.~\textit{Eq.}~1042 ἔφραζεν ὁ θεός σοι σαφῶς σῴζειν ἐμέ ; Ar.~\textit{Nub.}~1245 ἐγὼ γὰρ αὐτίκ᾽ ἀποκρινοῦμαί σοι σαφῶς ; Ar.~\textit{Ran.}~1434 ὁ μὲν σοφῶς γὰρ εἶπεν, ὁ δ᾽ ἕτερος σαφῶς ; Eur.~\textit{Hipp.}~346 οὐ μάντις εἰμὶ τἀφανῆ γνῶναι σαφῶς ; Eur.~\textit{IA}~1540 ἀλλ᾽, ὦ φίλη δέσποινα, πᾶν πεύσῃ σαφῶς ; Eur.~\textit{IT}~765 τὸ σῶμα σῴσας τοὺς λόγους σῴσεις ἐμοί ; Th.~1.74.1 ὃς αἰτιώτατος ἐν τῷ στενῷ ναυμαχῆσαι ἐγένετο, ὅπερ σαφέστατα ἔσωσε τὰ πράγματα ; Ar.~\textit{Vesp.}~53 οὕτως ὑποκρινόμενον σοφῶς ὀνείρατα ; Ar.~\textit{Eccl.}~1134 εὐδαιμονικόν γ᾽ ἄνθρωπον εἴρηκας σαφῶς ; Eur.~\textit{Hel.}~308 καὶ μὴν σαφῶς γ᾽ ἔλεξ᾽ ὀλωλέναι πόσιν ; Eur.~\textit{IA}~1608 ἡ παῖς σαφῶς σοι πρὸς θεοὺς ἀφίπτατο ; Aesch.~\textit{Ag.}~1636 τὸ γὰρ δολῶσαι πρὸς γυναικὸς ἦν σαφῶς ; Soph.~\textit{Phil.}~40 ἁνὴρ κατοικεῖ τούσδε τοὺς τόπους σαφῶς ; \textit{Il.}~1.80–83 κρείσσων γὰρ βασιλεὺς ὅτε χώσεται ἀνδρὶ χέρηϊ … σὺ δὲ φράσαι εἴ με σαώσεις ; Ar.~\textit{Plu.}~40 πεύσει. σαφῶς γὰρ ὁ θεὸς εἶπέ μοι τοδί ; Eur.~\textit{Heracl.}~872 θεοῖς ὁμιλεῖν νῦν ἐπίσταμαι σαφῶς ; Hsch.~σ~289 σάφα· σαφῶς ; Hsch.~σ~303 σαώσεις· σώσεις (Α 83) ; Hsch.~σ~1374 σοφῶς· εὐτάκτως, καλῶς}
\lp{186 σαφῶς σώσεις ἐμέ (τὰ σῖγμα τῶν Εὐριπίδου)]}{Pl. Com. fr.~30~Kock εὖ γέ σοι γένοιθ᾽, ὅτι / ἔσωσας ἐκ τῶν σῖγμα τῶν Εὐριπίδου ; Eub.~fr.~26~Kock Εὐριπίδου δ᾽ ἔσωσά σ᾽, ὡς ἴσασι⟨ν Ἑλλήνων ὅ⟩σοι … τὰ σῖγμα συλλέξαντες, ὡς αὐτοὶ σοφοί ; Eur.~\textit{Med.}~476 ἔσῳσά σ᾽, ὡς ἴσασιν Ἑλλήνων ὅσοι ; Ath.~10.455b–d Πίνδαρος δὲ πρὸς τὴν ἀσιγμοποιηθεῖσαν ᾠδήν … ἐποίησε· πρὶν μὲν εἷρπε σχοινοτένειά τ᾽ ἀοιδὰ καὶ τὸ σὰν κίβδηλον ἀνθρώποις ; Σ Eur.~\textit{Med.}~476 ἔσωσά σ᾽ ὡς ἴσασιν: πλεονάζει ὁ στίχος τῷ [σ]. ὅθεν καὶ Πλάτων ἐν ταῖς Ἑορταῖς φησιν· ἔσωσας ἐκ τῶν σῖγμα τῶν Εὐριπίδου ; D.H.~\textit{Comp.}~14 ἄχαρι δὲ καὶ ἀηδὲς τὸ ς καὶ πλεονάσαν σφόδρα λυπεῖ· θηριώδους γὰρ καὶ ἀλόγου μᾶλλον ἢ λογικῆς ἐφάπτεσθαι δοκεῖ φωνῆς ὁ συριγμός … εἰσὶ δ᾽ οἳ καὶ ἀσίγμους ὅλας ᾠδὰς ἐποίουν ; Ar.~\textit{Thesm.}~456 ἅτ᾽ ἐν ἀγρίοισι τοῖς λαχάνοις αὐτὸς τραφείς ; Hsch.~ψ~115 ψελλός· ὁ τὸ σίγμα παχύτερον λέγων}
\lp{187 μόνος γὰρ ἂν λέξειας ἀξίως ἐμοῦ]}{Ar.~\textit{Thesm.}~173–174 καὶ γὰρ ἐγὼ τοιοῦτος ἦν / ὢν τηλικοῦτος, ἡνίκ᾽ ἠρχόμην ποιεῖν ; Aesch.~\textit{Pers.}~631–632 εἰ γάρ τι κακῶν ἄκος οἶδε πλέον, / μόνος ἂν θνητῶν πέρας εἴποι ; Ar.~\textit{Ran.}~71–72 δέομαι ποιητοῦ δεξιοῦ. / οἱ μὲν γὰρ οὐκέτ᾽ εἰσίν, οἱ δ᾽ ὄντες κακοί ; Ar.~\textit{Ran.}~83–84 Ἀγάθων δὲ ποῦ ᾽στιν; — ἀπολιπών μ᾽ ἀποίχεται, / ἀγαθὸς ποιητὴς καὶ ποθεινὸς τοῖς φίλοις ; Ar.~\textit{Ran.}~954–958 ἔπειτα τουτουσὶ λαλεῖν ἐδίδαξα … / νοεῖν ὁρᾶν ξυνιέναι στρέφειν ἐρᾶν τεχνάζειν, / κάχ᾽ ὑποτοπεῖσθαι, περινοεῖν ἅπαντα ; Ar.~\textit{Ran.}~1413 τὸν μὲν γὰρ ἡγοῦμαι σοφὸν τῷ δ᾽ ἥδομαι ; Σ Ar.~\textit{Nub.}~144 σοφὸς Σοφοκλῆς, σοφώτερος δ᾽ Εὐριπίδης· / ἀνδρῶν δὲ πάντων Σωκράτης σοφώτατος ; \textit{idem} Suid.~σ~820 (σοφώτερος pro σοφώτατος) ; Ar.~\textit{Ach.}~428–429 ἀλλὰ κἀκεῖνος μὲν ἦν / χωλὸς προσαιτῶν στωμύλος δεινὸς λέγειν ; Ar.~\textit{Pax}~532–534 ἐπυλλίων Εὐριπίδου— … οὐ γὰρ ἥδεται / αὕτη ποιητῇ ῥηματίων δικανικῶν}
\lp{187 μόνος γὰρ ἄν]}{Eur.~\textit{Ba.}~962 μόνος γὰρ αὐτῶν εἰμ᾽ ἀνὴρ τολμῶν τόδε ; Eur.~\textit{El.}~285 εἷς ἂν μόνος νιν τῶν ἐμῶν γνοίη φίλων ; Eur.~\textit{Alc.}~122 μόνος δ᾽ ἄν, εἰ φῶς τόδ᾽ ἦν ; Eur.~\textit{Alc.}~293 μόνος γὰρ αὐτοῖς ἦσθα, κοὔτις ἐλπὶς ἦν}
\lp{187 λέξειας ἀξίως ἐμοῦ]}{Eur.~\textit{Hipp.}~345 πῶς ἂν σύ μοι λέξειας ἁμὲ χρὴ λέγειν ; \textit{idem} Ar.~\textit{Eq.}~16 ; Eur.~\textit{IA}~507 ὑπέθηκας ὀρθῶς τοὺς λόγους σοῦ τ᾽ ἀξίως ; Eur.~\textit{IA}~1225 ζῶσάν τε καὶ θάλλουσαν ἀξίως ἐμοῦ ; Ar.~\textit{Ran.}~391–393 πολλὰ μὲν γέλοιά μ᾽ εἰπεῖν, / πολλὰ δὲ σπουδαῖα, καὶ / τῆς σῆς ἑορτῆς ἀξίως ; Eur.~\textit{IA}~1407 εὖ γὰρ τόδ᾽ εἶπας ἀξίως τε πατρίδος ; D.~6.11 ταῦθ᾽ ἃ πάντες ἀεὶ γλίχονται λέγειν, ἀξίως δ᾽ οὐδεὶς εἰπεῖν δεδύνηται ; D.~60.1 ἀξίως εἰπεῖν τῶν τετελευτηκότων ; Σ Eur.~\textit{Hec.}~990 ὡς λίαν καλῶς καὶ ἀξίως σου λέγεις ; Ar.~\textit{Ach.}~677 οὐ γὰρ ἀξίως ἐκείνων ὧν ἐναυμαχήσαμεν ; Pl.~\textit{Ep.}~13.361a ἀξίως ἐμοῦ τε καὶ σοῦ ; Isoc.~10.13 οὐ γὰρ τῆς αὐτῆς γνώμης ἐστὶν ἀξίως εἰπεῖν περὶ ἑκατέρων αὐτῶν ; Hsch.~ε~2798 ἐνδίκως· δικαίως ; Hsch.~λ~649 λέξαι· εἰπεῖν … ἐπιλέξαι, ἐπειπεῖν}
\lp{188 ἔπειτα πῶς οὐκ]}{Ar.~\textit{Lys.}~24 κᾆτα πῶς οὐχ ἥκομεν ; Ephipp.~fr.~3.1–2~K.-A. (=~3~Kock ; Ath.~9.370c–371a) ἔπειτα πῶς οὐ στέφανος οὐδείς ἐστι πρόσθε τῶν θυρῶν ; Arist.~\textit{Rh.}~2.3.13 τί οὐκ ἀπολογεῖ ; Ar.~\textit{Av.}~278 εἶτα πῶς ἄνευ καμήλου Μῆδος ὤν εἰσέπτετο ; Ar.~\textit{Ran.}~647 κᾆτα πῶς οὐκ ἔπταρον ; Ar.~\textit{Eq.}~1207 τί οὐ διακρίνεις Δῆμ᾽ ὁπότερός ἐστι νῷν ; Alex.~fr.~125.1~K.-A. (=~120~Kock ; Ath.~3.114b–e) ἔπειτα πῶς ἦλθες ; Arist.~\textit{Pol.}~2.1268b ἔπειτα πῶς οὐκ ἔσται ταραχώδης ἡ κρίσις}
\lp{188 ἔπειτα]}{Ar.~\textit{Ach.}~917 ἔπειτα φαίνεις δῆτα διὰ θρυαλλίδα ; Ar.~\textit{Nub.}~226 ἔπειτ᾽ ἀπὸ ταρροῦ τοὺς θεοὺς ὑπερφρονεῖς ; Ar.~\textit{Nub.}~1249 ἔπειτ᾽ ἀπαιτεῖς τἀργύριον τοιοῦτος ὤν ; Antiph.~fr.~75.9~K.-A. (=~74.9~Kock ; Ath.~10.459a) ἔπειτα τοῦτο ζημιοῖς με; μηδαμῶς ; Ar.~\textit{Vesp.}~1133 ἔπειτα παῖδας χρὴ φυτεύειν καὶ τρέφειν ; Ar.~\textit{Lys.}~914 ἔπειτ᾽ ὀμόσασα δῆτ᾽ ἐπιορκήσω τάλαν ; Ar.~\textit{Thesm.}~637 κἄπειτ᾽ ἀποδύσετ᾽ ἐννέα παίδων μητέρα ; Ar.~\textit{Plu.}~827 μάλιστ᾽. — ἔπειτα τοῦ δέει ; Hsch.~ε~5289 ἔπειτα· εἶθ᾽ οὕτως. ἢ μετὰ ταῦτα}
\lp{188 αὐτὸς ἀπολογεῖ παρών]}{Ar.~\textit{Vesp.}~840 σὺ δὲ κατηγόρει παρών ; Hdt.~6.136.2 Μιλτιάδης δὲ αὐτὸς μὲν παρεὼν οὐκ ἀπελογέετο· ἦν γὰρ ἀδύνατος … ὑπεραπελογέοντο οἱ φίλοι ; Ar.~\textit{Vesp.}~944 ἀνάβαιν᾽, ἀπολογοῦ. τί σεσιώπηκας; λέγε ; Ar.~\textit{Vesp.}~949 πάρεχ᾽ ἐκποδών. ἐγὼ γὰρ ἀπολογήσομαι ; Ar.~\textit{Av.}~548 ἀλλ᾽ ὅ τι χρὴ δρᾶν, σὺ δίδασκε παρών ; Ar.~\textit{Plu.}~140–141 ἢν σὺ μὴ παρὼν / αὐτὸς διδῷς τἀργύριον ; Eur.~\textit{Heracl.}~561 πέπλοις δὲ σῶμ᾽ ἐμὸν κρύψον παρών ; Eur.~\textit{Hipp.}~1242 τίς ἄνδρ᾽ ἄριστον βούλεται σῶσαι παρών ; [Eur.]~\textit{Rh.}~179 καὶ μὴν λαφύρων γ᾽ αὐτὸς αἱρήσῃ παρών ; Soph.~\textit{Aj.}~1131 εἰ τοὺς θανόντας οὐκ ἐᾷς θάπτειν παρών ; Soph.~\textit{Aj.}~1156 τοιαῦτ᾽ ἄνολβον ἄνδρ᾽ ἐνουθέτει παρών ; Plu.~\textit{Alc.}~19.4 ἀπολογείσθω παρών ; Hsch.~α~6303 ἀποδικεῖν· ἀπολογεῖσθαι δίκην (Antiphan. fr. 303 K.–A.) ; Hsch.~α~6412 ἀπόκρισις· ἀπολογία (Lys. fr. 305 S. = 475 C. Antiph. or. fr. 154 Tha.) ; Ar.~\textit{Av.}~1215 ἐπέβαλεν ὀρνίθαρχος οὐδείς σοι παρών ; Ar.~\textit{Lys.}~284 ἐγὼ οὐκ ἄρα σχήσω παρὼν τολμήματος τοσούτου ; Eur.~\textit{Supp.}~391 στρατὸς δὲ θάσσει κἀξετάζεται παρὼν ; Aesch.~\textit{Cho.}~1014 νῦν αὐτὸν αἰνῶ, νῦν ἀποιμώζω παρών ; Soph.~\textit{Aj.}~742 πρὶν παρὼν αὐτὸς τύχῃ ; Th.~6.29.2 μὴ ἀπόντος πέρι αὐτοῦ διαβολὰς ἀποδέχεσθαι ; D.~58.52 τοῦτον ἀπολογεῖ ; Hsch.~α~6138 ἀπηγορήσομαι· ἀπολογήσομαι ; Hsch.~α~6457 ἀπολογία· πληροφορία ; Ammon.~\textit{Diff.}~62 ἀπολογεῖσθαι μὲν γάρ ἐστι τὸ ἐπὶ τῷ λόγῳ τὴν κατηγορίαν ἀνασκευάζειν, ἀπολογίζεσθαι δὲ τὸ ἀποδοῦναι τὸν λόγον τῶν δαπανημάτων}
\lp{188 αὐτὸς ἀπολογεῖ παρών (ἐρήμη δίκη)]}{Zen.~3.84 Ἐρήμη δίκη: ὅταν μὴ ἀπαντήσας ὁ διωκόμενος ἐπὶ τὴν κρίσιν καταδικασθῇ ; \textit{idem} Hsch.~ε~5773 (ἐρήμη· μονομερὴς δίκη ; καταδιαιτηθῇ) ; Th.~6.61.7 οἱ δ᾽ Ἀθηναῖοι ἐρήμῃ δίκῃ θάνατον κατέγνωσαν αὐτοῦ τε καὶ τῶν μετ᾽ ἐκείνου}
\lp{189 ἐγὼ φράσω σοι]}{Ar.~\textit{Ach.}~915 ἐγὼ φράσω σοι τῶν περιεστώτων χάριν ; Ar.~\textit{Eq.}~1211 ἐγὼ φράσω σοι. τὴν ἐμὴν κίστην ἰὼν ; Ar.~\textit{Ran.}~1177 σὺ δὲ πῶς ἐποίεις τοὺς προλόγους; — ἐγὼ φράσω ; Eur.~\textit{Hec.}~1180 ἅπαντα ταῦτα συντεμὼν ἐγὼ φράσω ; Philyll.~fr.~7~K.-A. (=~8~Kock ; Ath.~4.171d) βούλεσθε δῆτ᾽ ἐγὼ φράσω τίς εἰμ᾽ ἐγώ ; Eur.~\textit{Med.}~535 εἴληφας ἢ δέδωκας, ὡς ἐγὼ φράσω ; D.~4.20 καὶ πῶς ταῦτ᾽ ἐθελήσει ποιεῖν; ἐγὼ φράσω, καθ᾽ ἕκαστον τούτων διεξιὼν χωρίς ; Ar.~\textit{Nub.}~1354 ἐγὼ φράσω· ᾽πειδὴ γὰρ εἱστιώμεθ᾽, ὥσπερ ἴστε ; Antiph.~fr.~75.2~K.-A. (=~74~Kock ; Ath.~10.459a) ἀλλ᾽ ἐγὼ σαφῶς φράσω ; Eur.~\textit{Ion}~816 λάθρα πέφηνεν· ὡς λάθρα δ᾽, ἐγὼ φράσω ; Soph.~\textit{OC}~546 τί τοῦτο; — πρὸς δίκας τι. — τί γάρ; — ἐγὼ φράσω ; Soph.~\textit{Phil.}~1336 ὡς δ᾽ οἶδα ταῦτα τῇδ᾽ ἔχοντ᾽ ἐγὼ φράσω ; Hsch.~φ~844 φράσον· εἰπέ, λέξον (Aesch. Pers. 333). ἑρμήνευσον (Ev. Matth. 15, 15)}
\lp{189 γιγνώσκομαι]}{Ar.~\textit{Ach.}~441–443 εἶναι μὲν ὅσπερ εἰμί, φαίνεσθαι δὲ μή· / τοὺς μὲν θεατὰς εἰδέναι μ᾽ ὃς εἴμ᾽ ἐγώ, / τοὺς δ᾽ αὖ χορευτὰς ἠλιθίους παρεστάναι ; Ar.~\textit{Thesm.}~33–34 οὐχ ἑόρακας πώποτε; / μὰ τὸν Δί᾽ οὔτοι γ᾽ ὥστε καί μέ γ᾽ εἰδέναι ; Hsch.~α~1756 ὅταν θέλωσιν ἀλλήλοις μὴ γινώσκεσθαι ; Hsch.~γ~740 γνώριμος· οἰκεῖος. γινωσκόμενος (π 9)}
\lp{189–190 πρῶτα μὲν … ἔπειτα]}{Eur.~\textit{Med.}~548–549 ἐν τῷδε δείξω πρῶτα μὲν σοφὸς γεγώς, / ἔπειτα σώφρων, εἶτα σοὶ μέγας φίλος ; Ar.~\textit{Plu.}~771–772 καὶ προσκυνῶ γε πρῶτα μὲν τὸν ἥλιον, / ἔπειτα σεμνῆς Παλλάδος κλεινὸν πέδον ; Eur.~\textit{Med.}~720–721 γύναι, πρόθυμός εἰμι, πρῶτα μὲν θεῶν, / ἔπειτα παίδων ὧν ἐπαγγέλλῃ γονάς ; Ar.~\textit{Ach.}~648 ἠρώτησεν πρῶτα μὲν αὐτοὺς πότεροι ταῖς ναυσὶ κρατοῦσιν ; Ar.~\textit{Eq.}~774 ὃς πρῶτα μὲν ἡνίκ᾽ ἐβούλευον σοὶ χρήματα πλεῖστ᾽ ἀπέδειξα ; Ar.~\textit{Vesp.}~552 ὃν πρῶτα μὲν ἕρποντ᾽ ἐξ εὐνῆς τηροῦσ᾽ ἐπὶ τοῖσι δρυφάκτοις ; Ar.~\textit{Ran.}~1528 πρῶτα μὲν εὐοδίαν ἀγαθὴν ἀπιόντι ποιητῇ ; Ar.~\textit{Plu.}~728–729 καὶ πρῶτα μὲν δὴ τῆς κεφαλῆς ἐφήψατο, / ἔπειτα καθαρὸν ἡμιτύβιον λαβὼν}
\lp{190 πολιός εἰμι]}{Eur.~\textit{Alc.}~470 σχετλίω, πολιὰν ἔχοντε χαίταν ; Eur.~\textit{Ba.}~324 πολιὰ ξυνωρίς, ἀλλ᾽ ὅμως χορευτέον ; Σ Ar.~\textit{Thesm.}~190 γέρων γὰρ τότε Εὐριπίδης ἦν· ἕκτῳ γοῦν ἔτει ὕστερον τελευτᾷ ; Ar.~\textit{Vesp.}~1192 ἤδη γέρων ὢν καὶ πολιός, ἔχων δέ τοι ; Ar.~\textit{Lys.}~595 ὁ μὲν ἥκων γάρ, κἂν ᾖ πολιός, ταχὺ παῖδα κόρην γεγάμηκεν ; Eur.~\textit{Supp.}~166 πολιὸς ἀνὴρ τύραννος εὐδαίμων πάρος ; Hsch.~π~2773 πολιά· γῆρας ; Hsch.~π~2782 πολιόν· λευκόν (Ι 366 . .) ; Hsch.~π~2785 πολιός· γέρων ; Ar.~\textit{Ach.}~600 ὁρῶν πολιοὺς μὲν ἄνδρας ἐν ταῖς τάξεσιν ; Ar.~\textit{Ach.}~610 ἤδη πεπρέσβευκας σὺ πολιὸς ὢν ; Ar.~\textit{Ach.}~692 ταῦτα πῶς εἰκότα, γέροντ᾽ ἀπολέσαι πολιὸν ἄνδρα περὶ κλεψύδραν ; Eur.~\textit{Ba.}~185 καὶ κρᾶτα σεῖσαι πολιόν ; Eur.~\textit{Ba.}~258 εἰ μή σε γῆρας πολιὸν ἐξερρύετο ; Anacr.~fr.~43.1~Bergk⁴ (Stob. Flor.~118.13) Πολιοὶ μὲν ἡμὶν ἤδη κρόταφοι κάρη τε λευκόν ; Luc.~\textit{Merc.Cond.}~33 σκυθρωπῷ καὶ γέροντι ἀνδρὶ καὶ πολιῷ τὸ γένειον}
\lp{190 πολιός εἰμι καὶ πώγων᾽ ἔχω]}{Luc.~\textit{Salt.}~5 ἐν βαθεῖ τούτῳ τῷ πώγωνι καὶ πολιᾷ τῇ κόμῃ καθῆσθαι μέσον ἐν τοῖς γυναίοις ; \textit{Il.}~22.74 ἀλλ᾽ ὅτε δὴ πολιόν τε κάρη πολιόν τε γένειον ; Tyrt.~fr.~10.23~Bergk⁴ (Lycurg.~\textit{Leocr.}~107) ἤδη λευκὸν ἔχοντα κάρη πολιόν τε γένειον … / (νέοισι δὲ πάντ᾽ ἐπέοικεν …) / ἀνδράσι μὲν θηητὸς ἰδεῖν, ἐρατὸς δὲ γυναιξίν ; Pl.~\textit{Prm.}~127b τὸν μὲν οὖν Παρμενίδην εὖ μάλα ἤδη πρεσβύτην εἶναι, σφόδρα πολιόν, καλὸν δὲ κἀγαθὸν τὴν ὄψιν … Ζήνωνα δὲ … εὐμήκη δὲ καὶ χαρίεντα ἰδεῖν, καὶ λέγεσθαι αὐτὸν παιδικὰ τοῦ Παρμενίδου γεγονέναι ; Luc.~\textit{Merc.Cond.}~12 ἄνδρα ἐν βαθεῖ πώγωνι καὶ πολιᾷ τῇ κόμῃ ; Plu.~\textit{Ant.}~80.3 ὁ δὲ πώγωνα πολιὸν καθεὶς ; Luc.~\textit{Philops.}~23 μόνῃ τῇ πολιᾷ καὶ τῷ πώγωνι διαφέροντες τῶν βρεφῶν}
\lp{190 πολιός … πώγων᾽ ἔχω (ἀπὸ πώγωνος σοφοί)]}{Men.~fr.~639~Kock =~776~K.-A. (Stob. Flor.~52.10) οὐχ αἱ τρίχες ποιοῦσιν αἱ λευκαὶ φρονεῖν ; Apostol.~6.93e (CPG II, p.~390) Ἐκ πώγωνος σοφοὶ καὶ ἀπὸ πώγωνος σοφισταὶ καὶ ἀπὸ πώγωνος φιλόσοφοι: ἐπὶ τῶν δοκούντων μόνῳ τῷ σχήματι φιλόσοφοι εἶναι καὶ τῇ στολῇ ; Epict.~3.1.24 κατεκρίθην πώγωνα ἔχειν πολιὸν καὶ τρίβωνα ; Luc.~\textit{Demon.}~13 εἰ σὺ ἀπὸ τοῦ πώγωνος ἀξιοῖς κρίνεσθαι τοὺς φιλοσοφοῦντας ; Tosi n°~1093 Non aetate verum ingenio apiscitur sapientia ; Pl.~\textit{Tht.}~168e μεγάλους πώγωνας ἐχόντων ; Plu.~\textit{De\,Is.\,et\,Os.}~3 (352c) οὔτε γὰρ φιλοσόφους πωγωνοτροφίαι, ὦ Κλέα, καὶ τριβωνοφορίαι ποιοῦσιν ; Tosi n°~2183 Barba non facit philosophum ; ἀπὸ πώγωνος φιλόσοφοι / σοφοί / σοφισταί}
\lp{190 πώγων᾽ ἔχω]}{Ar.~\textit{Eccl.}~99–103 τὸν πώγωνά τε / ὅταν καθῶμεν ὃν περιδησόμεσθ᾽ ἐκεῖ, / τίς οὐκ ἂν ἡμᾶς ἄνδρας ἡγήσαιθ᾽ ὁρῶν; / Ἀγύρριος γοῦν τὸν Προνόμου πώγων᾽ ἔχων / λέληθε· καίτοι πρότερον ἦν οὗτος γυνή ; Alex.~fr.~266.6–7~K.-A. (=~264.5–6~Kock ; Ath.~13.565b) τί γὰρ αἱ τρίχες λυποῦσιν ἡμᾶς, πρὸς θεῶν; / δι᾽ ἃς ἀνὴρ ἕκαστος ἡμῶν φαίνεται ; Vit. Eur. (Schwartz I, p.~2,~10–11) ἐλέγετο δὲ καὶ βαθὺν πώγωνα θρέψαι ; Apostol.~1.19 (CPG II, p.~244) Ἀγύριός ἐστι Προνόμου πώγων᾽ ἔχων: ἐπὶ τῶν ψηνιζομένων καὶ ὀπυομένων ; App. Prov.~2.36 (CPG I, p.~400) Εἰς Τροιζῆνα δεῖ σε βαδίζειν: ἐπὶ τῶν κακογενείων καὶ σπανοπωγώνων εἰρηται. Πώγων γάρ ἐστι λιμὴν Τροιζήνιος ; Ar.~\textit{Eccl.}~68–71 ἔχετε δὲ τοὺς πώγωνας, οὓς εἴρητ᾽ ἔχειν / πάσαισιν ἡμῖν, ὁπότε συλλεγοίμεθα; / — νὴ τὴν Ἑκάτην καλόν γ᾽ ἔγωγε τουτονί. / — κἄγωγ᾽ Ἐπικράτους οὐκ ὀλίγῳ καλλίονα ; \textit{idem} Ar.~\textit{Eccl.}~25 (ἔχουσι pro ἔχετε δέ) ; Plu.~\textit{Mor.}~245c–f (Mul. virt.~4) ἔθεντο νόμον τὸν κελεύοντα πώγωνα δεῖν ἐχούσας συναναπαύεσθαι τοῖς ἀνδράσι τὰς γεγαμημένας ; Luc.~\textit{Eun.}~8 καὶ τὸ μέγιστον, πώγωνα βαθὺν ἔχειν αὐτὸν ; Σ Ar.~\textit{Eccl.}~102 ὁ δὲ Πρόνομος αὐλητὴς μέγαν ἔχων πώγωνα, Ἀγύρριος δὲ εὐρύπρωκτος. οὗτος οὖν τὸν Προνόμου ἔχων πώγωνα λανθάνει ὢν γυνή ; Suid.~π~2150 κἀγώ γ᾽ Ἐπικράτους οὐκ ὀλίγῳ καλλίονα πώγων᾽ ἔχω ; Erasm.~\textit{Adag.}~4.2.17 (LB II, col.~994 E–F) Προνόμου πώγων. i. e. Pronomi barba, de praelonga dicebatur. Aristophanes in Concionatricibus … Argyrrhius exercitus dux quispiam fuit effoeminatus \& cinaedus, nihil habens viri praeter barbam. … Convenit igitur in quosdam, qui alioqui mollissimi, vultus vestiumque severitate, vitae continentiam praeferunt ; Hsch.~π~4488 πώγων· γένειον. ἄλλοι τὴν ἀναφορὰν τῆς φλογός ; Suid.~π~2151 Πωγωνίας: ὁ μέγας πώγων ; Apostol.~3.53 (CPG II, p.~299) Ἄπλυτον πώγων᾽ ἔχεις: σκώπτων εἶπε}
\lp{191 σὺ δ᾽ εὐπρόσωπος … ἐξυρημένος]}{Ar.~\textit{Thesm.}~31–33 μῶν ὁ μέλας ὁ καρτερός; / οὔκ, ἀλλ᾽ ἕτερός τις· οὐχ ἑόρακας πώποτε; / μῶν ὁ δασυπώγων}
\lp{191 σὺ δ᾽ εὐπρόσωπος (ὁ καλὸς Ἀγάθων)]}{Pl.~\textit{Prt.}~315d–e νέον τι ἔτι μειράκιον, ὡς μὲν ἐγᾦμαι καλόν τε κἀγαθὸν τὴν φύσιν, τὴν δ᾽ οὖν ἰδέαν πάνυ καλός. ἔδοξα ἀκοῦσαι ὄνομα αὐτῷ εἶναι Ἀγάθωνα ; Luc.~\textit{Rh.Pr.}~11 πάναβρόν τινα Σαρδανάπαλλον ἢ Κινύραν ἢ αὐτὸν Ἀγάθωνα, τὸν τῆς τραγῳδίας ἐπέραστον ἐκεῖνον ποιητήν ; Pl.~\textit{Smp.}~174a ταῦτα δὴ ἐκαλλωπισάμην, ἵνα καλὸς παρὰ καλὸν ἴω ; Pl.~\textit{Smp.}~212e ἵνα ἀπὸ τῆς ἐμῆς κεφαλῆς τὴν τοῦ σοφωτάτου καὶ καλλίστου κεφαλὴν ἐὰν εἴπω οὑτωσὶ ἀναδήσω ; Pl.~\textit{Smp.}~213c ἀλλὰ διεμηχανήσω ὅπως παρὰ τῷ καλλίστῳ τῶν ἔνδον κατακείσῃ ; Plu.~\textit{Mor.}~177a–b τοῦ δὲ Εὐριπίδου τὸν καλὸν Ἀγάθωνα περιλαμβάνοντος ἐν τῷ συμποσίῳ καὶ καταφιλοῦντος ἤδη γενειῶντα ; Plu.~\textit{Mor.}~645d–e (\textit{Quaest.\,conv.}~3.1.1) κατηγοροῦντα τοῦ καλοῦ Ἀγάθωνος, ὃν πρῶτον εἰς τραγῳδίαν φασὶν ἐμβαλεῖν καὶ ὑπομῖξαι τὸ χρωματικόν ; Ael.~\textit{VH}~13.4 συγκλιθέντα αὐτῷ Ἀγάθωνα τὸν τῆς τραγῳδίας ποιητὴν περιλαβὼν κατεφίλει, τετταράκοντα ἐτῶν που γεγονότα ; Ath.~5.185a κατὰ τὸν καλὸν Ἀγάθωνα ; Erasm.~\textit{Adag.}~1.3.72 (LB II, col.~139 D–E) Τῶν καλῶν καὶ ὀπώρη καλή, i. e. Pulchrorum etiam autumnus pulcher. Metaphora proverbialis nata ex Archelai apophthegmate, quod ab eo dictum Plutarchus refert in Euripidem, qui jam pubescentem atque exoletum Agathonem in convivio suaviabatur ; Com. adesp. fr.~216~Kock τῶν γὰρ καλῶν τοι καὶ τὸ μετόπωρον καλόν ; Plu.~\textit{Mor.}~770c ἔφη γὰρ Ἀγάθωνα τὸν καλὸν ἢδη γενειῶντα περιβάλλων καὶ κατασπαζόμενος}
\lp{191 εὐπρόσωπος]}{Ar.~\textit{Plu.}~976 πενιχρὸν μέν, ἄλλως δ᾽ εὐπρόσωπον καὶ καλὸν ; Ar.~\textit{Pax}~617 ταῦτ᾽ ἄρ᾽ εὐπρόσωπος ἦν ; Ar.~\textit{Ran.}~412 νῦν δὴ κατεῖδον καὶ μάλ᾽ εὐπροσώπου ; Anaxandr.~fr.~9.5~K.-A. (=~9~Kock ; Ath.~13.570e) Θεολύτη μάλ᾽ εὐπρόσωπος καὶ καλή ; Aesch.~\textit{Sept.}~533 βλάστημα καλλίπρῳρον, ἀνδρόπαις ἀνήρ ; Pl.~\textit{Chrm.}~154d δόξει σοι ἀπρόσωπος εἶναι ; X.~\textit{Mem.}~1.3.10 τὸν Ἀλκιβιάδου υἱὸν φιλῆσαι, ὄντα εὐπροσωπότατον καὶ ὡραιότατον ; Hsch.~α~6858 ἀπρόσωπος· οὐκ εὐπρόσωπος (Plat. Charm. 154d) ; Hsch.~κ~484 καλλίπρωρον· εὐπρόσωπον (Aesch. Sept. 533) ; Cratin.~fr.~337~K.-A. (=~304~Kock ; Eust. in Od. p.~1430,23) εὐπρόσωπος ἦσθ᾽ ὀνόματος οὐδὲν ἐπὶ χεῖρας φέρων ; Eur.~\textit{Phoen.}~1336 οἰχόμεσθ᾽· οὐκ εὐπροσώποις φροιμίοις ἄρχῃ λόγου ; Aesch.~\textit{Cho.}~969–970 τύχαι δ᾽ εὐπροσωποκοῖται τὸ πᾶν / ἰδεῖν ; Soph.~\textit{Aj.}~1009 δέξαιτ᾽ ἂν εὐπρόσωπος ἵλεώς τ᾽ ἴσως ; Plu.~\textit{Mor.}~458c–459b (De cohib. ira~10) καὶ μὴν ἐδόκουν εὐπρόσωπος εἶναι ; Plu.~\textit{Mor.}~632b (\textit{Quaest.\,conv.}~2.1.5) διὸ καὶ Κριτόβουλον ὁ Σωκράτης εὐπροσωπότατον ὄντα προκαλούμενος εἰς σύγκρισιν εὐμορφίας ἔπαιζεν οὐκ ἐχλεύαζεν ; Luc.~\textit{Apol.}~3 τὴν ἀπολογίαν ἥτις ἂν εὐπρόσωπός σοι γένοιτο ; Σ Eur.~\textit{Or.}~1093–1096 τίνα εὐπρόσωπον ἀπολογίαν δώσω ; Σ Eur.~\textit{Phoen.}~1336 εὐπρόσωπον μειράκιον, ἤγουν εὐειδὲς ; Hsch.~ε~7090 εὔπρῳρον· εὐπρόσωπον. εὐκέφαλον ; Hsch.~ε~7215 εὔστομον· ἀπὸ μέρους, εὐπρόσωπον ; Hsch.~ω~296 ὡραῖος· εὔμορφος, εὔτροπος, εὐσχήμων, εὐπρόσωπος ; Ammon.~\textit{Diff.}~201 εὔμορφος μὲν γὰρ ὁ τὴν μορφὴν εὖ ἔχων οἷον εὐπρόσωπος ; Suid.~α~4235 Ἀστεῖος: εὐσύνετος, εὐπρόσωπος, χαρίεις, καλὸς, γελοιώδης}
\lp{191 λευκός]}{Ar.~\textit{Av.}~668 ὡς δ᾽ ἁπαλόν, ὡς δὲ λευκόν ; Ar.~\textit{Eccl.}~427–428 εὐπρεπὴς νεανίας / λευκός τις ἀνεπήδησ᾽ ὅμοιος Νικίᾳ ; Cratin.~fr.~195.3~K.-A. (Ath.~1.29d) οἴμ᾽ ὡς ἁπαλὸς καὶ λευκός· ἆρ᾽ οἴσει τρία ; Com. adesp. fr.~791~Kock οὐδὲν λευκῶν ἀνδρῶν ὄφελος ; \textit{idem} Suid.~ο~801 (add. ἢ σκυτοτομεῖν· ἐπὶ τῶν εἰς μηδὲν λυσιτελούντων …), App. Prov.~4.35 (ἀνδρῶν λευκῶν ; add. ἢ σκυτοτομεῖν· ἐπὶ τῶν εἰς μηδὲν λυσιτελούντων …) ; Eur.~\textit{Ba.}~457–459 λευκὴν δὲ χροιὰν ἐκ παρασκευῆς ἔχεις, / οὐχ ἡλίου βολαῖσιν, ἀλλ᾽ ὑπὸ σκιᾶς, / τὴν Ἀφροδίτην καλλονῇ θηρώμενος ; Pl.~\textit{Phdr.}~239c–d μαλθακόν τινα καὶ οὐ στερεὸν διώκων, οὐδ᾽ ἐν ἡλίῳ καθαρῷ τεθραμμένον ἀλλὰ ὑπὸ συμμιγεῖ σκιᾷ, πόνων μὲν ἀνδρείων καὶ ἱδρώτων ξηρῶν ἄπειρον, ἔμπειρον δὲ ἁπαλῆς καὶ ἀνάνδρου διαίτης ; Pl.~\textit{R.}~474e μέλανας δὲ ἀνδρικοὺς ἰδεῖν, λευκοὺς δὲ θεῶν παῖδας εἶναι ; X.~\textit{An.}~5.4.32 παῖδας τῶν εὐδαιμόνων σιτευτούς … ἁπαλοὺς καὶ λευκοὺς σφόδρα ; Plu.~\textit{Ages.}~9.5 τῶν δὲ σωμάτων λευκῶν καὶ ἁπαλῶν παντάπασι διὰ τὰς σκιατραφίας γυμνουμένων κατεγέλων ; Plu.~\textit{Thes.}~23.2 δύο θηλυφανεῖς μὲν ὀφθῆναι καὶ νεαρούς … λουτροῖς τε θερμοῖς καὶ σκιατραφίαις καὶ ταῖς περὶ κόμην καὶ λειότητα καὶ χροιὰν ἀλοιφαῖς ; Hsch.~λ~734 λευκοί· οἱ δειλοί ; Eust.~455.37 (ad Δ~141), citant Ael. Dion. ὅτι τὸ ἐν σώμασι λευκὸν οὐκ ἐν ἐπαίνῳ ἦν, καὶ ἐν ῥητορικοῖς δηλοῦται λεξικοῖς. ὧν ἐν θατέρῳ μὲν κεῖται, ὅτι λευκοὶ οἱ δειλοί ; Ar.~\textit{Ran.}~1092 λευκὸς πίων ὑπολειπόμενος ; Ar.~\textit{Eccl.}~387 ὡς λευκοπληθὴς ἦν ἰδεῖν ἡκκλησία ; Ar.~\textit{Eccl.}~699 καὶ καλλίστη καὶ λευκοτάτη ; Sosicr.~fr.~1.1–2~K.-A. (Poll.~9.57) ὅταν γάρ, οἶμαι, λευκὸς ἄνθρωπος, παχύς, / ἀργός, λάβῃ δίκελλαν, εἰωθὼς τρυφᾶν ; Callias Com. fr.~14~K.-A. (Πεδῆται ; =~11~Kock ; Σ Ar. Av.~151) (A.) τί δ᾽ ἄρα; τοὺς Μελανθίου τῷ γνώσομαι; / (B.) οὓς ἂν μάλιστα λευκοπρώκτους εἰσίδῃς ; Pl.~\textit{R.}~556d ἰσχνὸς ἀνὴρ πένης, ἡλιωμένος, παραταχθεὶς ἐν μάχῃ πλουσίῳ ἐσκιατροφηκότι, πολλὰς ἔχοντι σάρκας ἀλλοτρίας ; X.~\textit{Hell.}~3.4.19 λευκοὺς μὲν διὰ τὸ μηδέποτε ἐκδύεσθαι, μαλακοὺς δὲ καὶ ἀπόνους διὰ τὸ ἀεὶ ἐπ᾽ ὀχημάτων εἶναι, ἐνόμισαν οὐδὲν διοίσειν τὸν πόλεμον ἢ εἰ γυναιξὶ δέοι μάχεσθαι ; X.~\textit{Oec.}~4.2 ἀναγκάζουσαι καθῆσθαι καὶ σκιατραφεῖσθαι ; [Arist.]~\textit{Phgn.}~812a Οἱ δὲ λευκοὶ ἄγαν δειλοί· ἀναφέρεται ἐπὶ τὰς γυναῖκας ; Plu.~\textit{Aem.}~31.7 τῷ σώματι στίλβων ὑπὸ λειότητος καὶ σκιατραφίας ; Plu.~\textit{Caes.}~17.2 καὶ τὴν σάρκα λευκὸς καὶ ἁπαλός ; Plu.~\textit{Lyc.}~14.2 ἀφελὼν δὲ θρύψιν καὶ σκιατραφίαν καὶ θηλύτητα πᾶσαν ; Suid.~α~2158 ἀνδράριόν τι θαλαμηπόλον, σκιατραφές τε καὶ ἁβροδίαιτον, καὶ πόρρω τοῦ ἀρρενωποῦ τεταγμένον ; App. Prov.~3.61 (CPG I, p.~429) Λευκοὶ τοὺς δειλοὺς ὀνειδίζουσιν ; Ar.~\textit{Nub.}~1012,~1017 στῆθος λιπαρόν, χροιὰν λαμπράν ; Alex.~fr.~217.1~K.-A. (=~214~Kock ; Ath.~9.395a) λευκὸς Ἀφροδίτης εἰμὶ γὰρ περιστερός ; Eur.~\textit{Alc.}~159–160 ὕδασι ποταμίοις λευκὸν χρόα / ἐλούσατ᾽ ; Eur.~\textit{Hipp.}~768–770 κρεμαστὸν ἅψεται ἀμφὶ βρόχον / λευκᾷ καθαρμόζουσα δεί-/ρᾳ ; Soph.~\textit{Ant.}~1238–1239 καὶ φυσιῶν ὀξεῖαν ἐκβάλλει ῥοὴν / λευκῇ παρειᾷ φοινίου σταλάγματος ; \textit{Il.}~1.55 θεὰ λευκώλενος Ἥρη ; Luc.~\textit{DMar.}~1.3 ἰδὲ σεαυτὴν οὐδὲν ἄλλο ἢ χροίαν λευκὴν ἀκριβῶς}
\lp{191 λευκός (λευκόπυγος, μελάμπυγος)]}{Suet.~\textit{Blasph.}~53 Taillardat (Miller, Mélanges, p.~415) Λευκόπυγος δ᾽ ὁ ἄνανδρος· καὶ ἔμπαλιν μελάμπυγος, ὁ ἀνδρεῖος ; \textit{idem} Hsch.~λ~739 (λευκόπυγος· ὁ ἄνανδρος (Alex. fr. 322 K.-A.). ἔμπαλιν δὲ μελαμπύγους τοὺς ἀνδρείους ἔλεγον), Suid.~λ~335 (λευκοπύγους· δειλούς, ὡς μελαμπύγους τοὺς ἀνδρείους) ; Suid.~μ~449 ὅτι τοὺς λευκοπύγους ὡς γυναικώδεις ἐκωμῴδουν ; App. Prov.~3.62 (CPG I, p.~429) Λευκόπυγος: ἐπὶ τῶν δειλῶν καὶ ἀνάνδρων. Τοιοῦτοι γὰρ οἱ μὴ πονοῦντες ἐν ἔργοις. Τούτῳ δὲ ἐναντίον τὸ μελάμπυγος, ἐπὶ τῶν γενναίων ; Hsch.~μ~641 μελάμπυγος· ἀνδρεῖος. τοὺς γὰρ δασεῖς τὰς πυγὰς ἀνδρείους ἐνόμιζον (Ar. Lys. 802) ; Zen.~5.10 Μὴ σύ γε μελαμπύγου τύχοις: μή τινος ἀνδρείου καὶ ἰσχυροῦ τύχοις}
\lp{191 ἐξυρημένος]}{Ar.~\textit{Ach.}~119–121 ὦ θερμόβουλον πρωκτὸν ἐξυρημένε, / τοιόνδε γ᾽ ὦ πίθηκε τὸν πώγων᾽ ἔχων / εὐνοῦχος ἡμῖν ἦλθες ἐσκευασμένος ; Ar.~\textit{Eq.}~1373–1374 οὐδ᾽ ἀγοράσει γ᾽ ἀγένειος οὐδεὶς ἐν ἀγορᾷ. / ποῦ δῆτα Κλεισθένης ἀγοράσει καὶ Στράτων ; Ar.~\textit{Thesm.}~141–143 τίς δ᾽ αὐτὸς ὦ παῖ; πότερον ὡς ἀνὴρ τρέφει; / καὶ ποῦ πέος; ποῦ χλαῖνα; ποῦ Λακωνικαί; / ἀλλ᾽ ὡς γυνὴ δῆτ᾽· εἶτα ποῦ τὰ τιτθία ; Ar.~\textit{Thesm.}~215–216 ἀποξυρεῖν ταδί, / τὰ κάτω δ᾽ ἀφεύειν ; Ar.~\textit{Thesm.}~218 Ἀγάθων σὺ μέντοι ξυροφορεῖς ἑκάστοτε ; Ar.~\textit{Thesm.}~574–575 φίλαι γυναῖκες ξυγγενεῖς τοὐμοῦ τρόπου· / ὅτι μὲν φίλος εἴμ᾽ ὑμῖν, ἐπίδηλος ταῖς γνάθοις ; Ar.~\textit{Thesm.}~590–591 ἀφηῦσεν αὐτὸν κἀπέτιλ᾽ Εὐριπίδης / καὶ τἄλλ᾽ ἅπανθ᾽ ὥσπερ γυναῖκ᾽ ἐσκεύασεν ; Ar.~\textit{Ran.}~426–428 τὸν Κλεισθένους δ᾽ ἀκούω / ἐν ταῖς ταφαῖσι πρωκτὸν / τίλλειν ἑαυτοῦ καὶ σπαράττειν τὰς γνάθους ; Ar.~fr.~341~K.-A. (Synag.~α~1520) ἀντίθετον· τὸ σχῆμα τῆς φράσεως, ὡς ἡμεῖς, Ἀριστοφάνης Θεσμοφοριαζούσαις (fr. 341)· ῾καὶ καταγαθὸν ἀντίθετον ἐξευρημένον.᾽ ; Ath.~13.564f–565a (Chrysippe) τὸ ξύρεσθαι τὸν πώγωνα κατ᾽ Ἀλέξανδρον προῆκται, τῶν πρώτων οὐ χρωμένων αὐτῷ ; Ath.~13.565c (Diogène) μή τι ἔχεις ἐγκαλεῖν τῇ φύσει ὅτι ἄνδρα σὲ ἐποίησε καὶ οὐ γυναῖκα ; Hsch.~σ~15 σάβυττος· εἶδος ξυρήσεως εἰς καλλωπισμόν· πότερον δὲ τοῦ πώγωνος ἢ τῆς κεφαλῆς ἄδηλον. τινὲς δὲ τὸ γυναικεῖον (Eupol. fr. 313,2 K.–A.) ; Suid.~ε~2021 ὁ δὲ Κλεισθένης ἐπ᾽ αἰσχρότητι κωμῳδεῖται. ἦν δὲ λεῖος τὸ γένειον καὶ εὐνούχῳ εἰκάζειν αὐτόν. τὸν δὲ Κλεισθένην καὶ ἐπὶ τοῦ πασχητιᾶν κωμῳδοῦσι ; Suid.~κ~1756 οὗτος ἀεὶ ἐξυρᾶτο. λέγεται οὖν ἐπὶ τῶν ἀεὶ ξυρωμένων. ἦν δὲ Σιβυρτίου· καὶ τοῦτο ἔπραττε διὰ τὸ νέος φαίνεσθαι ἀεί ; Ar.~\textit{Nub.}~355 καὶ νῦν γ᾽ ὅτι Κλεισθένη εἶδον, ὁρᾷς, διὰ τοῦτ᾽ ἐγένοντο γυναῖκες ; Ar.~\textit{Thesm.}~219–220 χρῆσόν τί νυν ἡμῖν ξυρόν. — αὐτὸς λάμβανε / ἐντεῦθεν ἐκ τῆς ξυροδόκης ; Ar.~\textit{Thesm.}~235 ὁρᾷς σεαυτόν; — οὐ μὰ Δί᾽ ἀλλὰ Κλεισθένη ; Ar.~\textit{Thesm.}~1043–1044 ὃς ἔμ᾽ ἀπεξύρησε πρῶτον, / ὃς ἐμὲ κροκόεντ᾽ ἐνέδυσεν ; D.Chr.~36.17 εἷς δὲ ἐν αὐτοῖς μόνος ἐξυρημένος, καὶ τοῦτον ἐλοιδόρουν τε καὶ ἐμίσουν ἅπαντες ; Ath.~13.565f ἐρωμένους περιάγεσθαι ξυρουμένους τὴν ὑπήνην καὶ τὸν ὄρρον ; Plb.~8.9.9 ὧν οἱ μὲν ξυρόμενοι καὶ λεαινόμενοι διετέλουν ἄνδρες ὄντες, οἱ δ᾽ ἀλλήλοις ἐτόλμων ἐπανίστασθαι πώγωνας ἔχουσι ; \textit{idem} Theopomp. Hist. ap. Ath.~6.260e (οὐχ pro ὧν ; ξυρούμενοι) ; Suid.~ε~1871 Ἐξυρημένος: τὴν κόμην ἀποβαλών ; Ar.~\textit{Eccl.}~65 κἄγωγε· τὸ ξυρὸν δέ γ᾽ ἐκ τῆς οἰκίας ; Soph.~\textit{Aj.}~786 ξυρεῖ γὰρ ἐν χρῷ τοῦτο μὴ χαίρειν τινά ; Hdt.~2.36.1 Αἰγύπτιοι δὲ ὑπὸ τοὺς θανάτους ἀνιεῖσι τὰς τρίχας αὔξεσθαι τάς τε ἐν τῇ κεφαλῇ καὶ τῷ γενείῳ, τέως ἐξυρημένοι ; Hsch.~π~1420 πέξις· ψίλωσις, κουρά, ξύρησις ; Hsch.~π~3280 πριαμωθήσομαι· ξυρήσομαι. ἐπειδὴ τὸ τραγικὸν τοῦ Πριάμου πρόσωπον ξυρίας ἐστίν (com. ad. fr. 414 K.–A.)}
\lp{192 γυναικόφωνος]}{\textit{Test.} Poll.~2.111 βαρύφωνος, ὁμόφωνος, γυναικόφωνος ὡς Ἀριστοφάνης ; Ar.~\textit{Nub.}~979–980 οὐδ᾽ ἂν μαλακὴν φυρασάμενος τὴν φωνὴν πρὸς τὸν ἐραστὴν / αὐτὸς ἑαυτὸν προαγωγεύων τοῖς ὀφθαλμοῖς ἐβάδιζεν ; Ar.~\textit{Thesm.}~131 καὶ θηλυδριῶδες καὶ κατεγλωττισμένον ; Ar.~\textit{Thesm.}~136 ποδαπὸς ὁ γύννις ; Ar.~\textit{Thesm.}~267–268 ἢν λαλῇς δ᾽, ὅπως τῷ φθέγματι / γυναικιεῖς εὖ καὶ πιθανῶς ; [Arist.]~\textit{Phgn.}~2.15 τὰ μὲν ἀνδρεῖα βαρύφωνά ἐστι, τὰ δὲ δειλὰ ὀξύφωνα ; Plu.~\textit{Caes.}~10.1–3 καὶ τῇ φωνῇ γενομένου καταφανοῦς ; Plu.~\textit{Cic.}~28.2 μειράκιον ὢν ἔτι καὶ μήπω γενειῶν ὁ Κλώδιος ἤλπιζε λήσεσθαι διαδὺς πρὸς τὴν Πομπηίαν μετὰ τῶν γυναικῶν … συνεῖσα τὴν φωνὴν οὐ γυναικείαν οὖσαν ἀνέκραγε ; Plu.~\textit{Thes.}~23.2 διδάξαντα φωνὴν καὶ σχῆμα καὶ βάδισιν ὡς ἔνι μάλιστα παρθένοις ὁμοιοῦσθαι ; D.Chr.~33.38 ὥστε γυναικῶν λαβεῖν φωνὴν ἅπαντας καὶ μηδένα δύνασθαι μήτε νέον μήτε πρεσβύτερον ἀνδριστὶ μηδὲν εἰπεῖν ; Luc.~\textit{Eun.}~10 οὑτοσὶ ὁ τὰς γνάθους λεῖος καὶ τὸ φώνημα γυναικεῖος ; Luc.~\textit{Rh.Pr.}~11 διασεσαλευμένον τὸ βάδισμα, ἐπικεκλασμένον τὸν αὐχένα, γυναικεῖον τὸ βλέμμα, μελιχρὸν τὸ φώνημα, μύρων ἀποπνέοντα ; Hsch.~γ~1015 γύννις· δειλός, ἄνανδρος, γυναικώδης· Ἀριστοφάνης (Thesm. 136) ; Hsch.~θ~491 θηλυδριῶδες· τὸ κατακεκλασμένον (Ar. Thesm. 131) ; Ar.~\textit{Thesm.}~97–98 ἐγὼ γὰρ οὐχ ὁρῶ / ἄνδρ᾽ οὐδέν᾽ ἐνθάδ᾽ ὄντα, Κυρήνην δ᾽ ὁρῶ ; Ar.~\textit{Thesm.}~576 γυναικομανῶ γὰρ προξενῶ θ᾽ ὑμῶν ἀεί ; Eur.~\textit{Ba.}~855 γυναικόμορφον ἀγόμενον δι᾽ ἄστεως ; Aesch.~\textit{Ag.}~487 γυναικογήρυτον ὄλλυται κλέος ; Arist.~\textit{HA}~4.11 πάντα τὰ θήλεα λεπτοφωνότερα καὶ ὀξυφωνότερα ; Arist.~\textit{HA}~7.1 Ὅλως μὲν γὰρ γυνὴ ἀνδρὸς ὀξυφωνότερον ; Arist.~\textit{Pr.}~10.36 καὶ γὰρ φωνὴν θηλυκὴν ἴσχουσιν καὶ ὀξύτητα καὶ ἀναρθρίαν ; Arist.~\textit{Pr.}~11.16 οἱ μὲν οὖν ἄγονοι διὰ ταῦτα ὀξύφωνοί εἰσιν ; Ael.~\textit{NA}~6.20 φθέγγεται δὲ αὐτῶν τὰ μὲν γοερὰ καὶ θηλύφωνα ; Ael.~\textit{VH}~12.1 φώνημα δὲ εἶχεν ἡδὺ καὶ ἁπαλόν ; Poll.~4.114 (Bethe I, p.~234) εἴποις δ᾽ ἂν βαρύστονος ὑποκριτής … βαρύφωνος δὲ καὶ λεπτόφωνος καὶ γυναικόφωνος καὶ στρηνόφωνος ; Eur.~\textit{Ba.}~980 ἐπὶ τὸν ἐν γυναικομίμῳ στολᾷ}
\lp{192 ἁπαλός]}{Pl.~\textit{Smp.}~195c νέος μὲν οὖν ἐστι, πρὸς δὲ τῷ νέῳ ἁπαλός· ποιητοῦ δ᾽ ἔστιν ἐνδεὴς οἷος ἦν Ὅμηρος πρὸς τὸ ἐπιδεῖξαι θεοῦ ἁπαλότητα ; Ar.~\textit{Av.}~667–668 ὦ Ζεῦ πολυτίμηθ᾽ ὡς καλὸν τοὐρνίθιον, / ὡς δ᾽ ἁπαλόν, ὡς δὲ λευκόν ; \textit{AP}~12.161.1–2 (Asclep.) Δόρκιον ἡ φιλέφηβος ἐπίσταται, ὡς ἁπαλὸς παῖς, / ἕσθαι πανδήμου Κύπριδος ὠκὺ βέλος ; Pl.~\textit{Smp.}~196a νεώτατος μὲν δή ἐστι καὶ ἁπαλώτατος, πρὸς δὲ τούτοις ὑγρὸς τὸ εἶδος ; Max.Tyr.~18.4b Hobein Ἀλκιβιάδην τὸν γαυρότατον, καὶ Κριτόβουλον τὸν ὡραιότατον, καὶ Ἀγάθωνα τὸν ἁβρότατον, καὶ Φαῖδρον τὴν θείαν κεφαλήν ; Hsch.~α~5759 ἁπαλός· μικρός, τρυφερός ; Suid.~α~2885 ἁπαλοὶ θερμολουσίαις, ἁβροὶ μαλθακευνίαις. ἐπὶ τῶν ὑπὸ τρυφῆς καὶ ἁβρότητος διαρρεόντων ; \textit{idem} App. Prov.~1.34 (om. ἁβροὶ μαλθακευνίαις ; ἀντὶ τοῦ τρυφηλοὶ καὶ τὴν σάρκα διαρρέοντες pro ἐπὶ τῶν … διαρρεόντων) ; Ar.~\textit{Lys.}~417–418 τὸ δακτυλίδιον ξυμπιέζει τὸ ζυγὸν / ἅθ᾽ ἁπαλὸν ὄν ; Ar.~\textit{Eccl.}~901–902 τὸ τρυφερὸν γὰρ ἐμπέφυκε / τοῖς ἁπαλοῖσι μηροῖς ; Alex.~fr.~49.2–3~K.-A. (=~48~Kock ; Ath.~7.314d) ὅτι δεῖ γυναικὸς ἐπιφερούσης δακτύλους / ἁπαλοὺς ὑπ᾽ ἀκάνθης μηδὲ ἓν τούτους παθεῖν ; Timocl.~fr.~24.4–6~K.-A. (=~22~Kock ; Ath.~13.570f) δεῖν δέ τι / ἀγωνιᾶσαι καὶ ῥαπισθῆναί τε καὶ / πληγὰς λαβεῖν ἁπαλαῖσι χερσίν ; \textit{Od.}~21.150–151 πρὶν γὰρ κάμε χεῖρας ἀνέλκων / ἀτρίπτους ἁπαλάς ; Zonar. (Tittmann) col.~233 Ἁπαλός. μαλακός. παρὰ τὴν ἁφὴν ἁφαλὸς καὶ ἁπαλός}
\lp{192 ἁπαλός (Ἀγαθώνειος αὔλησις)]}{Hsch.~α~281 Ἀγαθώνειον αὔλησιν τὴν μαλακὴν (Ar. fr. 178 K.–A.). Ἀγαθὼν γὰρ ὁ τραγικὸς ἐπὶ μαλακίᾳ διεβάλλετο ; Suid.~α~124 Ἀγάθων· ὄνομα κύριον. τραγικός δὲ ἦν· διεβέβλητο δὲ ἐπὶ μαλακίᾳ ; Zen.~1.2 Ἀγαθώνιος αὔλησις: ἡ μαλακὴ, καὶ μήτε πικρὰ μήτε χαλαρὰ, ἀλλ᾽ εὔκρατος καὶ ἡδίστη ; \textit{idem} Suid.~α~125 (ἡ μαλακὴ καὶ ἐκλελυμένη· ἢ ἡ μήτε χαλαρά, μήτε πικρά), Diogenian.~1.7 (Ἀγαθώνειος ; μαλθακὴ καὶ μήτε χλιαρὰ μήτε πικρά ; om. εὔκρατος καὶ) ; Erasm.~\textit{Adag.}~2.4.63 (LB II, col.~543 E) Ἀγαθώνειος αὔλησις. i. e. Agathonia cantio. … Laboravit enim infamia molliciei. Auctor \& Suidas \& Zenodotus. Haud scio, an hic sit idem ille Agathon in Tragicorum certamine victor, in cujus epiniciis Plato celebrat illud nobile convivium. Hujus multis modis Aristophanes in Thesmophoriis carminum molliciem deridet}
\lp{192 εὐπρεπὴς ἰδεῖν]}{Ar.~\textit{Thesm.}~233 μὴ φροντίσῃ· ὡς εὐπρεπὴς φανεῖ πάνυ ; Ar.~\textit{Eccl.}~427 μετὰ τοῦτο τοίνυν εὐπρεπὴς νεανίας ; X.~\textit{Mem.}~2.1.22 τὴν μὲν ἑτέραν εὐπρεπῆ τε ἰδεῖν καὶ ἐλευθέριον φύσει ; Aeschin.~1.157 ὃς εὐπρεπὴς ὢν ἰδεῖν τοσοῦτον ἀπέχει τῶν αἰσχρῶν ; Ar.~\textit{Nub.}~1172 νῦν μέν γ᾽ ἰδεῖν εἶ πρῶτον ἐξαρνητικὸς ; Anaxandr.~fr.~35.1–2~K.-A. (=~34~Kock ; Ath.~6.242d–e) ὑμεῖς γὰρ ἀλλήλους ἀεὶ χλευάζετ᾽, οἶδ᾽ ἀκριβῶς· / ἂν μὲν γὰρ ᾖ τις εὐπρεπής, ἱερὸν γάμον καλεῖτε ; Eur.~\textit{IA}~822 γυναῖκα, μορφὴν εὐπρεπῆ κεκτημένην ; Xenoph.~fr.~3.5–6~Bergk⁴ (Ath.~12.526a) αὐχαλέοι, χαίτῃσιν ἀγαλλόμενοι εὐπρεπέεσσιν, / ἀσκητοῖς ὀδμὴν χρίμασι δευόμενοι ; Hdt.~1.60.4 ταύτην τὴν γυναῖκα σκευάσαντες πανοπλίῃ … σχῆμα οἷόν τι ἔμελλε εὐπρεπέστατον φανέεσθαι ἔχουσα ; Aeschin.~1.155 ὧν τοῖς μὲν διὰ τὴν εὐπρέπειαν πολλοὶ γεγόνασιν ἐρασταί ; Ar.~\textit{Av.}~1709–1710 οὔτε παμφαὴς / ἀστὴρ ἰδεῖν ἔλαμψε ; Ar.~\textit{Thesm.}~682 πᾶσιν ἐμφανὴς ὁρᾶν ἔσται γυναιξὶ καὶ βροτοῖσιν ; Eur.~\textit{IA}~386 εὐπρεπῆ γυναῖκα χρῄζεις, τὸ λελογισμένον παρεὶς ; Hsch.~ε~7085 εὐπρέπεια· εὐμορφία (2. Regn. 15, 25 . .) ; Hsch.~θ~105 Θαργηλία· ἔστιν ἡ Θαργηλία Μιλησία μὲν τὸ γένος, εὐπρεπὴς δὲ τὴν ὄψιν}
\lp{193 Εὐριπίδη — τί ἔστιν;]}{Ar.~\textit{Ach.}~410–411 Εὐριπίδη, — τί λέλακας; — ἀναβάδην ποιεῖς, / ἐξὸν καταβάδην; οὐκ ἐτὸς χωλοὺς ποιεῖς ; Ar.~\textit{Av.}~1010 Μέτων — — τί ἔστιν; — οἶσθ᾽ ὁτιὴ φιλῶ σ᾽ ἐγώ ; Ar.~\textit{Ran.}~1220–1221 Εὐριπίδη— — τί ἔσθ᾽; — ὑφέσθαι μοι δοκεῖ· / τὸ ληκύθιον γὰρ τοῦτο πνευσεῖται πολύ ; Ar.~\textit{Ach.}~1085 Δικαιόπολι. — τί ἔστιν; — ἐπὶ δεῖπνον ταχὺ ; Ar.~\textit{Nub.}~82 ἰδού. τί ἔστιν; — εἰπέ μοι, φιλεῖς ἐμέ ; Ar.~\textit{Thesm.}~95 σίγα. — τί δ᾽ ἔστιν; — Ἁγάθων ἐξέρχεται}
\lp{193 ἐποίησάς ποτε]}{\textit{Test.} Eust.~4.9–10 δῆλον ὅτι ποιεῖν παρὰ τοῖς παλαιοῖς, ὡς καὶ ὁ κωμικὸς ἐμφαίνει, τὸ ποιητικῶς ἤγουν ἐμμέτρως ᾄδειν καὶ γράφειν, ὡς τό, ἐποίησάς ποτε, ὦ Εὐριπίδη, ἀντὶ τοῦ, ἔγραψας ἐν τῇ ποιήσει σου ; Ar.~\textit{Av.}~920 ταυτὶ σὺ πότ᾽ ἐποίησας; ἀπὸ ποίου χρόνου ; Ar.~\textit{Thesm.}~149–150 χρὴ γὰρ ποιητὴν ἄνδρα πρὸς τὰ δράματα / ἃ δεῖ ποιεῖν, πρὸς ταῦτα τοὺς τρόπους ἔχειν ; Ar.~\textit{Ran.}~1044 οὐδ᾽ οἶδ᾽ οὐδεὶς ἥντιν᾽ ἐρῶσαν πώποτ᾽ ἐποίησα γυναῖκα ; Ar.~\textit{Ran.}~1177 σὺ δὲ πῶς ἐποίεις τοὺς προλόγους ; Hdt.~3.38.4 καὶ ὀρθῶς μοι δοκέει Πίνδαρος ποιῆσαι νόμον πάντων βασιλέα φήσας εἶναι ; Arist.~\textit{Po.}~14.12 καθάπερ καὶ Εὐριπίδης ἐποίησεν ἀποκτείνουσαν τοὺς παῖδας τὴν Μήδειαν ; Hdt.~4.16.1 ἐν αὐτοῖσι τοῖσι ἔπεσι ποιέων ἔφησε ; Aeschin.~1.135 ὅσα πεποίηκα ἐρωτικὰ εἴς τινας ποιήματα}
\lp{193–194 ἐποίησάς ποτε … (τοῖς αὑτῶν πτεροῖς)]}{Ar.~\textit{Thesm.}~275–276 μέμνησο τοίνυν ταῦθ᾽, ὅτι ἡ φρὴν ὤμοσεν, / ἡ γλῶττα δ᾽ οὐκ ὀμώμοκ᾽ ; Ar.~\textit{Ran.}~1471 ἡ γλῶττ᾽ ὀμώμοκ᾽, Αἰσχύλον δ᾽ αἱρήσομαι ; Ar.~\textit{Ran.}~1477 τίς οἶδεν εἰ τὸ ζῆν μέν ἐστι κατθανεῖν ; Aesch.~fr.~139~N.² =~R. ὧδ᾽ ἐστὶ μύθων τῶν Λιβυστικῶν κλέος, / πληγέντ᾽ ἀτράκτῳ τοξικῷ τὸν αἰετὸν / εἰπεῖν ἰδόντα μηχανὴν πτερώματος· / τάδ᾽ οὐχ ὑπ᾽ ἄλλων, ἀλλὰ τοῖς αὑτῶν πτεροῖς / ἁλισκόμεσθα ; \textit{idem} Ar.~\textit{Av.}~807–808 (praem. ταυτὶ μὲν ᾐκάσμεσθα κατὰ τὸν Αἰσχύλον· ; v. 4 tantum), D.H.~\textit{Dem.}~7 (λόγοις pro πτεροῖς ; add. ⟨κατὰ τὴν τραγῳδί⟩αν), Ath.~11.494b (καὶ ταῦτ᾽ οὐχ ὑπ᾽ ἄλλων, ἀλλὰ τοῖς αὑτοῦ πτεροῖς κατὰ τὸν θαυμάσιον Αἰσχύλον ἁλίσκῃ), Gal.~\textit{PHP}~4.5 (V, p.~395 Kühn) (καὶ τάχ᾽ οὐχ ὑπ᾽ ἄλλων … ἁλισκόμεθα), Macar.~8.57 (CPG II, p.~222) (τοῖς σαυτοῦ πτεροῖς ἥλως), Suid.~ο~992 (vv. 4–5 tantum ; om. τάδ᾽ ; αὐτῶν ; ἁλισκόμεθα ; add. ἀντὶ τοῦ ἑαυτοῖς ταῦτα πεποιήκαμεν) ; Aristid.~\textit{Or.}~45 (II, p.~17 Dindorf = p.~15 Jebb) ὥστ᾽ εἰ μὲν βούλει, τὸ τοῦ Εὐριπίδου καὶ τῶν αὐτοῦ νὴ Δία συμβαίνει, σοῦ τάδ᾽, οὐκ ἐμοῦ κλύεις, ὦ τοσοῦτον ἐνθυμηθείς· εἰ δ᾽ αὖ βούλει, τὸ τοῦ Αἰσχύλου, τάδ᾽ οὐχ ὑπ᾽ ἄλλων, ἀλλὰ τοῖς αὑτοῦ πτεροῖς, ἃ θρυλεῖς ἐν ἐκείνοις τοῖς λόγοις, ἁλίσκει ; Tosi n°~2242 Hunc telo suo sibi… pellere ; τοῖς αὐτῶν πτεροῖς ; Erasm.~\textit{Adag.}~1.6.52 (LB II, col.~243 C–E) Τοῖς αὐτῶν πτεροῖς ἁλισκόμεσθα. i. e. Nostris ipsorum pennis capimur. … Convenit in eos, qui ipsi praebent occasionem sui mali ; Eur.~\textit{Hipp.}~352 σοῦ τάδ᾽, οὐκ ἐμοῦ κλύεις ; Aristid.~\textit{Or.}~53 (II, p.~624 Dindorf) οὐ μόνον ὅτι τοῖς οἰκείοις πτεροῖς καὶ λόγοις ἁλίσκει, ὃ δεινότατόν ἐστι πάντων καὶ συμφορὰ τοῖς χρωμένοις ; Hld.~\textit{Aeth.}~2.33 ἀλλὰ τὸ χαλεπώτατον, τοῖς ἐμοῖς, τὸ τοῦ λόγου, κατ᾽ ἐμοῦ κέχρηται πτεροῖς ; Thdt.~\textit{HE}~3.8.2 τοῖς οἰκείοις γάρ, φησί, πτεροῖς κατὰ τὴν παροιμίαν βαλλόμεθα· ἐκ γὰρ τῶν ἡμετέρων συγγραμμάτων καθοπλιζόμενοι τὸν καθ᾽ ἡμῶν ἀναδέχονται πόλεμον ; Σ Luc.~\textit{Hermot.}~53 ἐπεὶ καὶ αὐτὸς σὺ ὁ ταῦτα ληρῶν ἄνθρωπος … καὶ οὕτω τοῖς οἰκείοις ἀλώσῃ πτεροῖς}
\lp{194 χαίρεις ὁρῶν φῶς, πατέρα δ᾽ οὐ χαίρειν δοκεῖς;]}{Ar.~\textit{Nub.}~1415 κλάουσι παῖδες, πατέρα δ᾽ οὐ κλάειν δοκεῖς ; Ar.~\textit{Av.}~1244 ἔχ᾽ ἀτρέμα. φέρ᾽ ἴδω, πότερα Λυδὸν ἢ Φρύγα ; Eur.~\textit{Alc.}~690–691 μὴ θνῇσχ᾽ ὑπὲρ τοῦδ᾽ ἀνδρός, οὐδ᾽ ἐγὼ πρὸ σοῦ. / χαίρεις ὁρῶν φῶς· πατέρα δ᾽ οὐ χαίρειν δοκεῖς ; \textit{idem} Σ Ar.~\textit{Nub.}~1415 (praem. παρὰ τὰ ἐξ Ἀλκήστιδος Εὐριπίδου ; v. 691 tantum) ; Eur.~\textit{Hec.}~1256 ἀλγεῖς· τί δ᾽; ἦ ᾽μὲ παιδὸς οὐκ ἀλγεῖν δοκεῖς ; Eur.~fr.~710~Kannicht τὸν δὲ Τήλεφον / οὐκ οἰόμεσθα; νοῦς ἄρ᾽ ὑμῖν οὐκ ἔνι ; \textit{idem} Ar.~\textit{Ach.}~555–556 (praem. ταῦτ᾽ οἶδ᾽ ὅτι ἂν ἐδρᾶτε· ; ἡμῖν) ; Epict.~2.22.11 ἐρεῖς τὰς φωνὰς τὰς τοῦ Ἀδμήτου πατρός· [χαίρεις ὁρῶν φῶς, πατέρα δ᾽ οὐ χαίρειν δοκεῖς;] θέλεις βλέπειν φῶς, πατέρα δ᾽ οὐ θέλειν δοκεῖς ; Σ Ar.~\textit{Thesm.}~194 ἐξ Ἀλκήστιδος Εὐριπίδου ; Ar.~\textit{Ach.}~893 ἀλλ᾽ ἔσφερ᾽ αὐτήν· μηδὲ γὰρ θανών ποτε ; Ar.~\textit{Eq.}~1251–1252 σὲ δ᾽ ἄλλος τις λαβὼν κεκτήσεται, / κλέπτης μὲν οὐκ ἂν μᾶλλον, εὐτυχὴς δ᾽ ἴσως ; Ar.~\textit{Lys.}~763–764 ποθεῖτ᾽ ἴσως τοὺς ἄνδρας· ἡμᾶς δ᾽ οὐκ οἴει / ποθεῖν ἐκείνους ; Eur.~\textit{Alc.}~675–676 ὦ παῖ, τίν᾽ αὐχεῖς, πότερα Λυδὸν ἢ Φρύγα / κακοῖς ἐλαύνειν ἀργυρώνητον σέθεν ; Ar.~\textit{Vesp.}~751–752 μή μοι τούτων μηδὲν ὑπισχνοῦ. / κείνων ἔραμαι, κεῖθι γενοίμαν}
\lp{194 χαίρεις … δοκεῖς; (γέρων φιλόζωος)]}{Eur.~\textit{Alc.}~669–672 μάτην ἄρ᾽ οἱ γέροντες εὔχονται θανεῖν, / γῆρας ψέγοντες καὶ μακρὸν χρόνον βίου· / ἢν δ᾽ ἐγγὺς ἔλθῃ θάνατος, οὐδεὶς βούλεται / θνῄσκειν, τὸ γῆρας δ᾽ οὐκέτ᾽ ἔστ᾽ αὐτοῖς βαρύ ; Arist.~\textit{Rh.}~2.13.8 (1389b30–35) καὶ φιλόζωοι, καὶ μᾶλλον ἐπὶ τῇ τελευταίᾳ ἡμέρᾳ διὰ τὸ τοῦ ἀπόντος εἶναι τὴν ἐπιθυμίαν, καὶ οὗ ἐνδεεῖς, τούτου μάλιστα ἐπιθυμεῖν ; Cic.~\textit{Sen.}~24 nemo enim est tam senex qui se annum non putet posse vivere ; \textit{idem} Tosi n°~1085 (om. enim) ; Sen.~\textit{Ep.}~12.6 Deinde nemo tam senex est, ut inprobe unum diem speret ; Σ Eur.~\textit{Hec.}~497 (Dindorf) γέρων μέν εἰμι, ἀλλὰ φιλόζωος. … ἢ τὸ “γέρων μέν εἰμι᾽ διὰ τὸ τοῖς γέρουσι περισπουδαστότερον εἶναι τὸ ζῆν, ἐπειδὴ ὁρῶσιν ὡς ἐγγύς εἰσι τοῦ θανάτου}
\lp{194 ὁρῶν φῶς]}{Eur.~\textit{Alc.}~722 φίλον τὸ φέγγος τοῦτο τοῦ θεοῦ, φίλον ; Eur.~\textit{IA}~1218 μή μ᾽ ἀπολέσῃς ἄωρον· ἡδὺ γὰρ τὸ φῶς ; Eur.~\textit{IA}~1250 τὸ φῶς τόδ᾽ ἀνθρώποισιν ἥδιστον βλέπειν ; Eur.~\textit{IT}~608 ὃν οὐδὲν ἧσσον ἢ ᾽μὲ φῶς ὁρᾶν θέλω ; \textit{Il.}~18.61 ὄφρα δέ μοι ζώει καὶ ὁρᾷ φάος ἠελίοιο ; Luc.~\textit{DMort.}~22.9 ἡδὺ γὰρ ἦν τὸ φῶς καὶ τὸ τεθνάναι δεινὸν ; Eur.~\textit{Heracl.}~969 χρῆν τόνδε μὴ ζῆν μηδ᾽ ὁρᾶν φάος ἔτι ; Eur.~\textit{IA}~1394 εἷς γ᾽ ἀνὴρ κρείσσων γυναικῶν μυρίων ὁρᾶν φάος ; Σ Eur.~\textit{Or.}~386 (Dindorf) οὐ γὰρ ζῷ : οὐκ ἄλλο κέκτημαι τῶν ζώντων εἰ μὴ μόνον τὸ ὁρᾶν τὸ φάος}
\lp{195–197 μή νυν … οἰκείως φέρε]}{\textit{Test.} Suid.~υ~709 Ὑφέξειν· ὑπενεγκεῖν, ὑπομεῖναι. μὴ νῦν ἐλπίσῃς τὸ σὸν κακὸν ἡμᾶς ὑφέξειν. καὶ γὰρ ἂν μαινοίμεθ᾽ ἄν. αὐτὸς ὅ γε σόν ἐστιν οἰκείως φέρε}
\lp{195 ἔγωγε]}{Ar.~\textit{Ach.}~776–777 ἦ λῇς ἀκοῦσαι φθεγγομένας; — νὴ τοὺς θεοὺς / ἔγωγε ; Ar.~\textit{Eq.}~32–33 ἐτεὸν ἡγεῖ γὰρ θεούς; — ἔγωγε. — ποίῳ χρώμενος τεκμηρίῳ ; Ar.~\textit{Eccl.}~90–91 νὴ τὴν Ἄρτεμιν / ἔγωγε ; Soph.~\textit{Trach.}~1247–1248 πράσσειν ἄνωγας οὖν με πανδίκως τάδε; — ἔγωγε· τούτων μάρτυρας καλῶ θεούς}
\lp{195 μή νυν ἐλπίσῃς]}{Eur.~\textit{Alc.}~1077 μή νυν ὑπέρβαλλ᾽, ἀλλ᾽ ἐναισίμως φέρε ; Eur.~\textit{Heracl.}~1051 δοῦναι κτανόντας· μὴ γὰρ ἐλπίσῃς ὅπως ; Ar.~\textit{Nub.}~716 μή νυν βαρέως ἄλγει λίαν ; Eur.~\textit{Hipp.}~1456 μή νυν προδῷς με, τέκνον, ἀλλὰ καρτέρει ; Hdt.~7.157.3 μὴ γὰρ ἐλπίσῃς, ἢν ἡμέας καταστρέψηται ὁ Πέρσης μάχῃ κρατήσας, ὡς οὐκὶ ἥξει παρὰ σέ γε ; Suid.~ε~912 Ἐλπίσας: οὐ μόνον ἐπ᾽ ἀγαθῷ, ἀλλ᾽ ἁπλῶς ἐπὶ τῇ τοῦ μέλλοντος ἐκβάσει λέγεται παρὰ Ἀττικοῖς ; Eur.~\textit{Hel.}~1419 μή νυν ἄγαν σὸν δάκρυσιν ἐκτήξῃς χρόα ; Aesch.~\textit{Pers.}~746 ἤλπισε σχήσειν ῥέοντα, Βόσπορον ῥόον θεοῦ ; Soph.~\textit{Ant.}~705 μή νυν ἓν ἦθος μοῦνον ἐν σαυτῷ φόρει ; Hsch.~ε~4038 ἔολπεν· ἐλπίζει, προσδοκᾷ, οἴεται}
\lp{195 τὸ σὸν κακόν]}{Epict.~3.18.8 σὸν κακόν ἐστι τὸ κακῶς ἀπολογηθῆναι· τοῦτο φυλάσσου μόνον· κριθῆναι δ᾽ ἢ μὴ κριθῆναι ὥσπερ ἄλλου ἐστὶν ἔργον, οὕτως κακὸν ἄλλου ἐστίν ; Eur.~\textit{Tro.}~982 τὸ σὸν κακὸν κοσμοῦσα, μὴ οὐ πείσῃς σοφούς ; Pl.~\textit{R.}~1.347a μηδένα ἐθέλειν ἑκόντα ἄρχειν καὶ τὰ ἀλλότρια κακὰ μεταχειρίζεσθαι ἀνορθοῦντα ; Epict.~2.16.45 οὐ δύνασαι καθαίρειν τὰ ἀλλότρια κακά, ἀλλ᾽ οὐδὲ Θησεύς, ἵνα τὰ τῆς Ἀττικῆς καθάρῃς· τὰ σαυτοῦ κάθαρον ; Hsch.~υ~523 ὕπεχε· ὑπόφερε. πάρεχε. ὑπόμενε (Ar. Pac. 431) ; Erasm.~\textit{Adag.}~1.1.84 (LB II, col.~59 F –~60 A) In me haec cudetur faba. Terentius in Eunucho : At enim isthaec in me cudetur faba : h. e. Donato interprete, in me malum hoc recidet, in me haec vindicabitur culpa}
\lp{196 ἡμᾶς ὑφέξειν]}{Soph.~\textit{OT}~551–552 εἴ τοι νομίζεις ἄνδρα συγγενῆ κακῶς / δρῶν οὐχ ὑφέξειν τὴν δίκην, οὐκ εὖ φρονεῖς ; Ar.~\textit{Lys.}~839–841 σὸν ἔργον ἤδη τοῦτον ὀπτᾶν καὶ στρέφειν / κἀξηπεροπεύειν καὶ φιλεῖν καὶ μὴ φιλεῖν, / καὶ πάνθ᾽ ὑπέχειν πλὴν ὧν σύνοιδεν ἡ κύλιξ ; Ar.~\textit{Thesm.}~183 τίς οὖν παρ᾽ ἡμῶν ἐστιν ὠφέλειά σοι ; Eur.~\textit{Andr.}~358 αὐτοὶ τὴν δίκην ὑφέξομεν ; Eur.~\textit{Hec.}~1253 δούλης ὑφέξω τοῖς κακίοσιν δίκην ; Eur.~\textit{Ion}~1308 σὺ δ᾽ οὐχ ὑφέξεις ζημίαν, κτείνουσ᾽ ἐμέ ; Soph.~\textit{Trach.}~1274 τῷ τήνδ᾽ ἄτην ὑπέχοντι ; Hsch.~υ~887 ὑφέξω· δώσω (Eur. Hec. 1253 . .) ; Ar.~\textit{Thesm.}~756 ὕπεχ᾽ αὐτό, χαριοῦμαι γὰρ ἕν γε τοῦτό σοι ; D.Chr.~46.1 μείζω ὑμῖν ὑφέξω τιμωρίαν ἢ αὐτοὶ ζητεῖτε ; Zen.~3.20 Δίκην ὑφέξει κἂν ὄνος δάκῃ κύνα: ἐπὶ τῶν ἐπὶ μικροῖς συκοφαντουμένων}
\lp{196 καὶ γὰρ ἂν μαινοίμεθ᾽ ἄν]}{Ar.~\textit{Plu.}~1070 μὰ τὴν Ἑκάτην οὐ δῆτα· μαινοίμην γὰρ ἄν ; Eur.~\textit{IA}~1256 φιλῶ τ᾽ ἐμαυτοῦ τέκνα· μαινοίμην γὰρ ἄν ; D.~21.208 καὶ γὰρ ἂν μαινοίμην ; \textit{idem} D.~35.40 (μαινοίμην γὰρ ἄν), D.~\textit{Ex.}~45.4 (praem. ἐγὼ δ᾽ οὐ χείρους ὑμᾶς ἡγοῦμαι φύσει Θηβαίων), Lys.~21.22 (μαινοίμην γὰρ ⟨ἄν⟩), Longus~3.31.2 (add. εἰ μὴ γέρων τε ὢν ἤδη) ; Ar.~\textit{Thesm.}~469–470 καὐτὴ γὰρ ἔγωγ᾽, οὕτως ὀναίμην τῶν τέκνων, / μισῶ τὸν ἄνδρ᾽ ἐκεῖνον, εἰ μὴ μαίνομαι ; Hsch.~μ~87 μαίνεται· ὀργίζεται. ὑλακτεῖ}
\lp{196 ἂν … ἄν]}{Ar.~\textit{Ran.}~581 οὐκ ἂν γενοίμην Ἡρακλῆς ἄν ; Ar.~\textit{Ran.}~1450 τἀναντί᾽ ⟨ἂν⟩ πράττοντες οὐ σῳζοίμεθ᾽ ἄν ; Eur.~\textit{Hipp.}~270 σοῦ δ᾽ ἂν πυθέσθαι καὶ κλύειν βουλοίμεθ᾽ ἄν ; Soph.~\textit{OT}~339 τίς γὰρ τοιαῦτ᾽ ἂν οὐκ ἂν ὀργίζοιτ᾽ ἔπη ; Ar.~\textit{Ach.}~307 πῶς δέ γ᾽ ἂν καλῶς λέγοις ἄν ; Ar.~\textit{Lys.}~191 τίς ἂν οὖν γένοιτ᾽ ἂν ὅρκος ; Eur.~\textit{Hipp.}~961 τῆσδ᾽ ἂν γένοιντ᾽ ἄν, ὥστε σ᾽ αἰτίαν φυγεῖν ; Eur.~\textit{Tro.}~1244 ἀφανεῖς ἂν ὄντες οὐκ ἂν ὑμνήθημεν ἂν}
\lp{197 ἀλλ᾽ αὐτὸς ὅ γε σόν ἐστιν οἰκείως φέρε]}{Eur.~\textit{Alc.}~685–686 σαυτῷ γὰρ εἴτε δυστυχὴς εἴτ᾽ εὐτυχὴς / ἔφυς ; Tosi n°~385 Πράττων τὰ σαυτοῦ, μὴ τὰ τῶν ἄλλων φρόνει ; Soph.~\textit{Phil.}~1247 καὶ πῶς δίκαιον, ἅ γ᾽ ἔλαβες βουλαῖς ἐμαῖς, πάλιν μεθεῖναι ταῦτα ; Ep.Gal.~6.5 ἕκαστος γὰρ τὸ ἴδιον φορτίον βαστάσει ; App. Prov.~2.37 (CPG I, p.~400) Ἕκαστος τὸν ἑαυτοῦ ὀβελίσκον τρέφει: ἐπὶ τῶν φροντίδα τῶν οἰκείων ποιουμένων ; Soph.~\textit{El.}~441 οὐκ ἄν ποθ᾽ ὅν γ᾽ ἔκτεινε, τῷδ᾽ ἐπέστεφε ; Soph.~\textit{El.}~923 πῶς δ᾽ οὐκ ἐγὼ κάτοιδ᾽ ἅ γ᾽ εἶδον ἐμφανῶς ; Soph.~\textit{OT}~342 οὐκοῦν ἅ γ᾽ ἥξει καὶ σὲ χρὴ λέγειν ἐμοί ; Phryn.~\textit{Ecl.}~409 Fischer Τὰ ἴδια πράττω καὶ τὰ ἴδια πρᾶττε οἱ πολλοὶ λέγουσιν εἰκῇ, δέον τὰ ἐμαυτοῦ πράττω καὶ τὰ σαυτοῦ πρᾶττε λέγειν ὡς οἱ παλαιοί}
\lp{197 ἀλλ᾽ αὐτὸς … οἰκείως φέρε (ἣν τις ἔμαξε μᾶζαν)]}{Ter.~\textit{Phorm.}~318 Tute hoc intristi: tibi omne est exedendum ; \textit{idem} Tosi n°~890 (omnest) ; Macar.~4.50 (CPG II, p.~171) Ἣν τις ἔμαξε μᾶζαν ταύτην καὶ ἐσθιέτω ; Erasm.~\textit{Adag.}~1.1.85 (LB II, col.~60 A–B) Tute hoc intristi, omne tibi exedendum est. Non diversum superiori Donatus admonuit, illud quod apud eundem poetam legitur in Phormione : Tute hoc intristi, tibi omne est exedendum}
\lp{197 οἰκείως φέρε]}{Ar.~\textit{Lys.}~1118 ἀλλ᾽ ὡς γυναῖκας εἰκός, οἰκείως πάνυ ; Eur.~\textit{Tro.}~406 ὡς ἡδέως κακοῖσιν οἰκείοις γελᾷς ; Soph.~\textit{Ant.}~1187 χαλῶσα, καί με φθόγγος οἰκείου κακοῦ ; Soph.~\textit{OC}~765 πρόσθεν τε γάρ με τοῖσιν οἰκείοις κακοῖς ; Hsch.~ο~232 οἰκεία· ἰδία ; Men.~fr.~236.12~K.-A. (Μισογύνης ; Stob.~4.22.71) ἔθαψε, περιέστειλεν οἰκείως ; Damox.~fr.~2.58~K.-A. (Ath.~3.103a) καὶ ταῖς ἐπιφοραῖς εὐθὺς οἰκείως πλέκω ; Th.~2.60.6 ὅ τε ἔχων ἀμφότερα, τῇ δὲ πόλει δύσνους, οὐκ ἂν ὁμοίως τι οἰκείως φράζοι ; Lys.~1.39 φίλος ὢν Σώστρατος καὶ οἰκείως διακείμενος}
\lp{198–199 τὰς συμφορὰς γὰρ … ἀλλὰ τοῖς παθήμασιν]}{\textit{Test.} Suid.~τ~440 Τεχνάσμασι: δόλοις, πανουργίαις. τὰς συμφορὰς γὰρ οὐχὶ τοῖς τεχνάσμασι φέρειν δίκαιον, ἀλλὰ τοῖς παθήμασι. καὶ μὴν σύ γ᾽, ὦ καταπῦγον, εὐρύπρωκτος εἶ, οὐ τοῖς λόγοισιν, ἀλλὰ τοῖς παθήμασιν ; Agatho TrGF~39 F~34? τὰς συμφορὰς γὰρ οὐχὶ τοῖς τεχνάσμασιν / φέρειν δίκαιον, ἀλλὰ τοῖς παθήμασιν}
\lp{198–199 οὐχὶ τοῖς τεχνάσμασιν … ἀλλὰ τοῖς παθήμασιν (Ἀγάθωνος ἀντίθετα)]}{Ar.~\textit{Thesm.}~146–147 τοῦ φθόνου μὲν τὸν ψόγον / ἤκουσα, τὴν δ᾽ ἄλγησιν οὐ παρεσχόμην ; Agatho fr.~6~N.² (Arist.~\textit{EN}~6.4) τέχνη τύχην ἔστερξε καὶ τύχη τέχνην ; Agatho fr.~7~N.² (Arist.~\textit{EE}~3.1) φαῦλοι βροτῶν γὰρ τοῦ πονεῖν ἡσσώμενοι, / θανεῖν ἐρῶσιν ; Agatho fr.~8~N.² (Arist.~\textit{Rh.}~2.19.13) τὰ δὲ / ἡμῖν ἀνάγκῃ καὶ τύχῃ προσγίγνεται ; Agatho fr.~11~N.² (Ath.~5.185a) τὸ μὲν πάρεργον ἔργον ὣς ποιούμεθα, τὸ δ᾽ ἔργον ὡς πάρεργον ἐκπονούμεθα ; Agatho fr.~12~N.² (Ath.~5.211e) εἰ μὲν φράσω τἀληθές, οὐχί σ᾽ εὐφρανῶ· εἰ δ᾽ εὐφρανῶ τί σ᾽, οὐχὶ τἀληθὲς φράσω ; Agatho fr.~14~N.² (Ath.~13.584a) γυνὴ τοι σώματος δι᾽ ἀργίαν ψυχῆς φρόνησιν ἐντὸς οὐκ ἀργὸν φορεῖ ; Pl.~\textit{Smp.}~197d οὗτος δὲ ἡμᾶς ἀλλοτριότητος μὲν κενοῖ, οἰκειότητος δὲ πληροῖ ; Pl.~\textit{Smp.}~198c καὶ γάρ με Γοργίου ὁ λόγος ἀνεμίμνῃσκεν ; Ael.~\textit{VH}~14.13 πολλοῖς καὶ πολλάκις χρῆται τοῖς ἀντιθέτοις ὁ Ἀγάθων ; Philostr.~\textit{VS}~1.9.1 καὶ Ἀγάθων δὲ ὁ τῆς τραγῳδίας ποιητής, ὃν ἡ κωμῳδία σοφόν τε καὶ καλλιεπῆ οἶδε, πολλαχοῦ τῶν ἰάμβων γοργιάζει ; Ath.~5.187c χλευάζει τε τὰ ἰσόκωλα τὰ Ἀγάθωνος καὶ τὰ ἀντίθετα ; Agatho fr.~5~N.² (Arist.~\textit{EN}~6.2) μόνου γὰρ αὐτοῦ καὶ θεὸς στερίσκεται, / ἀγένητα ποιεῖν}
\lp{198 τὰς συμφορὰς γὰρ … φέρειν]}{Antiph.~fr.~320~K.-A. (=~278~Kock ; Stob. Flor.~108.29a~Meineke) τὸ μὲν ἀτυχῆσαι παντὸς εἶναί μοι δοκεῖ, / ἀνδρὸς δ᾽ ἐνεγκεῖν ἀτυχίαν ὀρθῷ τρόπῳ ; \textit{idem} Apostol.~16.10a (CPG II, p.~658) (praem. τὰ τῆς τύχης φέρειν δεῖ ἠπίως τὸν εὐμενῆ (cf. Antiph. fr. 321) ; ὀρθοῦ τρόπου ; add. Ἀντιφάνους) ; Antiph.~fr.~321~K.-A. (=~281~Kock ; Stob.~4.44.22 = Flor.~108.29~Meineke) τὰ τύχης φέρειν δεῖ γνησίως τὸν εὐγενῆ ; Alex.~fr.~254~K.-A. (Stob. Flor.~108.52~Meineke) Σοφοῦ γὰρ ἀνδρὸς τὰς τύχας ὀρθῶς φέρειν ; \textit{idem} Tosi n°~118 ; Men. frr.~874–875~K.-A. (=~816–817~Kock ; Stob.~4.56.6a–b = Flor.~124.6~Meineke) ἀνθρωπίνως δεῖ τὰς τύχας φέρειν, ξένε. / τὰ κοινὰ κοινῶς χρὴ φέρειν συμπτώματα ; \textit{idem} Apostol.~3.1c (CPG II, p.~288) (δεῖ pro χρή ; add. Μενάνδρου) ; Philippid.~fr.~18.1–4~K.-A. (=~18~Kock ; Stob.~4.44.10 = Flor.~108.10~Meineke) ὅταν ἀτυχεῖν σοι συμπέσῃ τι, δέσποτα, / Εὐριπίδου μνήσθητι, καὶ ῥᾴων ἔσῃ· / οὐκ ἔστιν ὅστις πάντ᾽ ἀνὴρ εὐδαιμονεῖ. / εἶναι δ᾽ ὑπόλαβε καὶ σὲ τῶν πολλῶν ἕνα ; Eur.~\textit{Alc.}~416 Ἄδμητ᾽, ἀνάγκη τάσδε συμφορὰς φέρειν ; Eur.~\textit{Alc.}~1078 ῥᾷον παραινεῖν ἢ παθόντα καρτερεῖν ; Eur.~\textit{HF}~1227–1228 ὅστις εὐγενὴς βροτῶν, / φέρει τά γ᾽ ἐκ θεῶν πτώματ᾽ οὐδ᾽ ἀναίνεται ; Eur.~\textit{HF}~1349–1350 ap. Stob. Flor.~108.12~Meineke Τὰς συμφορὰς γὰρ ὅστις οὐκ ἐπίσταται / θνητὸς πεφυκὼς ὃν τρόπον χρεὼν φέρειν, / οὐδ᾽ ἀνδρὸς ἂν δύναιθ᾽ ὑποστῆναι βέλος ; Eur.~\textit{Med.}~1018 κούφως φέρειν χρὴ θνητὸν ὄντα συμφοράς ; Eur.~\textit{Phoen.}~1763 τὰς γὰρ ἐκ θεῶν ἀνάγκας θνητὸν ὄντα δεῖ φέρειν ; Eur.~fr.~37~Kannicht τὰς δὲ δαιμόνων τύχας / ὅστις φέρει κάλλιστ᾽ ἀνὴρ οὗτος σοφός ; Eur.~fr.~98~Kannicht ἀλλ᾽ εὖ φέρειν χρὴ συμφορὰς τὸν εὐγενῆ ; Aesch.~\textit{PV}~103–105 τὴν πεπρωμένην δὲ χρὴ / αἶσαν φέρειν ὡς ῥᾷστα, γιγνώσκονθ᾽ ὅτι / τὸ τῆς ἀνάγκης ἔστ᾽ ἀδήριτον σθένος ; Aesch.~\textit{Pers.}~293–294 ὅμως δ᾽ ἀνάγκη πημονὰς βροτοῖς φέρειν / θεῶν διδόντων ; Soph.~\textit{OC}~1694 τὸ φέρον ἐκ θεοῦ φέρειν ; \textit{idem} App. Prov.~4.98 (add. καλῶς ; add. χρή) ; Soph.~fr.~526~N.² =~585~R. (Stob. Flor.~108.58) ἀλγεινά, Πρόκνη, δῆλον· ἀλλ᾽ ὅμως χρεὼν / τὰ θεῖα θνητοὺς ὄντας εὐπετῶς φέρειν ; Archil.~fr.~9.9–10~Bergk⁴ ἀλλὰ τάχιστα / τλῆτε γυναικεῖον πένθος ἀπωσάμενοι ; Hor.~\textit{Carm.}~1.24.19–20 durum : sed levius fit patientia / quidquid corrigere est nefas ; Tosi n°~1205 Facile omnes … aegrotis damus (Ter. Andr. 309) ; ῥᾷον παραινεῖν ἢ παθόντα καρτερεῖν (Eur. Alc. 1078) ; Erasm.~\textit{Adag.}~1.3.14 Feras non culpes quod vitari non potest ; Antiph.~fr.~103~K.-A. (=~104~Kock ; Stob. Flor.~108.28~Meineke) ἀνδρὸς διαφέρει τοῦτ᾽ ἀνήρ· ὁ μὲν κακῶς / πράττων τὸ λυποῦν ἤγαγ᾽ εἰς παράστασιν, / ὁ δ᾽ ἐμφρόνως δεξάμενος ἤνεγκεν καλῶς ; Men.~fr.~855~K.-A. (=~812~Kock ; Stob. Flor.~108.5~Meineke) πειρῶ Τύχης ἄνοιαν ἀνδρείως φέρειν ; Men.~fr.~856~K.-A. (Stob.~4.44.6) ἄνδρα τὸν ἀληθῶς εὐγενῆ καὶ τἀγαθὰ / καὶ τὰ κακὰ δεῖ πταίοντα γενναίως φέρειν ; Eur.~\textit{Phoen.}~382 ἀτὰρ τί ταῦτα; δεῖ φέρειν τὰ τῶν θεῶν ; Eur.~fr.~505~Kannicht τὰ προσπεσόντα δ᾽ ὅστις εὖ φέρει βροτῶν, / ἄριστος εἶναι σωφρονεῖν τ᾽ ἐμοὶ δοκεῖ ; Aesch.~\textit{PV}~265–267 ἐλαφρὸν ὅστις πημάτων ἔξω πόδα / ἔχει παραινεῖν νουθετεῖν τε τὸν κακῶς / πράσσοντ᾽ ; Soph.~\textit{Phaedra} (Stob. Flor.~108.53) αἴσχη μέν, ὦ γυναῖκες, οὐδ᾽ ἂν εἷς φύγοι / βροτῶν ποθ᾽, ᾧ καὶ Ζεὺς ἐφορμήσῃ κακά. / νόσους δ᾽ ἀνάγκη τὰς θεηλάτους φέρειν ; Soph.~\textit{Phil.}~1316–1317 ἀνθρώποισι τὰς μὲν ἐκ θεῶν / τύχας δοθείσας ἔστ᾽ ἀναγκαῖον φέρειν ; Soph.~fr.~296~N.² =~319~R. (Orion Flor.~7.10) Πρὸς ἀνδρὸς ἐσθλοῦ πάντα γενναίως φέρειν ; \textit{Il.}~24.49 τλητὸν γὰρ Μοῖραι θυμὸν θέσαν ἀνθρώποισιν ; \textit{Od.}~18.134–135 ἀλλ᾽ ὅτε δὴ καὶ λυγρὰ θεοὶ μάκαρες τελέσωσι, / καὶ τὰ φέρει ἀεκαζόμενος τετληότι θυμῷ ; Pl.~\textit{Menex.}~247d καὶ φέροντες μὲν ἀνδρείως τὰς συμφορὰς ; Pl.~\textit{Menex.}~248c βαρέως φέροντες τὰς συμφοράς· κούφως δὲ καὶ μετρίως ; Isoc.~2.39 ἀλλὰ καλῶς καὶ μετρίως καὶ τὰς συμφορὰς καὶ τὰς εὐτυχίας φέρειν ἐπισταμένους ; Stob. Flor.~108 (titre) ὅτι δεῖ γενναίως φέρειν τὰ προσπίπτοντα ὄντας ἀνθρώπους καὶ κατ᾽ ἀρετὴν ζῆν ὀφείλοντας ; Lys.~3.4 οὗτος δὲ βέλτιστος ἂν εἴη καὶ σωφρονέστατος, ὅστις κοσμιώτατα τὰς συμφορὰς φέρειν δύναται ; Hsch.~ξ~114 ξυμφορά· συντυχία. καὶ ἐπὶ ἀγαθοῦ τάσσεται ⟨ὡς παῤ Αἰσχύλῳ ἐν Καβείροις (fr. 96a)⟩, καὶ ἐπὶ κακοῦ ⟨παρὰ Σοφοκλεῖ (fr. dub. 1151)⟩ ; Hsch.~σ~2356 συμφορά· συντυχία. σύμπτωμα. ἀτυχία ; Synag.~σ~303 συμφορά· ἀτυχία, κακῶν συντυχία ; \textit{idem} Suid.~σ~1407 ; Men.~fr.~236.14–16~K.-A. (Μισογύνης ; Stob.~4.44.35) οὕτω γὰρ οἴσεις πᾶν τὸ πρᾶγμ᾽· ἂν δ᾽ ἐκλέγῃ / ἀεὶ τὸ λυποῦν, μηδὲν ἀντιπαρατιθείς / τῶν προσδοκωμένων, ὀδυνήσῃ διὰ τέλους ; Ter.~\textit{Phorm.}~138 Quod fors feret feremus aequo animo ; D.~2.9 καὶ συμπονεῖν καὶ φέρειν τὰς συμφορὰς καὶ μένειν ἐθέλουσιν ἅνθρωποι ; Plu.~\textit{Cons.Apoll.}~102a τότε μὲν οὖν ὑπὸ τὸν τῆς τελευτῆς καιρὸν ἐντυγχάνειν σοι καὶ παρακαλεῖν ἀνθρωπίνως φέρειν τὸ συμβεβηκὸς ἀνοίκειον ἦν ; Phryn.~\textit{Ecl.}~397 Fischer Μετριάζειν· τοῦτο οἱ μὲν ἀρχαῖοι ἐπὶ τοῦ τὰ συμβαίνοντα μετρίως φέρειν τιθέασιν, Μένανδρος (fr. 881 Koe.) δ᾽ ἐπὶ τοῦ ἀσθενεῖν παρὰ τὴν τῶν δοκίμων χρῆσιν}
\lp{198 τοῖς τεχνάσμασιν]}{Ar.~\textit{Thesm.}~94 τοῦ γὰρ τεχνάζειν ἡμέτερος ὁ πυραμοῦς ; \textit{idem} Suid.~τ~439 (s.v. τεχνάζειν· τὸ πανούργως τι διαπράττεσθαι ; praem. v. 93 τὸ πρᾶγμα κομψὸν καὶ σφόδρ᾽ ἐκ τοῦ σοῦ τρόπου) ; Eur.~\textit{Or.}~1560 τεχνάσματ᾽ ἐστὶ ταῦτα καὶ πολὺς γέλως ; Poll.~9.133 (Bethe II, p.~183) ἀπάτη ἐξαπάτη, … σόφισμα, τέχνη τέχνασμα, γοητεία, ἐπιβουλή, δόλος, κλοπή, … τεχνάζειν, γοητεύειν, δολοῦν, κλέπτειν ; Ammon.~\textit{Diff.}~467 τεχνάσματα καὶ τεχνήματα διαφέρουσι. τεχνάσματα μὲν γὰρ τὰ κακουργήματα, τεχνήματα δὲ αἱ τέχναι ; Eur.~\textit{IT}~1355 δόλια τεχνήματ᾽, εἰχόμεσθα τῆς ξένης ; Aesch.~fr.~322~N.² κάπηλα προσφέρων τεχνήματα ; Soph.~\textit{Phil.}~928 δεινῆς τέχνημ᾽ ἔχθιστον, οἷά μ᾽ εἰργάσω ; X.~\textit{Hell.}~6.4.7 ὡς ταῦτα πάντα τεχνάσματ ; Σ Eur.~\textit{Or.}~1560 τεχνάσματ᾽ : εἰπὼν τεχνάσματα ἐπάγει ὡς ταυτὸν καὶ πολὺς γέλως … Gr. ὑποβολαὶ, πανουργίαι ; Hsch.~τ~722 τέχνη· ἐπιστήμη (Γ 61 . .). ἢ δόλος (δ 455 . .) ; Hsch.~τ~723 τεχνήματα· κατασκευάσματα ; Suid.~τ~441 Τέχνη: ὁ δόλος. ὅθεν καὶ ἀτεχνῶς, ἀντὶ τοῦ ἀληθῶς καὶ ἄνευ δόλου ; Ar.~\textit{Lys.}~398 τοιαῦτ᾽ ἀπ᾽ αὐτῶν ἐστιν ἀκολαστάσματα ; Eur.~\textit{Or.}~1053 καὶ μνῆμα δέξαιθ᾽ ἕν, κέδρου τεχνάσματα ; Hsch.~κ~2314 κέρδος· σύνεσις. τέχνασμα. ὠφέλεια (Κ 225) ; Hsch.~τ~721 τεχνάζω· τεχνάζομαι}
\lp{199 τοῖς παθήμασιν]}{Men.~fr.~179~Kock (Stob. Flor.~108.57) οὐδὲν πέπονθας δεινόν, ἂν μὴ προσποιῇ ; \textit{idem} Plu.~\textit{De\,tranq.\,an.}~475c (praem. τὸ τοῦ Μενάνδρου πρόχειρον), Diogenian.~7.38, Macar.~6.62 (CPG II, p.~196) (add. ἐπὶ τῶν ὑποκρινομένων δεινὰ πεπονθέναι) ; Pl.~\textit{Smp.}~196b–c οὔτε γὰρ αὐτὸς βίᾳ πάσχει, εἴ τι πάσχει—βία γὰρ Ἔρωτος οὐχ ἅπτεται· οὔτε ποιῶν ποιεῖ ; Σ Ar.~\textit{Thesm.}~199 ἐκφεύγειν ; Hsch.~π~38 παθήματα· συμφοραί (Ep. Rom. 8, 18) ; Hsch.~φ~302 φέρειν· λαμβάνειν. βαστάζειν. ὑπομένειν (Aesch. Pers. 293. Soph. OT 1211) ; Soph.~\textit{OC}~7 στέργειν γὰρ αἱ πάθαι με χὠ χρόνος ξυνὼν ; Soph.~\textit{Phil.}~538 ἐγὼ δ᾽ ἀνάγκῃ προύμαθον στέργειν κακά ; Thgn.~1029–1030 (Stob. Flor.~124.9) τόλμα, θυμέ, κακοῖσιν ὅμως ἄτλητα πεπονθώς· / δειλῶν τοι κραδίη γίγνεται ὀξυτέρη ; Gorg.~\textit{Hel.}~9 ἴδιόν τι πάθημα διὰ τῶν λόγων ἔπαθεν ἡ ψυχή ; Σ Soph.~\textit{OC}~7 πάθαι γὰρ αἱ συμφοραί ; Synag.~τ~233 τραγικὸν πάθος· συμφορᾶς μεστόν ; Eur.~\textit{Hipp.}~570 ὦ δυστάλαινα τῶν ἐμῶν παθημάτων ; Soph.~\textit{Trach.}~141–142 πεπυσμένη μέν, ὡς ἀπεικάσαι, πάρει / πάθημα τοὐμόν ; Th.~1.23.1 παθήματά τε ξυνηνέχθη γενέσθαι ἐν αὐτῷ τῇ Ἑλλάδι οἷα οὐχ ἕτερα ἐν ἴσῳ χρόνῳ ; Hsch.~τ~237 τὰ στίγματα· τὰ παθήματα, τὰ ἕλκη (Ep. Gal. 6, 17)}
\lp{199 τοῖς παθήμασιν (πάθει μάθος)]}{Aesch.~\textit{Ag.}~177 πάθει μάθος ; Hdt.~1.207.1 τὰ δέ μοι παθήματα ἐόντα ἀχάριτα μαθήματα γέγονε ; Mant. Prov.~2.92 (CPG II, p.~772) Τὰ παθήματα τοῖς ἀνθρώποις μαθήματα γίνεται ; \textit{idem} Tosi n°~1201 (praem. πολλάκις ; γίγνονται) ; Soph.~\textit{OT}~403 εἶναι, παθὼν ἔγνως ἂν οἷά περ φρονεῖς ; Hes.~\textit{Op.}~218 ἐς τέλος ἐξελθοῦσα· παθὼν δέ τε νήπιος ἔγνω ; Pl.~\textit{Smp.}~222b ἃ δὴ καὶ σοὶ λέγω, ὦ Ἀγάθων, μὴ ἐξαπατᾶσθαι ὑπὸ τούτου, ἀλλ᾽ ἀπὸ τῶν ἡμετέρων παθημάτων γνόντα εὐλαβηθῆναι, καὶ μὴ κατὰ τὴν παροιμίαν ὥσπερ νήπιον παθόντα γνῶναι ; Diogenian.~2.31 Ἁλιεὺς πληγεὶς νοῦν οἴσει: … ὅμοιον δέ ἐστι τῷ παρ᾽ Ἡσιόδῳ, Παθὼν δέ τε νήπιος ἔγνω ; Erasm.~\textit{Adag.}~1.1.31 (LB II, col.~38 F –~39 E) Malo accepto stultus sapit. … Plato in Symposio … i. e. Quae sane tibi quoque dico, Agathon, ne ab hoc fallaris … nec juxta proverbium veluti stultus post acceptum malum sapias. … Rursum hoc pacto: ἐξ ὧν ἔπαθες, ἔμαθες. i. e. ut sententiam magis quam verba reddam, Quae nocent, docent}
\lp{200–201 καὶ μὴν σύ γ᾽ … ἀλλὰ τοῖς παθήμασιν]}{\textit{Test.} Suid.~τ~440 καὶ μὴν σύ γ᾽, ὦ καταπῦγον, εὐρύπρωκτος εἶ, οὐ τοῖς λόγοισιν, ἀλλὰ τοῖς παθήμασιν}
\lp{200 καὶ μὴν σύ γ᾽]}{Ar.~\textit{Nub.}~909 καταπύγων εἶ κἀναίσχυντος ; Ar.~\textit{Thesm.}~35 καὶ μὴν βεβίνηκας σύ γ᾽, ἀλλ᾽ οὐκ οἶσθ᾽ ἴσως ; Aesch.~\textit{PV}~982 καὶ μὴν σύ γ᾽ οὔπω σωφρονεῖν ἐπίστασαι ; Soph.~\textit{Aj.}~531 καὶ μὴν φόβοισί γ᾽ αὐτὸν ἐξελυσάμην}
\lp{200 ὦ κατάπυγον]}{Ar.~\textit{Ach.}~79 ἡμεῖς δὲ λαικαστάς τε καὶ καταπύγονας ; Ar.~\textit{Vesp.}~687 ὅταν εἰσελθὸν μειράκιόν σοι κατάπυγον, Χαιρέου υἱός ; Ar.~\textit{Lys.}~137 ὦ παγκατάπυγον θἠμέτερον ἅπαν γένος ; Ar.~fr.~128~K.-A. (=~130~Kock ; Poll.~6.69 ; D.L.~4.18–19) καταπυγοσύνη ταῦτ᾽ ἐστὶ πρὸς κρέας μέγα ; \textit{idem} Suid.~ο~76 (praem. ὀζωτά· Ἀριστοφάνης περὶ Εὐριπίδου φησίν· ὀζωτὰ καὶ στιλφωτά· ἅπερ, ὡς αὐτός φησι) ; Sophr.~fr.~63~K.-A. (=~63 Kaibel ; Ath.~7.281e) τὸ καταπυγοτέραν τ᾽ ἀλφηστᾶν … ἀπὸ τοῦ οὖν κατὰ τὴν πυγὴν θατέρῳ τὸν ἕτερον ἀκολουθεῖν τῶν ἀρχαίων τινὲς τοὺς ἀκρατεῖς καὶ καταφερεῖς οὕτω καλοῦσιν ; Luc.~\textit{Ind.}~23 ἀλλ᾽, ὦ κατάπυγον, οἴει τοσοῦτον μανδραγόραν κατακεχύσθαι αὐτοῦ ; Luc.~\textit{Tim.}~22 ἢ καταπύγων οἰκέτης ἐκ παιδικῶν τίμιος, ὑπεξυρημένος ἔτι τὴν γνάθον ; Σ Ar.~\textit{Eq.}~639 καταπύγων δὲ, εὐρύπρωκτος ; Hsch.~κ~1365 καταπύγου· κιναίδου, ἤγουν ἀσελγοῦς ; Hsch.~κ~1366 καταπύγων· ὁ μέσος δάκτυλος. καὶ ὁ κατωφερής ; Poll.~6.126–127 (Bethe II, p.~35) Κίναιδος, πόρνος, ἀκάθαρτος, … καταπύγων, θηλυδρίας, γυναικίας, … ἀνδρόγυνος, … ἐκτεθηλυσμένος ἢ μεμαλαγμένος, λελυγισμένος, γύννις, θῆλυς τὴν ψυχήν ; Phot.~κ~343 Καταπύγων: ὁ λάγνης· ἡ δὲ μεταφορὰ ἀπὸ τῶν ἀλφηστῶν οὕτως λεγομένων ἰχθύων, ὅτι ἕπονται κατ᾽ οὐράν ; Luc.~\textit{Ind.}~23 (proverbe) καὶ ὅλως κατὰ τὴν παροιμίαν, θᾶττον ἂν πέντε ἐλέφαντας ὑπὸ μάλης κρύψειας ἢ ἕνα κίναιδον ; \textit{idem} Apostol.~8.82 (CPG II, p.~454) (πέντ᾽ ἐλάφους pro ἂν πέντε ἐλέφαντας ; add. ἐπὶ τῶν ἀναισχύντων) ; Ar.~\textit{Ach.}~664 δειλὸς καὶ λακαταπύγων ; Ar.~\textit{Nub.}~529 ὁ σώφρων τε χὠ καταπύγων ἄριστ᾽ ἠκουσάτην ; Ar.~\textit{Vesp.}~84 ἐπεὶ καταπύγων ἐστὶν ὅ γε Φιλόξενος ; Cratin.~fr.~58.1–2~K.-A. (=~53~Kock ; Ael. NA~12.10) φέρε νῦν σοι ἐξ αἰθρίας καταπυγοσύνην μυὸς ἀστράψω Ξενοφῶντος ; Cratin.~fr. *259~K.-A. (=~241~Kock ; Plu. Per.~24.6) Ἥραν τέ οἱ Ἀσπασίαν τίκτει Καταπυγοσύνη παλλακὴν κυνώπιδα ; Luc.~\textit{Lexiph.}~12 τὸν καταπύγονα καὶ λακκοσχέαν ; Alciphr.~3.9 (III~45) τὸν καταπύγονα καὶ θηλυδρίαν ; Σ Ar.~\textit{Ach.}~716 Ἀλκιβιάδην τὸν Κλεινίου ὡς καταπύγονα κωμῳδοῦσιν ; Σ Ar.~\textit{Lys.}~776 καταπυγωνέστερον : Μαλακώτερον καὶ πορνικώτερον τῶν χελιδόνων, τουτέστι τῶν γυναικῶν ; Σ Ar.~\textit{Vesp.}~684 κατάπυγον δὲ πεπορνευμένον ; Hsch.~β~317 Βάταλος· καταπύγων καὶ ἀνδρόγυνος, κίναιδος. ἔκλυτος (Dem. 18, 180) ; Hsch.~κ~1364 καταπυγοσύνη· ἡδονὴ μεγάλη ; Poll.~2.184 (Dindorf) Ἀπὸ δὲ πυγῶν, εὔπυγος, καλλίπυγος, καταπύγων ; Hdn.~\textit{Περὶ\,κλίσεως\,ὀνομάτων}~2.208 πυγή κατάπυγοϲ καταπύγων καταπύγονοϲ (ὁ κίναιδοϲ) ; Phot. s.v.~καταπυγίζειν (Porson, p.~122) Καταπυγίζειν: τὸ τὴν πυγὴν ἐπιπολὺ μεταφέρειν ἐν τῷ βαδίζειν ; Suid.~α~2684 Ἀντιμάχου καταπυγοσύνης ἀναπλήσει. ἀντὶ τοῦ ἡ μαλακία. οὗτος εἰς θηλύτητα καὶ εὐμορφίαν κωμῳδεῖται ; Suid.~γ~467 Γρύττος· τῶν ἐπὶ μαλακίᾳ διαβαλλομένων οὗτος ἦν. οἱ δὲ Γρῦπον λέγουσι διὰ τοῦ π· ὃς ἦν καταπύγων. ἤτοι λαγνός ; Suid.~κ~736 Καταπύγους: κιναίδους ; Suid.~κ~738 Καταπύγων: ὁ λάγνης ; Suid.~κ~739 Καταπυγωνέστερον: μαλακώτερον, πορνικώτερον. οὐκέτι δόξειέ τι ἄλλο τι καταπυγωνέστερον εἶναι ; Diogenian.~3.87 Γέλως Ἰωνικός: ἐπὶ τῶν κιναίδων. Παρόσον οἱ Ἴωνες ἐπὶ τούτῳ διαβάλλονται ; Ar.~\textit{Eq.}~639 ἐκ δεξιᾶς ἀπέπαρδε καταπύγων ἀνήρ ; Ar.~\textit{Nub.}~1023 καταπυγοσύνης ἀναπλήσει ; Ar.~\textit{Lys.}~776 ὄρνεον οὐδ᾽ ὁτιοῦν καταπυγωνέστερον εἶναι ; Hsch.~ζ~163 ζιτᾶνα· καταπύγονα ; Hsch.~κ~4641 κυπάται· κίναιδοι, μαλακοί ; Hsch.~ο~68 ὀδαχᾶς· καταπύγων. Ταραντῖνοι (gl. Ital. 140 K.–A.) ; Phryn.~\textit{Ecl.}~406 Fischer Καταφερής· ἐπὶ τῶν πρὸς ἀφροδίσια ἀκολάστων λέγουσιν οἱ πολλοί, οὐδαμῶς οὕτω τῶν δοκίμων χρωμένων}
\lp{200 ὦ κατάπυγον (Παυσανίας καὶ Ἀγάθων)]}{Pl.~\textit{Smp.}~193b–c καὶ μή μοι ὑπολάβῃ Ἐρυξίμαχος, κωμῳδῶν τὸν λόγον, ὡς Παυσανίαν καὶ Ἀγάθωνα λέγω—ἴσως μὲν γὰρ καὶ οὗτοι τούτων τυγχάνουσιν ὄντες καί εἰσιν ἀμφότεροι τὴν φύσιν ἄρρενες ; X.~\textit{Symp.}~8.32 καίτοι Παυσανίας γε ὁ Ἀγάθωνος τοῦ ποιητοῦ ἐραστὴς ἀπολογούμενος ὑπὲρ τῶν ἀκρασίᾳ ἐγκαλινδουμένων ; Lib.~\textit{Or.}~64.83 καὶ ὑπὸ τούτων ἐγένοντο θηλυδρίαι Κλεισθένης, Φιλόξενος, Ἀμυνίας, Ἀριστόδημος, Βάθυλλος, Διόγνητος, Δῆμος, Κλεινίας, Ἀγάθων}
\lp{200 εὐρύπρωκτος εἶ]}{Ar.~\textit{Ach.}~716 τοῖς νέοισι δ᾽ εὐρύπρωκτος καὶ λάλος χὠ Κλεινίου ; Ar.~\textit{Nub.}~1084–1085 ἕξει τινὰ γνώμην λέγειν τὸ μὴ εὐρύπρωκτος εἶναι; / ἢν δ᾽ εὐρύπρωκτος ᾖ, τί πείσεται κακόν ; Ar.~\textit{Nub.}~1090–1094 ἐξ εὐρυπρώκτων ; Ar.~\textit{Nub.}~1330 ὦ λακκόπρωκτε. — πάττε πολλοῖς τοῖς ῥόδοις ; Ar.~\textit{Vesp.}~1070 σχῆμα κεὐρυπρωκτίαν ; Eub.~fr.~118.7~K.-A. (Ath.~1.25c) οἵτινες πόλιν μίαν λαβόντες εὐρυπρωκτότεροι πολὺ τῆς πόλεος ἀπεχώρησαν ἧς εἷλον τότε ; Cephisod.~fr.~3.4~K.-A. (=~3~Kock ; Ath.~15.689f) (Ξα.) ὦ λακκόπρωκτε, βάκχαριν τοῖς σοῖς ποσὶν ; Hsch.~α~7248 Ἀριστόδημος· Ἀριστόδημον οἱ κωμικοὶ τὸν πρωκτόν (Cratin. fr. 160 K.–A. Ar. fr. 242 K.–A.) καὶ Θεόδωρον (com. ad. fr. 351 K.–A.) καὶ Τιμησιάνακτα (com. ad. fr. 283 K.–A.) ἔλεγον, ἀπὸ τῶν ἡταιρηκότων ; Hsch.~κ~4739 κυσοχήνη· εἶδος δεσμοῦ. οἱ δὲ εὐρυπρωκτίαν. οἱ δὲ ξύλον [ἐν] ἐν ᾧ ἁμαρτάνουσαι αἱ πόρναι ἐδεσμεύοντο (Hippon. fr. 82, 2 W. = 83, 2 Degani) ; \textit{idem} Phot.~κ~1264 (εὐρυπρωκτία tantum) ; Hsch.~π~1089 πασχητιᾷ· πάσχειν θέλει [ἢ] ἀκόλαστος, ἢ αἰσχρᾶς ἡδονῆς ἡττᾶται ; Ar.~\textit{Ach.}~104 οὐ λῆψι χρῦσο χαυνόπρωκτ᾽ Ἰαοναῦ ; Ar.~\textit{Ach.}~843 οὐδ᾽ ἐξομόρξεται Πρέπις τὴν εὐρυπρωκτίαν σοι ; \textit{idem} Suid.~ε~1121 (ἐναπομόρξεται Πρέμις ; add. τουτέστι τὴν μαλακίαν ἐναποψήσεται, ἀπομάξει. καταπύγων δὲ ἐκωμῳδεῖτο ὁ Πρέμις), Zonar. (Tittmann) col.~911 (s.v. εὐρυπρωκτίαν· μαλακίαν ; τὴν εὐρυπρωκτίαν σοι ἐναποψήσεται, ὅ ἐστιν ἀπομάξει) ; Ar.~\textit{Nub.}~1099 τοὺς εὐρυπρώκτους· τουτονὶ ; Eup.~fr.~385.3–4~K.-A. (=~351~Kock ; Ath.~1.17d) (Ἀλκ.) ⏑ – ⏑ ὃς δὲ πρῶτος ἐξηῦρον τὸ πρῲ ᾽πιπίνειν / (Β.) πολλήν γε λακκοπρωκτίαν ἡμῖν ἐπίστασ᾽ εὑρών ; Σ Luc.~42.4 Rabe ὁ Ἀριστόδημος δὲ μιαρὸς καὶ καταπύγων ἐς ὑπερβολήν, ἀφ᾽ οὗ καὶ ὁ πρωκτὸς Ἀριστόδημος καλεῖται ; Hsch.~β~467 βινητιᾶν· τὸ πασχητιᾶν. καὶ τὸ ὄρεξιν ἐπὶ συνουσίαν ἔχειν (Plat. com. fr. 188, 21 K.–A.) ; Hsch.~ε~3839 Ἐξήκεστος· ἡταιρηκώς. ὅθεν καὶ τοὺς πρώκτους ὁμωνύμως Ἐξηκέστους ἔλεγον (com. ad. fr. 337 K.–A.) ; Hsch.~θ~267 Θεόδωρος· Θεοδώρους ἔλεγον οἱ κωμικοὶ τοὺς πρώκτους, ἀπὸ Θεοδώρου τινὸς οὐκ εὖ τῇ ἑαυτοῦ ὥρᾳ χρησαμένου (com. ad. fr. 351 K.–A.) ; Hsch.~λ~209 λακκοσκάπερδον· λακκόπρωκτον (com. ad. fr. 514 K.–A.) ; Suid.~λ~59 Λακκόπρωκτος· Ἀριστοφάνης· ὦ λακκόπρωκτε ; Zen.~1.52 Ἀκεσίας ἰάσατο: ἐπὶ τῶν ἐπὶ τὸ χεῖρον ἰωμένων. Ὅλην δὲ Ἀριστοφάνης ἐν τετραμέτροις † ἐκφέρει λέγων, Ἀκεσίας τὸν πρωκτὸν ἰάσατο ; Hsch.~ε~6767 «εὕδοντι δ᾽ αἱρεῖ πρωκτός»· παροιμία ἀπὸ τῆς εὕδοντι κύρτὸς αἱρεῖ ; Hsch.~π~4144 πρωκτὸς λουτροῦ περιγίνῃ· παροιμία, ὅταν τις μὴ δύνηται ἀπονίψασθαι}
\lp{201 οὐ τοῖς λόγοισιν ἀλλὰ τοῖς παθήμασιν]}{Eur.~\textit{El.}~893 ἥκω γὰρ οὐ λόγοισιν ἀλλ᾽ ἔργοις κτανὼν ; Soph.~\textit{OT}~516–517 ταῖς νῦν νομίζει πρός γ᾽ ἐμοῦ πεπονθέναι / λόγοισιν εἴτ᾽ ἔργοισιν εἰς βλάβην φέρον ; Tosi n°~2079 Facta non verba ; Soph.~\textit{OC}~782 λόγῳ μὲν ἐσθλά. τοῖσι δ᾽ ἔργοισιν κακά ; Soph.~\textit{OC}~1143–1144 οὐ γὰρ λόγοισι τὸν βίον σπουδάζομεν / λαμπρὸν ποεῖσθαι μᾶλλον ἢ τοῖς δρωμένοις ; Th.~6.18.6 οὐ λόγῳ ἀλλ᾽ ἔργῳ ; \textit{idem} Pl.~\textit{Ap.}~32d ; Din.~1.32 χρήσιμος ὑμῖν οὐ λόγοις ἀλλ᾽ ἔργοις βουλόμενος γενέσθαι ; Democr.~fr.~145~D.-K. (Plu.~\textit{Mor.}~9f) λόγος γὰρ ἔργου σκιή κατὰ Δημόκριτον ; \textit{idem} D.L.~9.37 (τούτου ἐστὶ καὶ τὸ λόγος ἔργου σκιή)}
\lp{201 τοῖς παθήμασιν (πάσχειν)]}{Ar.~\textit{Nub.}~1085 ἢν δ᾽ εὐρύπρωκτος ᾖ, τί πείσεται κακόν ; Ar.~\textit{Thesm.}~59–60 ὃς ἕτοιμος σοῦ τοῦ τε ποιητοῦ / τοῦ καλλιεποῦς κατὰ τοῦ θριγκοῦ ; Luc.~\textit{Gall.}~32 καταπυγοσύνην καὶ πασχητιασμόν τινα καὶ ἀσέλγειαν ; Luc.~\textit{Ind.}~23 ὡς χρὴ λεαίνεσθαι καὶ παρατίλλεσθαι καὶ πάσχειν καὶ ποιεῖν ἐκεῖνα ; Hsch.~ε~5373 ἐπὶ Φαληνίου· τὸν Ἀλκιβιάδην φησὶν ὁ Ἀριστοφάνης (fr. 244 K.–A.) ἐπὶ Φαληνίου γεγενῆσθαι, σκώπτων παρὰ τὸν φάλητα· ἐπασχητία γὰρ ; Suid.~π~767 Πασχητιᾷ: πάσχειν ἐθέλει ; Ar.~\textit{Thesm.}~843–845 καὶ τόκον πράττοιτο, διδόναι μηδέν᾽ ἀνθρώπων τόκον, … ἀξία γοῦν εἶ τόκου τεκοῦσα τοιοῦτον τόκον ; Arist.~\textit{EN}~1148b31–33 ὥσπερ οὐδὲ τὰς γυναῖκας, ὅτι οὐκ ὀπύουσιν ἀλλ᾽ ὀπύονται ; [Luc.]~\textit{Am.}~26 αὐτὸς ἔμοιγε δοκεῖ πασχητιᾶν ἀμφίβολον Ἀφροδίτην μεταδιώκων ; Ath.~5.187c καὶ τὸν Ἀλκιβιάδην δὲ παράγει λέγοντα ὅτι πασχητιᾷ ; D.L.~6.47 παῦσαι γάρ, ἔφη, καὶ σὺ τὰ δείγματα τοῦ πασχητιῶντος περιφέρων ; Hsch.~κ~4732 κυσιᾶι· πασχητιᾶι (com. ad. fr. 378 K.–A.) ; Hdn.~\textit{De\,pros.\,cath.}~20 οὐκ ἀγνοῶ δὲ ὡς ἔσθ᾽ ὅτε οἱ κωμικοὶ καὶ ἀρσενικῷ γένει τὴν σύνταξιν ποιοῦνται ἀλλ᾽ εἰκὸς αὐτοὺς θηλυκῇ προσηγορίᾳ σκώπτειν τοὺς πασχητιῶντας ; Suid.~ε~2021 τὸν δὲ Κλεισθένην καὶ ἐπὶ τοῦ πασχητιᾶν κωμῳδοῦσι ; Ar.~\textit{Thesm.}~232 οἴμοι κακοδαίμων, ψιλὸς αὖ στρατεύσομαι ; Hsch.~κ~2444 †κεωρεῖν· πασχητιᾶν ; Hsch.~ο~435 οἰφόλις· γυνὴ καταφερής, μαχλός, πασχητιῶσα}
\lp{202 τί δ᾽ ἔστιν ὅτι δέδοικας]}{Eur.~\textit{Or.}~1520 μὴ πέτρος γένῃ δέδοικας ὥστε Γοργόν᾽ εἰσιδών ; Ar.~\textit{Vesp.}~628–629 καὶ σὺ δέδοικάς με μάλιστ᾽ αὐτός· / νὴ τὴν Δήμητρα δέδοικας ; Ar.~\textit{Lys.}~830 τί δ᾽ ἔστιν; εἰπέ μοι τίς ἡ βοή ; Ar.~\textit{Thesm.}~1186 αὔλει σὺ θᾶττον· ἔτι δέδοικας τὸν Σκύθην ; Eur.~\textit{Ba.}~1029 τί δ᾽ ἔστιν; ἐκ βακχῶν τι μηνύεις νέον ; Soph.~\textit{Trach.}~457 κεἰ μὲν δέδοικας, οὐ καλῶς ταρβεῖς ; Hsch.~δ~400 δέδοικεν· ἐφοβήθη (Eur. Phoen. 959 . .)}
\lp{202 ἐλθεῖν αὐτόσε]}{Ar.~\textit{Lys.}~873 μὰ Δί᾽ ἐγὼ μὲν αὐτόσ᾽ οὔ ; Antiph.~fr.~57.19–20~K.-A. (=~55~Kock ; Ath.~15.667b) ἀλλ᾽ ἐγὼ μὲν σφενδόνῃ οὐκ ἂν ἐφικοίμην αὐτόσ᾽ ; Metag.~fr.~6.3–4~K.-A. (=~6~Kock ; Θουριοπέρσαι ; Ath.~6.269f) ὁ δ᾽ ἕτερος ὠθεῖ κῦμα ναστῶν καὶ κρεῶν ἑφθῶν τε βατίδων εἰλυομένων αὐτόσε ; Pl.~\textit{Phd.}~68b ἀγανακτήσει τε ἀποθνῄσκων καὶ οὐχ ἅσμενος εἶσιν αὐτόσε ; Hsch.~α~8461 αὐτόσε· ἐκεῖσε ; \textit{idem} Suid.~α~4522 (add. λέγεται δὲ καὶ αὐτοῦσε) ; Th.~4.1.4 καὶ ἄλλαι αἱ πληρούμεναι ἔμελλον αὐτόσε ἐγκαθορμισάμεναι τὸν πόλεμον ἐντεῦθεν ποιήσεσθαι ; X.~\textit{An.}~4.7.2 συνεληλυθότες δ᾽ ἦσαν αὐτόσε καὶ ἄνδρες καὶ γυναῖκες καὶ κτήνη πολλά}
\lp{203 κάκιον ἀπολοίμην ἂν ἢ σύ]}{Ar.~\textit{Ach.}~476 κάκιστ᾽ ἀπολοίμην, εἴ τί σ᾽ αἰτήσαιμ᾽ ἔτι ; Ar.~\textit{Ach.}~756 ὅπως τάχιστα καὶ κάκιστ᾽ ἀπολοίμεθα ; Ar.~\textit{Thesm.}~757 κακῶς ἀπόλοι᾽, ὡς φθονερὸς εἶ καὶ δυσμενής ; Ar.~\textit{Ran.}~579 κάκιστ᾽ ἀπολοίμην, Ξανθίαν εἰ μὴ φιλῶ ; Σ Ar.~\textit{Thesm.}~203 ἐὰν τοῦτο ποιήσωμεν ; Eub.~fr.~115.1–2~K.-A. (=~116~Kock ; Ath.~13.559b) κακὸς κακῶς ἀπόλοιθ᾽ ὅστις γυναῖκα δεύτερος ἔγημε ; Hsch.~α~6780 ἀποφθίμην· ἀπολοίμην. ἀποθάνοιμι ; Hsch.~κ~310 κάκιον· εὐτελέστερον}
\lp{203 πῶς; — ὅπως;]}{Ar.~\textit{Nub.}~664 πῶς δὴ φέρ᾽; — ὅπως; ἀλεκτρυὼν κἀλεκτρυών ; Ar.~\textit{Vesp.}~48 πῶς; — ὅπως ; \textit{idem} Ar.~\textit{Plu.}~139 (add. οὐκ ἔσθ᾽ ὅπως) ; Ar.~\textit{Ach.}~594–595 ἀλλὰ τίς γὰρ εἶ; / ὅστις; πολίτης χρηστός ; Ar.~\textit{Eq.}~128 καὶ πῶς; — ὅπως; ὁ χρησμὸς ἄντικρυς λέγει ; Ar.~\textit{Eq.}~1073 πῶς δὴ τριήρης ἐστὶ κυναλώπηξ; — ὅπως ; Ar.~\textit{Nub.}~677 ἀτὰρ τὸ λοιπὸν πῶς με χρὴ καλεῖν; — ὅπως ; Ar.~\textit{Nub.}~689–690 πῶς ἂν καλέσειας ἐντυχὼν Ἀμυνίᾳ; / ὅπως ἄν; ὡδί, δεῦρο δεῦρ᾽ Ἀμυνία}
\lp{204 γυναικῶν ἔργα]}{Ar.~\textit{Lys.}~708 κακῶν γυναικῶν ἔργα καὶ θήλεια φρὴν ; \textit{Il.}~6.289 πέπλοι παμποίκιλα ἔργα γυναικῶν ; Hes.~\textit{Op.}~521 οὔ πω ἔργα ἰδυῖα πολυχρύσου Ἀφροδίτης ; \textit{idem} \textit{h.Ven.}~1 (μοῦσά μοι ἐννέπε ἔργα πολυχρύσου Ἀφροδίτης) ; Sol.~fr.~26~Bergk⁴ (Plu.~\textit{Sol.}~31.3) ἔργα δὲ Κυπρογενοῦς νῦν μοι φίλα καὶ Διονύσου ; Manetho~3.395 πάσχουσ᾽ αἴσχεα πολλὰ γυναικῶν ἔργα τελοῦντες ; Ptol.~\textit{Tetr.}~3.14 οἱ δὲ ἄνδρες μαλακοί τε καὶ σαθροὶ πρὸς τὰς παρὰ φύσιν συνουσίας καὶ γυναικῶν ἔργα ; Vett. Val.~2.16.12 ἑσπέριοι δὲ θηλύνουσι τοὺς ἄνδρας· ἄλλοτε μὲν γὰρ ἀνδράσι συγκοιμώμενοι γυναικῶν ἔργα ἐπιτελοῦσι ; Longus~4.19.5 παροίνημα δὲ Γνάθωνος οὐ δύναμαι περιιδεῖν γενόμενον, ὃς ἐς Μυτιλήνην αὐτὸν ἄγειν ἐπὶ γυναικῶν ἔργα σπουδάζει ; Greg. Cypr.~3.36 Πολλαὶ κυνὸς ἄῤῥενος εὐναί: ἐπὶ τῶν κατωφερῶν εἰς τὰ τὴς Ἀφροδίτης ἔργα ; \textit{Il.}~5.429 ἀλλὰ σύ γ᾽ ἱμερόεντα μετέρχεο ἔργα γάμοιο ; \textit{Od.}~11.246 αὐτὰρ ἐπεί ῥ᾽ ἐτέλεσσε θεὸς φιλοτήσια ἔργα ; Theoc.~\textit{Ep.}~4.4 παιδογόνῳ δυνατὸν Κύπριδος ἔργα τελεῖν ; Q.S.~1.652 καὶ νόος ἐξορόθυνε γυναικῶν ἔργα λιποῦσαν ; Suid.~μ~1467 ἀκμαίη πρὸς ἔρωτα καὶ ἡδέα Κύπριδος ἔργα}
\lp{204 γυναικῶν ἔργα νυκτερείσια]}{Ar.~\textit{Thesm.}~280–281 ὦ Θρᾷττα θέασαι, καομένων τῶν λαμπάδων / ὅσον τὸ χρῆμ᾽ ἀνέρχεθ᾽ ὑπὸ τῆς λιγνύος}
\lp{204 νυκτερείσια]}{Eur.~\textit{Ba.}~485–487 τὰ δ᾽ ἱερὰ νύκτωρ ἢ μεθ᾽ ἡμέραν τελεῖς; / νύκτωρ τὰ πολλά· σεμνότητ᾽ ἔχει σκότος. / τοῦτ᾽ ἐς γυναῖκας δόλιόν ἐστι καὶ σαθρόν ; Eur.~\textit{Hec.}~831–832 [ἐκ τοῦ σκότου τε τῶν τε νυκτερησίων] / [φίλτρων μεγίστη γίγνεται βροτοῖς χάρις.] ; Aesch.~\textit{Ag.}~22 ὦ χαῖρε λαμπτὴρ νυκτός, ἡμερήσιον ; Luc.~\textit{Alex.}~53 χρησμὸς νυκτερήσιος ; S.E.~\textit{M.}~10.188 οὔτε ἡμερήσιον οὔτε νυκτερήσιον φάντασμα ; Ar.~\textit{Ach.}~866 ἐπὶ τὴν θύραν μοι Χαιριδῆς βομβαύλιοι ; Ar.~\textit{Nub.}~166 ὦ τρισμακάριος τοῦ διεντερεύματος ; Ar.~\textit{Av.}~299–300 κειρύλος. / κειρύλος γάρ ἐστιν ὄρνις; — οὐ γάρ ἐστι Σποργίλος ; Eur.~\textit{Or.}~271 βεβλήσεταί τις θεῶν βροτησίᾳ χερί ; Σ Luc.~\textit{Alex.}~53 χρησμὸς νυκτερείσιος* ; Hdt.~4.101.3 ἡ δὲ ὁδὸς ἡ ἡμερησίη ; Hdt.~8.98.1 τοσοῦτοι ἵπποι τε καὶ ἄνδρες διεστᾶσι κατὰ ἡμερησίην ὁδὸν ἑκάστην ; Hsch.~ν~696 νυκτερεύειν· ἀγρυπνεῖν. Ἀττικοί ; Hsch.~ν~709 νύκτωρ· νυκτός (Eur. Bacch. 486 . .)}
\lp{204 νυκτερείσια (ἐρείδειν)]}{Ar.~\textit{Thesm.}~488 ἐξῆλθον ὡς τὸν μοιχόν· εἶτ᾽ ἠρειδόμην ; Ar.~\textit{Eccl.}~616 ἐπὶ τὴν ὡραιοτάτην αὐτῶν καὶ ζητήσουσιν ἐρείδειν ; Ar.~fr.~715~K.-A. (=~695~Kock ; Ath.~2.48c) ὅστις ἐν ἡδυόσμοις στρώμασι παννυχίζων τὴν δέσποιναν ἐρείδεις ; Phot.~ε~1903 (Theodoridis II, p.~182) ἐρείδειν· τὸ ἐσθίειν σφοδρῶς ἢ ἄλλο τι ἐνεργεῖν συντόνως. Ταγηνισταῖς (fr. 509 K.-A.)· „ἐρείδετον· κἀγὼ κατόπιν σφῷν ἕψομαι“ ; Ar.~\textit{Pax}~31 ἔρειδε, μὴ παύσαιο μηδέποτ᾽ ἐσθίων ; Ar.~fr.~76~K.-A. (=~74~Kock ; Suid.~σ~436) Ἀριστοφάνης ἐν Βαβυλωνίοις· μέσην ἔρειδε πρὸς τὸ σιμόν ; Hsch.~ε~5705 ἔρειδε· συντόμως ἔσθιε. ἡ μεταφορὰ ἀπὸ τῶν ἐρεσσόντων}
\lp{205 κλέπτειν ὑφαρπάζειν τε θήλειαν Κύπριν]}{Ar.~\textit{Eccl.}~722 τὴν τῶν ἐλευθέρων ὑφαρπάζειν Κύπριν ; Eub.~fr.~67.8–9~K.-A. (Ath.~13.568f–569a) καὶ μὴ λαθραίαν Κύπριν, αἰσχίστην νόσων πασῶν, διώκειν ; Eur.~\textit{Ba.}~459 τὴν Ἀφροδίτην καλλονῇ θηρώμενος ; Eur.~\textit{Ba.}~688 θηρᾶν καθ᾽ ὕλην Κύπριν ἠρημωμένας ; Soph.~\textit{El.}~114 αἳ τοὺς εὐνὰς ὑποκλεπτομένους ; Epict.~2.4.10 οὕτως καὶ αἱ γυναῖκες φύσει κοιναί. ὅταν δ᾽ ὁ νομοθέτης ὡς ἑστιάτωρ διέλῃ αὐτάς, οὐ θέλεις καὶ αὐτὸς ἴδιον μέρος ζητεῖν, ἀλλὰ τὸ ἀλλότριον ὑφαρπάζεις καὶ λιχνεύεις ; Alciphr.~4.13 μηδὲ αἰδουμένως τῆς ἀφροδίτης παρακλέπτειν ; Hsch.~κ~2938 κλέψαι· κλέπτειν. τὸ τῶν ἀλλότριων ἐπιθυμεῖν ; Hsch.~κ~4653 Κύπρις· πόρνη ; Eur.~\textit{Cret.} fr.~82.7~Austin λαθραίαν ἐμπολωμένη Κύπριν ; Eur.~\textit{El.}~364 πρὸς θεῶν, ὅδ᾽ ἁνὴρ ὃς συνεκκλέπτει γάμους ; Hsch.~κ~2928 κλέπτει· πονηρεύεται. ἐπιθυμεῖ. ἐξαπατᾶται ; Hdt.~1.105.4 τοῖσι δὲ τῶν Σκυθέων συλήσασι τὸ ἱρὸν τὸ ἐν Ἀσκάλωνι καὶ τοῖσι τούτων αἰεὶ ἐκγόνοισι ἐνέσκηψε ὁ θεὸς θήλεαν νοῦσον}
\lp{205 κλέπτειν]}{\textit{AP}~12.21 (Strato) Κλέψομεν ἄχρι τίνος τὰ φιλήματα, καὶ τὰ λαθραῖα ; Lyc.~1116 ὡς κλεψίνυμφον κοὐ δορίκτητον γέρας ; Ar.\,Byz.~\textit{Epit.}~2.272 Αἰθίοπι δὲ τὴν εὐνὴν κλέπτουσαν ; Hsch.~κ~3039 κλοτοπεύειν· παραλογίζεσθαι. ἀπατᾶν. κλεψιγαμεῖν ; Suid.~κ~1737 Κλέπτειν· αἰτιατικῇ. τὸ τὰ ἀλλότρια λάθρα ἀφαιρεῖσθαι, τὸ παραλογίζεσθαι ; Suid.~ν~532 Νοθεία: ἡ κλεψιγαμία}
\lp{205 ὑφαρπάζειν]}{Ar.~\textit{Eccl.}~921–922 ἀλλ᾽ οὐκ ἄν ποθ᾽ ὑφαρπάσαιο / τἀμὰ παίγνια ; Ar.~\textit{Eq.}~1200 οἴμοι τάλας ἀδίκως γε τἄμ᾽ ὑφήρπασας ; Philem.~fr.~3.13–15~K.-A. (=~4~Kock ; Ath.~13.569e) οὐκ ἔστ᾽ οὐδὲ εἷς / ἀκκισμός, οὐδὲ λῆρος, οὐδ᾽ ὑφήρπασεν, / ἀλλ᾽ εὐθὺς ἣν βούλει σὺ χὢν βούλει τρόπον ; Hsch.~υ~616 ὑποδραμών· ὑφαρπάσας (Ar. Equ. 742 codd.) ; Ar.~\textit{Eq.}~56 πανουργότατά πως περιδραμὼν ὑφαρπάσας ; X.~\textit{Cyr.}~8.4.16 ἐπεὶ καὶ τὴν ἕδραν σου ὑφήρπασε ; Arist.~\textit{HA}~9.1 ἡ μὲν γὰρ τῆς μεσημβρίας, διὰ τὸ μὴ ὀξὺ βλέπειν τὴν γλαῦκα τῆς ἡμέρας, κατεσθίει ὑφαρπάζουσα αὐτῆς τὰ ᾠά, ἡ δὲ γλαὺξ τῆς νυκτὸς τὰ τῆς κορώνης}
\lp{205 θήλειαν Κύπριν]}{Eur.~\textit{IA}~569–570 γυναιξὶ μὲν κατὰ Κύπριν / κρυπτάν ; \textit{AP}~12.86 (Mel.) ἁ Κύπρις θήλεια γυναικομανῆ] φλόγα βάλλει· / ἄρσενα δ᾽ αὐτὸς Ἔρως ἵμερον ἁνιοχεῖ ; Pl.~\textit{Smp.}~181c ὁ δὲ τῆς Οὐρανίας πρῶτον μὲν οὐ μετεχούσης θήλεος ἀλλ᾽ ἄρρενος μόνον—καὶ ἔστιν οὗτος ὁ τῶν παίδων ἔρως— ; [Luc.]~\textit{Am.}~5 ἀλλ᾽ ὁ μὲν ὑπερφυῶς παιδικοῖς ἥδετο τὴν θήλειαν Ἀφροδίτην βάραθρον ἡγούμενος, ὁ δ᾽ ἁγνεύων ἔρωτος ἄρρενος ἐς γυναῖκας ἐπτόητο ; Eur.~\textit{Andr.}~179 ἀλλ᾽ εἰς μίαν βλέποντες εὐναίαν Κύπριν ; Eur.~\textit{Ba.}~314–315 οὐχ ὁ Διόνυσος σωφρονεῖν ἀναγκάσει / γυναῖκας ἐς τὴν Κύπριν ; Eur.~\textit{Ba.}~773 οἴνου δὲ μηκέτ᾽ ὄντος οὐκ ἔστιν Κύπρις ; Eur.~\textit{Tro.}~368 οἳ διὰ μίαν γυναῖκα καὶ μίαν Κύπριν ; Aesch.~\textit{PV}~649–651 Ζεὺς γὰρ ἱμέρου βέλει / πρὸς σοῦ τέθαλπται καὶ συναίρεσθαι Κύπριν / θέλει ; Hsch.~θ~488 θῆλυ· ἀσθενές (Eur. Med. 928). γυνή ; Ar.~\textit{Ach.}~988 ὦ Κύπριδι τῇ καλῇ καὶ Χάρισι ταῖς φίλαις ξύντροφε Διαλλαγή ; Eur.~\textit{Andr.}~223 σοὶ καὶ ξυνήρων, εἴ τί σε σφάλλοι Κύπρις ; Eur.~fr.~951~N.² ἢν οἱ τεκόντες τοῦτο γιγνώσκωσ᾽ ὅτι / νέοι ποτ᾽ ἦσαν, ἠπίως τὴν τῶν τέκνων / οἴσουσι κύπριν ; Ibyc.~fr.~2.1–3~Bergk⁴ Ἔρος αὖτέ με κυανέοισιν ὑπὸ βλεφάροις τακέρ᾽ ὄμμασι δερκόμενος / κηλήμασι παντοδαποῖς ἐς ἄπειρα / δίκτυα Κύπριδός με βάλλει ; Manetho~1.30 θηλυτέρης Κύπριδός τε ; Alciphr.~1.6 καὶ τοῦ λάγνος εἶναι καὶ θηλυμανὴς ἀπόσχου ; Hsch.~θ~492 θηλυμανεῖς· πόρνοι (Ierem. 5, 8) ; Zen. Ath. IV (Miller, Mélanges, p.~381) Ἐν πλησμονῇ τοι Κύπρις, ἐν πεινῶσι δ᾽ οὔ. Ἐπὶ τῶν πλουσίων, καὶ διὰ τοῦτο τρυφώντων, καὶ ὑπὸ ἐρώτων δεδουλωμένων}
\lp{206 ἰδού γε κλέπτειν]}{Ar.~\textit{Eq.}~87 ἰδού γ᾽ ἄκρατον. περὶ πότου γοῦν ἐστί σοι ; Ar.~\textit{Nub.}~818 ἰδού γ᾽ ἰδού, Δί᾽ Ὀλύμπιον· τῆς μωρίας ; Ar.~\textit{Nub.}~1469 ἰδού γε Δία πατρῷον· ὡς ἀρχαῖος εἶ ; Ar.~\textit{Lys.}~441 ἰδού γ᾽ἐπιχεσεῖ. ποῦ ᾽στιν ἕτερος τοξότης ; Ar.~\textit{Ran.}~1205 ἰδού, σὺ δείξεις ; Ar.~\textit{Plu.}~1046 ποίου χρόνου ταλάνταθ᾽, ὃς παρ᾽ ἐμοὶ χθὲς ἦν ; Ar.~\textit{Eccl.}~136 ἰδού γε σοὶ πίνουσι ; Suid.~ι~128–129 Ἰδού: ἀντὶ τοῦ λάβε, κράτει. … Ἰδού, θέασαι}
\lp{206 νὴ Δία βινεῖσθαι μὲν οὖν]}{Ar.~\textit{Lys.}~876 ἐγὼ οὐ δεόμενος; ἐπιτετριμμένος μὲν οὖν ; Ar.~\textit{Plu.}~287 νὴ τοὺς θεοὺς Μίδαις μὲν οὖν, ἢν ὦτ᾽ ὄνου λάβητε ; Ar.~\textit{Nub.}~1112 ὠχρὸν μὲν οὖν οἶμαί γε καὶ κακοδαίμονα ; Ar.~\textit{Lys.}~715 βινητιῶμεν, ᾗ βράχιστον τοῦ λόγου ; Ar.~\textit{Thesm.}~257 κεκρυφάλου δεῖ καὶ μίτρας. — ἡδὶ μὲν οὖν ; Ar.~\textit{Thesm.}~861 πατὴρ ἐκεῖνός ἐστι; Φρυνώνδας μὲν οὖν ; Ar.~\textit{Plu.}~412 πολὺ μὲν οὖν νὴ τοὺς θεούς ; Ar.~\textit{Eq.}~338 οὐκ αὖ μ᾽ ἐάσεις; — μὰ Δία. — ναὶ μὰ Δία. — μὰ τὸν Ποσειδῶ ; Ar.~\textit{Ran.}~241 παύσασθε. — μᾶλλον μὲν οὖν ; Ar.~\textit{Eccl.}~1074 δευρὶ μὲν οὖν}
\lp{206 βινεῖσθαι]}{Ar.~\textit{Eq.}~877 ἔπαυσα τοὺς βινουμένους, τὸν Γρύττον ἐξαλείψας ; Ar.~\textit{Lys.}~1092 οὐκ ἔσθ᾽ ὅπως οὐ Κλεισθένη βινήσομεν ; Ar.~\textit{Thesm.}~35 καὶ μὴν βεβίνηκας σύ γ᾽, ἀλλ᾽ οὐκ οἶσθ᾽ ἴσως ; Ar.~\textit{Thesm.}~50 μῶν βινεῖσθαι ; Ar.~fr.~393~K.-A. (=~377~Kock ; Suid.~π~1531) κείσεσθον ὥσπερ πηνίω βινουμένω ; Ar.~\textit{Eccl.}~228 βινούμεναι χαίρουσιν ὥσπερ καὶ πρὸ τοῦ ; Ar.~\textit{Eccl.}~525 τί δ᾽; οὐχὶ βινεῖται γυνὴ κἄνευ μύρου ; Eup.~fr. *104.1–2~K.-A. (=~100~Kock ; Σ Aristid. III, p.~672 Dindorf) καὶ μηκέτ᾽, ὦναξ Μιλτιάδη καὶ Περίκλεες, / ἐάσατ᾽ ἄρχειν μειράκια κινούμενα ; Philetaer.~fr.~9.4~K.-A. (=~9~Kock ; Ath.~13.587e) οὐχὶ Λαὶς μὲν τελευτῶσ᾽ ἀπέθανεν βινουμένη ; Hsch.~β~466 βινεῖν· παρὰ Σόλωνι (legg. F 27) τὸ βίᾳ μίγνυσθαι. τὸ δὲ κατὰ νόμον ὀπύειν ; Suid.~β~287 Βινεῖν: τὸ συνουσιάζειν· ἢ τὸ πιπίζειν ; Suid.~β~288 Βινεῖν: περαίνειν}
\lp{206 ἰδού γε κλέπτειν· … βινεῖσθαι μὲν οὖν (τὴν σκάφην σκάφην λέγω)]}{Ar.~fr.~927~K.-A. (901b, adesp.~227~Kock) ἄγροικός εἰμι· τὴν σκάφην σκάφην λέγω ; Tosi n°~2272 Τὰ σῦκα σῦκα, τὴν σκάφην δὲ σκάφην ὀνομάζων}
\lp{207 ἀτὰρ … γε]}{Ar.~\textit{Ach.}~448 ἀτὰρ δέομαί γε πτωχικοῦ βακτηρίου ; Ar.~\textit{Av.}~144 ἀτὰρ ἔστι γ᾽ ὁποίαν λέγετον εὐδαίμων πόλις ; Ar.~\textit{Thesm.}~87 ἀτὰρ τίν᾽ ἐκ τούτων σὺ μηχανὴν ἔχεις ; Ar.~\textit{Thesm.}~215 ἀτὰρ τί μέλλεις δρᾶν μ᾽ ; Ar.~\textit{Eccl.}~1067 ἀτὰρ ἥτις εἶ γε, πόλλ᾽ ἀγαθὰ γένοιτό σοι ; Suid.~α~4319 Ἀτάρ: καὶ δή· ἢ πλὴν, ἅτε δή· ἢ ὅμως· ἢ ἀντὶ τοῦ δέ}
\lp{207 ἡ πρόφασις]}{Ar.~\textit{Ach.}~392–394 ὡς σκῆψιν ἁγὼν οὗτος οὐκ ἐσδέξεται. / ὥρα ᾽στὶν ἤδη καρτερὰν ψυχὴν λαβεῖν, / καί μοι βαδιστέ᾽ ἐστὶν ὡς Εὐριπίδην ; Men.~fr.~171~K.-A. (Stob.~3.4.39) Μικρά γε πρόφασίς ἐστι τοῦ πρᾶξαι κακῶς ; Aristopho fr.~9.6~K.-A. (=~9~Kock ; Ath.~4.161c–d) τῆς εὐτελείας πρόφασιν εὑρόντες καλὴν ; D.~21.98 καὶ τίν᾽, ὢ πρὸς τῶν θεῶν, ἕξετ᾽ εἰπεῖν πρόφασιν δικαίαν ἢ καλήν; ὅτι νὴ Δί᾽ ἀσελγής ἐστι καὶ βδελυρός ; Antiph.~5.26 οὔτε τῷ ἀπάγοντι νύκτωρ μακρὰν ὁδὸν ἡ πρόφασις ἂν εἰκότως ἐγίγνετο ; Zen.~2.45 Ἀγὼν πρόφασιν οὐκ ἐπιδέχεται, οὔτε φιλία ; App. Prov.~3.5 Ἡ Θρᾷττα ἐπὶ τῷ ταρίχει: ἐπὶ τῶν κρυπτόντων τὰ ἴδια πάθη, ἢ ἐπὶ τῶν προφασιζομένων ἵνα τύχοιεν τινός ; Alex.~fr.~132~K.-A. (=~127~Kock ; Ath.~4.170a) καὶ μὴ προφάσεις ἐνταῦθά μοι, μηδ᾽ «οὐκ ἔχω.» ; Philem.~fr.~96.9–10~K.-A. (=~93~Kock ; Stob. Flor.~98.14~Meineke) οὐκ ἔστ᾽ ἀποτυχεῖν / κακοῦ, πρόφασιν δ᾽ ἀεί τιν᾽ ἐξευρίσκομεν ; Eur.~\textit{Hec.}~340 ἔχεις δὲ πρόφασιν ; Eur.~\textit{IA}~1180 ἐπεὶ βραχείας προφάσεως ἔδει μόνον ; Σ Eur.~\textit{Hec.}~340 πρόφασις κυρίως ὁ προτεινόμενος εὐπρόσωπος λόγος ἐπὶ κατηγορίᾳ τινὸς ; Hsch.~π~3940 πρόσχημα· πρόφασις. ὑπόκρισις. προκάλυμμα ; Hsch.~σ~965 σκῆψις· πρόφασις ; Suid.~σ~589 Σκῆψις: παρακάλυμμα. πρόφασις, ἢ φροντίς, ἀναβολή ; Diogenian.~7.47 Πατρόκλειος πρόφασις: ἐπὶ τῶν μὴ δυναμένων διὰ φόβον θρηνεῖν τὰς οἰκείας συμφορὰς, ἐξ ἑτέρων δὲ θλίψεων αὐτὰς ἀποκλαιόντων ; App. Prov.~2.41 Ἐκ παντὸς πρόφασις: ἐπὶ τῶν προφάσεως οὐκ ἀπορούντων ; Th.~1.23.6 τὴν μὲν γὰρ ἀληθεστάτην πρόφασιν, ἀφανεστάτην δὲ λόγῳ ; Hsch.~α~2163 αἰτία· πρόφασις ; Hsch.~π~4063 πρόφασις· αἰτία, ἀφορμή ; Hsch.~π~4078 προχάνη· σκῆψις. πρόφασις … καὶ καλύπτρα ; Suid.~μ~1417 Μύνη: ἡ πρόφασις}
\lp{207 ἡ πρόφασις (προφάσεως δεῖται μόνον ἡ πονηρία)]}{Arist.~\textit{Rh.}~1373a ὥσπερ γὰρ ἡ παροιμία, προφάσεως δεῖται μόνον ἡ πονηρία ; \textit{idem} Diogenian.~7.87, Greg. Cypr.~3.47, Tosi n°~2213 ; Erasm.~\textit{Adag.}~2.1.68 (LB II, col.~432 A–D) Occasione duntaxat opus improbitati. Aristoteles lib. Rhetor. II : … Ὥσπερ γὰρ ἡ παροιμία, Προφάσεως δεῖται μόνον ἡ πονηρία. … Quadrabit in nocendi cupidos, qui, ne temere laesisse videantur, causas praetexunt}
\lp{207 εἰκότως ἔχει]}{Ar.~\textit{Plu.}~1021 εἰ Θάσιον ἐνέχεις, εἰκότως γε νὴ Δία ; Eur.~\textit{IT}~911 σθένειν τὸ θεῖον μᾶλλον εἰκότως ἔχει ; D.~19.215 ἀλλ᾽ ὅμως ἐρεῖ, καὶ νὴ Δί᾽ εἰκότως γε ; D.~23.107 ὅτι νὴ Δί᾽ εἰκότως ἐφενακίσθητε καὶ παρεκρούσθητε ; Hsch.~ε~842 εἰκότως· πρεπόντως. εὐλόγως. δικαίως ; Ar.~\textit{Av.}~273 εἰκότως ⟨γε⟩· καὶ γὰρ ὄνομ᾽ αὐτῷ ᾽στὶ φοινικόπτερος ; Aeschin.~2.33 εἰκότως ἔχειν}
\lp{208 τί οὖν; ποιήσεις ταῦτα;]}{Ar.~\textit{Lys.}~861 τί οὖν; δώσεις τί μοι ; Luc.~\textit{DMeretr.}~6.4 τί φής; ποιήσεις ταῦτα; ποιήσεις, οἶδα ἐγώ ; Pl.~\textit{Prm.}~135c τί οὖν ποιήσεις φιλοσοφίας πέρι ; Ar.~\textit{Lys.}~506 σὺ δέ μοι λέγε. — ταῦτα ποιήσω ; Ar.~\textit{Ran.}~1515 ταῦτα ποιήσω· σὺ δὲ τὸν θᾶκον ; Ar.~\textit{Plu.}~94 ὁμολογῶ σοι. — φέρε τί οὖν}
\lp{208 μὴ δόκει γε σύ]}{Eur.~\textit{Med.}~365 ἀλλ᾽ οὔτι ταύτῃ ταῦτα, μὴ δοκεῖτέ πω ; Ar.~\textit{Vesp.}~15 ἀτὰρ σὺ λέξον πρότερος. — ἐδόκουν αἰετὸν ; Eur.~\textit{Hipp.}~119 μάταια βάζει· μὴ δόκει τούτων κλύειν ; [Eur.]~\textit{Rh.}~940 ἔδρασε δράσας — μὴ δόκει λεληθέναι ; Soph.~\textit{El.}~312 ἦ κάρτα· μὴ δόκει μ᾽ ἄν, εἴπερ ἦν πέλας ; Σ Eur.~\textit{Hipp.}~119 μὴ δόκει τούτου κλύειν: νόμιζε μηδὲ ἀκούειν τοῦ τοιούτου λόγου ; Hsch.~δ~2112 δοκεῖ· οἴεται, νομίζει ; Ar.~\textit{Vesp.}~922 μή νυν ἀφῆτέ γ᾽ αὐτόν ; Eur.~fr.~641 τὰ δ᾽ ἄλλα μὴ δόκει ξυνιέναι}
